% Leçon 23 — Expression écrite : la démocratie — Chinois Tle A — Cours signé OnBuch+
% Programme officiel MINESEC. Compiler avec : tectonic ZH23-ecrit-la-democratie.tex
\newcommand{\DOCMATIERE}{Chinois}
\newcommand{\DOCNIVEAU}{Tle A}
\newcommand{\DOCMODULE}{Module 3 : Citoyenneté et ouverture au monde}
\newcommand{\DOCLECON}{ZH23}
\newcommand{\DOCTITRE}{Leçon 23 — Expression écrite : la démocratie}
% OnBuch+ — préambule « cours signés » ENRICHI (compilé avec tectonic / XeLaTeX)
% Système visuel « Pop » d'OnBuch+ : fond crème (jamais blanc), bordures encre
% épaisses, ombres « dures » décalées, titres Archivo Black en capitales.
% Version MATHÉMATIQUES : graphes (pgfplots/tikz), boîtes pédagogiques
% (définition, propriété, méthode, attention, le savais-tu, exercice résolu,
% exercice, TP, bilan), page de garde et signature OnBuch+ ancrée en bas.
%
% Macros attendues dans main.tex AVANT \input{preamble} :
%   \DOCMATIERE  \DOCNIVEAU  \DOCMODULE  \DOCLECON (numéro)  \DOCTITRE
\documentclass[11pt]{article}

\usepackage{fontspec}
\usepackage[french]{babel} % français = langue principale ; le chinois via \zh{..} / textebox (xeCJK)
\usepackage{xeCJK}
\usepackage{pifont} % \ding{51} \ding{55} utilisés par les modèles
\usepackage[a4paper,margin=2cm,top=3.4cm,bottom=2.8cm,headheight=1.4cm,footskip=1.3cm]{geometry}
\usepackage{amsmath,amssymb,amsfonts}
\usepackage{mathtools} % \xmapsto, \coloneqq... souvent utilisés par les modèles
\usepackage{cancel} % simplifications barrées (bilans de forces)
% \abs{z} (module, valeur absolue) et \norm{u} (norme), utilisés par les modèles
\providecommand{\abs}{}\let\abs\relax\DeclarePairedDelimiter{\abs}{\lvert}{\rvert}
\providecommand{\norm}{}\let\norm\relax\DeclarePairedDelimiter{\norm}{\lVert}{\rVert}
\usepackage[table]{xcolor} % [table] : fournit aussi \rowcolors (alternance de lignes)
\usepackage{pagecolor}
\usepackage{tikz}
\usetikzlibrary{arrows.meta,positioning,calc,shapes.geometric,shapes.misc,shapes.arrows,decorations.pathmorphing,decorations.pathreplacing,decorations.markings,patterns,shadows,mindmap,trees,fit,backgrounds,matrix,intersections,angles,quotes}
\usepackage{pgfplots}
\usepackage[version=4]{mhchem} % équations nucléaires \ce{^{235}_{92}U}
\usepackage[european,siunitx]{circuitikz} % schémas électriques
\pgfplotsset{compat=1.18}
\usepgfplotslibrary{fillbetween,groupplots} % aires entre courbes (calcul intégral)
\usepackage{enumitem}
\usepackage{fancyhdr}
% [most] charge toutes les bibliothèques tcolorbox (listings, minted, raster,
% documentation...) et gonfle inutilement la mémoire TeX sur un document long
% (~300 boîtes) : on ne charge que ce qui est réellement utilisé.
\usepackage[skins,breakable]{tcolorbox}
\usepackage{graphicx}
\usepackage{colortbl}
\usepackage{booktabs}
\usepackage{tabularx}
\usepackage{diagbox} % case d'en-tête diagonale des tableaux à double entrée
% Colonnes extensibles centrée / gauche / droite (C, L, R), souvent utilisées par les modèles
\newcolumntype{C}{>{\centering\arraybackslash}X}
\newcolumntype{L}{>{\raggedright\arraybackslash}X}
\newcolumntype{R}{>{\raggedleft\arraybackslash}X}
\usepackage{array}
\usepackage{multicol}
\usepackage{multirow}
\usepackage{siunitx}
% parse-numbers=false : accepte aussi \SI{2,5 \times 10^{-3}}{...}
\sisetup{locale=FR,output-decimal-marker={,},group-minimum-digits=4,parse-numbers=false}
\DeclareSIUnit\ohm{\ensuremath{\symup{\Omega}}} % U+2126 absent de la police
\DeclareSIUnit\division{div} % réglages d'oscilloscope (ms/div, V/div)
\DeclareSIUnit\tr{tr} % tours (vitesse de rotation, ex. \SI{1500}{\tr\per\minute})
\DeclareSIUnit\molar{M} % concentration molaire (ex. \SI{3}{\molar}, \SI{1}{\micro\molar})
\DeclareSIUnit\an{an} % année (le modèle confond parfois avec \year, un compteur TeX) — ex. \SI{5}{\kilo\meter\per\an}
\DeclareSIUnit\FCFA{FCFA} % franc CFA (ex. \SI{500}{\FCFA\per\kilo\gram})
\DeclareSIUnit\degreeBrix{\textdegree Brix} % teneur en sucre d'un sirop/jus (réfractométrie)
\DeclareSIUnit\month{mois} % le modèle utilise parfois \month, un compteur TeX, comme unité de durée
\usepackage{qrcode}
% \degree : commande fréquemment utilisée par les modèles pour un angle ou une
% température en dehors de \SI{}{} (non fournie par les paquets chargés).
\providecommand{\degree}{^{\circ}}
% \qed (fin de démonstration), utilisé par les modèles sans amsthm
\providecommand{\qed}{\hfill$\square$}
% \barre : double barre des valeurs interdites dans un tableau de variations (tkz-tab)
\providecommand{\barre}{\ensuremath{\|}}
% environnement proof (amsthm non chargé) utilisé par les modèles
\ifdefined\proof\else\newenvironment{proof}[1][Démonstration]{\par\noindent\textit{#1.}\ }{\qed\par}\fi
\usepackage{lastpage}
\usepackage{hyperref}
\hypersetup{hidelinks}

% ============================================================
% Palette « Pop » OnBuch+ — mêmes hex que la classe P (plus_ui.dart)
% ============================================================
\definecolor{poporange}{HTML}{F59321}
\definecolor{popdark}{HTML}{E07A0C}
\definecolor{popcream}{HTML}{FFF6E8}
\definecolor{popink}{HTML}{1C1714}
\definecolor{poppurple}{HTML}{7A5AE0}
\definecolor{popgreen}{HTML}{2E9E6A}
\definecolor{poppink}{HTML}{E0457B}
\definecolor{popblue}{HTML}{2F7DE1}
\definecolor{popgold}{HTML}{C98A12}
\definecolor{popmuted}{HTML}{8A7F76}
\definecolor{popline}{HTML}{EADFD2}
\definecolor{popgray}{HTML}{8A7F76}
\colorlet{popgrey}{popgray}
% Alias tolérants (le modèle nomme parfois les couleurs différemment)
\colorlet{popwhite}{white}
\colorlet{popblack}{popink}
\colorlet{popred}{poppink}
\colorlet{popyellow}{popgold}
% Teintes claires (fonds de tableaux/encadrés) — le modèle nomme parfois
% les couleurs avec un suffixe L (ex. popblueL) sans les avoir définies.
\colorlet{poporangeL}{poporange!20}
\colorlet{popdarkL}{popdark!20}
\colorlet{poppurpleL}{poppurple!20}
\colorlet{popgreenL}{popgreen!20}
\colorlet{poppinkL}{poppink!20}
\colorlet{popblueL}{popblue!20}
\colorlet{popgoldL}{popgold!20}
\colorlet{popredL}{popred!20}
\colorlet{popgrayL}{popgray!20}
% Teintes claires pour fonds de tableaux / zones
\colorlet{popblueL}{popblue!10!white}
\colorlet{poporangeL}{poporange!14!white}
\colorlet{popgreenL}{popgreen!12!white}
\colorlet{poppurpleL}{poppurple!12!white}
\colorlet{poppinkL}{poppink!12!white}
\colorlet{popgoldL}{popgold!14!white}

\pagecolor{popcream}

% ============================================================
% Typographie
% ============================================================
\defaultfontfeatures{Ligatures=TeX}
\setmainfont{PlusJakartaSans-Medium.ttf}[Path=../../../fonts/,BoldFont=PlusJakartaSans-Bold.ttf]
% Symboles, API et emojis absents de la police principale : repli sur DejaVu Sans (caractères actifs)
\newfontface\symfont{DejaVuSans.ttf}[Path=../../../fonts/]
\makeatletter
\newcommand\symactive[1]{\catcode#1=13 \begingroup\lccode`\~=#1 \lowercase{\endgroup\def~}{{\symfont\char#1}}}
\newcommand\symblank[1]{\catcode#1=13 \begingroup\lccode`\~=#1 \lowercase{\endgroup\def~}{}}
\newcommand\symhyph[1]{\catcode#1=13 \begingroup\lccode`\~=#1 \lowercase{\endgroup\def~}{\mbox{-}}}
\newcount\symc
\newcommand\symrange[2]{\symc=#1 \@whilenum\symc<#2 \do{\expandafter\symactive\expandafter{\the\symc}\advance\symc by 1 }}
\newcommand\symblankrange[2]{\symc=#1 \@whilenum\symc<#2 \do{\expandafter\symblank\expandafter{\the\symc}\advance\symc by 1 }}
\makeatother
\symrange{"0250}{"02FF}\symactive{"2713}\symactive{"2714}\symactive{"2717}\symactive{"2718}\symactive{"2610}\symactive{"2611}\symactive{"2612}
\symrange{"2600}{"27BF}
\catcode"2705=13 \begingroup\lccode`\~="2705 \lowercase{\endgroup\def~}{{\symfont\char"2713}}\symblank{"26BD}\symblank{"26A1}
\symrange{"2460}{"257F}\symactive{"223C}\symactive{"2248}\symactive{"2260}
\symactive{"01F8}\symactive{"01F9}\symactive{"1E3E}\symactive{"1E3F}
\symhyph{"2011}\symhyph{"2E1A}\symblankrange{"1F300}{"1FAFF}\symblank{"FE0F}
% Caractères chinois : WenQuanYi Zen Hei (sous-ensemble GB2312 + pinyin), gras/italique simulés
\setCJKmainfont{WenQuanYiZenHei-subset.ttf}[Path=../../../fonts/,AutoFakeBold=3,AutoFakeSlant=0.2]
\xeCJKsetup{CJKspace=false}
% Pinyin : voyelles à caron (ǎ ǐ ǒ ǔ ǖ ǘ ǚ ǜ) absentes de la police principale -> caractères actifs
\newfontface\pyfont{WenQuanYiZenHei-subset.ttf}[Path=../../../fonts/]
\newcommand\pyactive[1]{\catcode#1=13 \begingroup\lccode`\~=#1 \lowercase{\endgroup\def~}{{\pyfont\char#1}}}
\pyactive{"01CD}\pyactive{"01CE}\pyactive{"01CF}\pyactive{"01D0}\pyactive{"01D1}\pyactive{"01D2}\pyactive{"01D3}\pyactive{"01D4}\pyactive{"01D5}\pyactive{"01D6}\pyactive{"01D7}\pyactive{"01D8}\pyactive{"01D9}\pyactive{"01DA}\pyactive{"01DB}\pyactive{"01DC}
% Maths en sans-serif (Fira Math) avec chiffres et lettres droites de Plus
% Jakarta, pour que les formules se fondent dans le texte.
\usepackage[math-style=ISO,bold-style=ISO]{unicode-math}
\setmathfont{FiraMath-Regular.otf}[Path=../../../fonts/]
\setmathfont{PlusJakartaSans-Medium.ttf}[Path=../../../fonts/,range={up/num,up/{Latin,latin},\mathup}]
\usepackage{icomma} % virgule décimale sans espace en mode math
% Caractères mathématiques Unicode absents des polices de texte (Plus Jakarta,
% Archivo Black) : rendus par leur équivalent mathématique, en texte comme en
% formule — sinon ils s'affichent en carré vide (ex. « Arithmétique dans ℤ »).
\usepackage{newunicodechar}
\newunicodechar{ℝ}{\ensuremath{\mathbb{R}}}
\newunicodechar{ℤ}{\ensuremath{\mathbb{Z}}}
\newunicodechar{ℂ}{\ensuremath{\mathbb{C}}}
\newunicodechar{ℕ}{\ensuremath{\mathbb{N}}}
\newunicodechar{ℚ}{\ensuremath{\mathbb{Q}}}
\newunicodechar{𝔻}{\ensuremath{\mathbb{D}}}
\newunicodechar{α}{\ensuremath{\alpha}}
\newunicodechar{β}{\ensuremath{\beta}}
\newunicodechar{γ}{\ensuremath{\gamma}}
\newunicodechar{δ}{\ensuremath{\delta}}
\newunicodechar{ε}{\ensuremath{\varepsilon}}
\newunicodechar{θ}{\ensuremath{\theta}}
\newunicodechar{λ}{\ensuremath{\lambda}}
\newunicodechar{μ}{\ensuremath{\mu}}
\newunicodechar{σ}{\ensuremath{\sigma}}
\newunicodechar{τ}{\ensuremath{\tau}}
\newunicodechar{φ}{\ensuremath{\varphi}}
\newunicodechar{ψ}{\ensuremath{\psi}}
\newunicodechar{ω}{\ensuremath{\omega}}
\newunicodechar{Σ}{\ensuremath{\Sigma}}
% majuscules produites par \MakeUppercase dans les titres (α-aminés -> Α)
\newunicodechar{Α}{\ensuremath{\alpha}}
\newunicodechar{Β}{\ensuremath{\beta}}
\newunicodechar{Γ}{\ensuremath{\Gamma}}
\newunicodechar{Δ}{\ensuremath{\Delta}}
\newunicodechar{Ε}{\ensuremath{\varepsilon}}
\newunicodechar{Θ}{\ensuremath{\Theta}}
\newunicodechar{Λ}{\ensuremath{\Lambda}}
\newunicodechar{Μ}{\ensuremath{\mu}}
\newunicodechar{Τ}{\ensuremath{\tau}}
\newunicodechar{Φ}{\ensuremath{\Phi}}
\newunicodechar{Ψ}{\ensuremath{\Psi}}
\newunicodechar{Ω}{\ensuremath{\Omega}}
\newunicodechar{↦}{\ensuremath{\mapsto}}
\newunicodechar{∈}{\ensuremath{\in}}
\newunicodechar{∉}{\ensuremath{\notin}}
\newunicodechar{∧}{\ensuremath{\wedge}}
\newunicodechar{⇔}{\ensuremath{\Leftrightarrow}}
\newunicodechar{⟺}{\ensuremath{\Longleftrightarrow}}
\newunicodechar{⇒}{\ensuremath{\Rightarrow}}
\newunicodechar{⟹}{\ensuremath{\Longrightarrow}}
\newunicodechar{∘}{\ensuremath{\circ}}
\newunicodechar{∪}{\ensuremath{\cup}}
\newunicodechar{∩}{\ensuremath{\cap}}
\newunicodechar{≡}{\ensuremath{\equiv}}
\newunicodechar{⊂}{\ensuremath{\subset}}
\newunicodechar{⊥}{\ensuremath{\perp}}
\newunicodechar{∃}{\ensuremath{\exists}}
\newunicodechar{∀}{\ensuremath{\forall}}
\newunicodechar{∛}{\ensuremath{\sqrt[3]{\,}}}
\newunicodechar{∜}{\ensuremath{\sqrt[4]{\,}}}
\newunicodechar{⃗}{\ensuremath{{}^{\to}}}
\newunicodechar{ⁿ}{\ifmmode^{n}\else\textsuperscript{\ensuremath{n}}\fi}
\newunicodechar{ˣ}{\ifmmode^{x}\else\textsuperscript{\ensuremath{x}}\fi}
\newunicodechar{ᵇ}{\ifmmode^{b}\else\textsuperscript{\ensuremath{b}}\fi}
\newunicodechar{ʳ}{\ifmmode^{r}\else\textsuperscript{\ensuremath{r}}\fi}
\newunicodechar{ᵘ}{\ifmmode^{u}\else\textsuperscript{\ensuremath{u}}\fi}
\newunicodechar{ʸ}{\ifmmode^{y}\else\textsuperscript{\ensuremath{y}}\fi}
\newunicodechar{ᵃ}{\ifmmode^{a}\else\textsuperscript{\ensuremath{a}}\fi}
\newunicodechar{ᵉ}{\ifmmode^{e}\else\textsuperscript{\ensuremath{e}}\fi}
\newunicodechar{ᵏ}{\ifmmode^{k}\else\textsuperscript{\ensuremath{k}}\fi}
\newunicodechar{ᵖ}{\ifmmode^{p}\else\textsuperscript{\ensuremath{p}}\fi}
\newunicodechar{ᴺ}{\ifmmode^{N}\else\textsuperscript{\ensuremath{N}}\fi}
\newunicodechar{ᶜ}{\ifmmode^{c}\else\textsuperscript{\ensuremath{c}}\fi}
\newunicodechar{ʰ}{\ifmmode^{h}\else\textsuperscript{\ensuremath{h}}\fi}
\newunicodechar{ᵠ}{\ifmmode^{q}\else\textsuperscript{\ensuremath{q}}\fi}
\newunicodechar{ₙ}{\ifmmode_{n}\else\textsubscript{\ensuremath{n}}\fi}
\newunicodechar{ₐ}{\ifmmode_{a}\else\textsubscript{\ensuremath{a}}\fi}
\newunicodechar{ᵢ}{\ifmmode_{i}\else\textsubscript{\ensuremath{i}}\fi}
\newunicodechar{ᵣ}{\ifmmode_{r}\else\textsubscript{\ensuremath{r}}\fi}
\newunicodechar{ₖ}{\ifmmode_{k}\else\textsubscript{\ensuremath{k}}\fi}
\newunicodechar{ᵦ}{\ifmmode_{\beta}\else\textsubscript{\ensuremath{\beta}}\fi}
\newunicodechar{ₓ}{\ifmmode_{x}\else\textsubscript{\ensuremath{x}}\fi}
\newfontfamily\popdisplay{ArchivoBlack-Regular.ttf}[Path=../../../fonts/]
\newfontfamily\poplogo{SpaceGrotesk-Bold.ttf}[Path=../../../fonts/]
\newfontfamily\popsemi{PlusJakartaSans-SemiBold.ttf}[Path=../../../fonts/]
\newfontfamily\popxbold{PlusJakartaSans-ExtraBold.ttf}[Path=../../../fonts/]
\newfontfamily\popxboldls{PlusJakartaSans-ExtraBold.ttf}[Path=../../../fonts/,LetterSpace=3.0]

\setlist[enumerate]{leftmargin=1.7em,itemsep=5pt,topsep=4pt}
\setlist[itemize]{leftmargin=1.7em,itemsep=4pt,topsep=4pt,label={\color{poporange}\textbullet}}
\setlength{\parindent}{0pt}
\setlength{\parskip}{5pt}
\renewcommand{\arraystretch}{1.25}

% pgfplots : style Pop par défaut
\pgfplotsset{
  every axis/.append style={axis line style={popink,line width=1pt},tick style={popink},
    grid style={popline,line width=0.5pt},label style={font=\small},tick label style={font=\footnotesize}},
  popcurve/.style={poporange,line width=1.8pt,no markers,smooth},
}
% Même style disponible dans un \draw TikZ simple (hors pgfplots)
\tikzset{popcurve/.style={poporange,line width=1.8pt}}
\tikzset{popgrid/.style={popline,very thin}}
\colorlet{popcurve}{poporange} % le nom du style est parfois utilisé comme couleur

% ============================================================
% Pill / squiggle
% ============================================================
\newcommand{\poppill}[3]{%
  \tikz[baseline=-0.55ex]\node[draw=popink,line width=1.2pt,rounded corners=2.6pt,%
    fill=#2,inner xsep=6.5pt,inner ysep=3pt]{%
    \popxboldls\fontsize{7}{7}\selectfont\color{#3}\MakeUppercase{#1}};%
}
\newcommand{\popsquiggle}[1]{%
  \tikz[baseline=-0.4ex]{
    \draw[#1,line width=2.6pt,line cap=round]
      plot [smooth] coordinates {(0,0.09) (0.34,0.01) (0.68,0.21) (1.02,0.01) (1.36,0.10)};
  }%
}
% Mot-clé surligné (marqueur orange sous le texte)
\newcommand{\cle}[1]{{\popxbold\color{popdark}#1}}

% ============================================================
% En-tête / pied
% ============================================================
\pagestyle{fancy}
\fancyhf{}
\renewcommand{\headrulewidth}{0pt}
\renewcommand{\footrulewidth}{0pt}
\lhead{\raisebox{-0.15ex}{\poplogo\fontsize{13}{13}\selectfont\color{popink}OnBuch\color{poporange}+}}
\rhead{\poppill{\DOCMATIERE\ \textbullet\ \DOCNIVEAU\ \textbullet\ Leçon \DOCLECON}{white}{popink}}
\lfoot{\popxbold\fontsize{7.6}{7.6}\selectfont\color{popmuted}\MakeUppercase{Cours signé OnBuch+}\ \textbullet\ {\popsemi\DOCTITRE}}
\rfoot{\tikz[baseline=-0.6ex]\node[rounded corners=8.5pt,fill=popink,minimum height=17pt,minimum width=17pt,inner xsep=5pt,inner sep=0]{\popxbold\fontsize{8}{8}\selectfont\color{popcream}\thepage\,/\,\pageref*{LastPage}};}

% ============================================================
% Titres
% ============================================================
\newcommand{\chaptitre}[2]{%
  \begin{tcolorbox}[enhanced,colback=poporange,colframe=popink,boxrule=2.6pt,arc=11pt,
      drop shadow=popink,left=15pt,right=15pt,top=12pt,bottom=11pt,before skip=3pt,after skip=18pt]
    \poppill{#2}{popink}{popcream}\par\vspace{9pt}
    {\raggedright\hyphenpenalty=10000\exhyphenpenalty=10000\popdisplay\fontsize{22}{24}\selectfont\color{popcream}\MakeUppercase{#1}\par}
  \end{tcolorbox}
}
% Section numérotée : pastille orange avec le numéro + titre Archivo
\newcounter{popsec}
\newcommand{\coursec}[1]{%
  \stepcounter{popsec}\setcounter{popsub}{0}%
  \par\vspace{16pt}\noindent
  \begin{minipage}{\linewidth}
  \tikz[baseline=-0.7ex]\node[draw=popink,line width=1.6pt,fill=poporange,rounded corners=4pt,minimum size=22pt,inner sep=0]{\popdisplay\fontsize{12}{12}\selectfont\color{popcream}\thepopsec};%
  \hspace{8pt}{\popdisplay\fontsize{15}{17}\selectfont\color{popink}\MakeUppercase{#1}}\par
  \vspace{4pt}\noindent{\color{poporange}\rule{36pt}{3.6pt}}
  \end{minipage}\par\nopagebreak\vspace{8pt}
}
\newcounter{popsub}
\newcommand{\courssub}[1]{%
  \stepcounter{popsub}%
  \par\vspace{9pt}\noindent{\popxbold\fontsize{11.5}{13}\selectfont\color{popink}{\color{poporange}\thepopsec.\thepopsub}\hspace{6pt}#1}\par\nopagebreak\vspace{4pt}
}

% ============================================================
% Boîtes pédagogiques — même recette : fond blanc, bordure épaisse colorée,
% ombre dure encre, badge plein. Titre optionnel [#1].
% ============================================================
\tcbset{popbase/.style={breakable,enhanced,colback=white,boxrule=2.3pt,arc=9pt,
  shadow={3pt}{-3pt}{0pt}{popink},left=14pt,right=14pt,top=11pt,bottom=11pt,before skip=11pt,after skip=15pt,
  fontupper=\popsemi\fontsize{10.6}{14.5}\selectfont\color{popink},
  fontlower=\popsemi\fontsize{10.6}{14.5}\selectfont\color{popink}}}
% Coche dessinée (le glyphe ✓ n'existe pas dans les polices du document)
\newcommand{\popcheck}[1]{\tikz[baseline=-0.5ex]\draw[#1,line width=1.6pt,line cap=round,line join=round](-0.13,0.0)--(-0.04,-0.09)--(0.13,0.1);}
\providecommand{\coche}{\popcheck{popgreen}}
\providecommand{\resultat}[1]{\par\smallskip\noindent\fcolorbox{popgreen}{popgreenL}{\parbox{\dimexpr\linewidth-2\fboxsep-2\fboxrule\relax}{\textbf{Résultat.} #1}}\par\smallskip}
\newcommand{\popboxhead}[3]{\poppill{#1}{#2}{white}\ifx\relax#3\relax\else\hspace{6pt}{\popxbold\fontsize{10.6}{12}\selectfont\color{popink}#3}\fi\par\vspace{7pt}}

\newtcolorbox{aretenir}[1][]{popbase,colframe=popgreen,before upper={\popboxhead{À retenir}{popgreen}{#1}}}
% Texte en chinois (document de lecture, dialogue, extrait) : cadre
% d'étude. Un titre optionnel (source, auteur) s'affiche dans l'en-tête.
\newtcolorbox{textebox}[1][]{popbase,colframe=popblue,before upper={\popboxhead{文本 · Texte}{popblue}{#1}},after upper={}}
% Marques ✓ / ✗ (forme correcte / fautive) souvent utilisées par les modèles
\providecommand{\cross}{\textcolor{poppink}{\ensuremath{\times}}}
\providecommand{\xmark}{\cross}\providecommand{\crossmark}{\cross}\providecommand{\cmark}{\ensuremath{\checkmark}}
% Ligne de réponse / justification à compléter (inventée par les modèles)
\providecommand{\justification}[1]{\par\noindent\textit{Justification~:} \dotfill\par}
% \textipa (package tipa non chargé) : la police ne couvre pas l'API -> simple passage
\providecommand{\textipa}[1]{#1}
% Soulignement (ulem non chargé) : \uline, \ul
\providecommand{\uline}[1]{\underline{#1}}\providecommand{\ul}[1]{\underline{#1}}
% \glossaire{\item ...} inventée par les modèles : liste de glossaire
\providecommand{\glossaire}[1]{\begin{itemize}#1\end{itemize}}
% \alert (beamer) inventée par les modèles : texte mis en évidence
\providecommand{\alert}[1]{\textbf{\textcolor{poppink}{#1}}}
% variantes ASCII d'ordinaux écrites en macro
\providecommand{\iere}{\textsuperscript{re}}\providecommand{\ieme}{\textsuperscript{e}}\providecommand{\ere}{\textsuperscript{re}}\providecommand{\eme}{\textsuperscript{e}}\providecommand{\er}{\textsuperscript{er}}\providecommand{\ie}{\textsuperscript{e}}
% Ordinaux français écrits en macro par les modèles : 1\ière, 3\ième
\newcommand{\ière}{\textsuperscript{re}}\newcommand{\ième}{\textsuperscript{e}}
% \exemple (inventée par les modèles) : étiquette « Exemple : »
\providecommand{\exemple}{\textbf{Exemple~:}\ }
% \square utilisable aussi hors mode math (cases à cocher)
\AtBeginDocument{\let\squareMath\square\renewcommand{\square}{\ensuremath{\squareMath}}}
% Mot ou phrase en chinois à l'intérieur d'un paragraphe français
\newcommand{\zh}[1]{\ifmmode\text{#1}\else{#1}\fi}
\let\eng\zh \let\esp\zh \let\all\zh % alias tolérés
% flèches d'intonation inventées par les modèles
\providecommand{\textsearrow}{}\renewcommand{\textsearrow}{\ensuremath{\searrow}}\providecommand{\textnearrow}{}\renewcommand{\textnearrow}{\ensuremath{\nearrow}}
% \fr{..} : mot français (inventé par les modèles)
\providecommand{\fr}[1]{\textit{#1}}
% zhbox : encadré de texte chinois inventé par les modèles
\newenvironment{zhbox}{\begin{quote}}{\end{quote}}
% \courssubsub : titre de niveau 3 inventé par les modèles
\providecommand{\courssubsub}[1]{\par\medskip\noindent\textbf{#1}\par\smallskip}
% \vs : « versus » inventé par les modèles
\providecommand{\vs}{\textit{vs}\ }
% \blank : case à compléter inventée par les modèles
\providecommand{\blank}[1][]{\underline{\hspace{1.6em}}}
\newtcolorbox{exemplebox}[1][]{popbase,colframe=poporange,before upper={\popboxhead{Exemple}{poporange}{#1}}}
\newtcolorbox{definition}[1][]{popbase,colframe=popblue,before upper={\popboxhead{Définition}{popblue}{#1}}}
\newtcolorbox{propriete}[1][]{popbase,colframe=poppurple,before upper={\popboxhead{Propriété}{poppurple}{#1}}}
\newtcolorbox{methode}[1][]{popbase,colframe=popgold,before upper={\popboxhead{Méthode}{popgold}{#1}}}
\newtcolorbox{attention}[1][]{popbase,colframe=poppink,before upper={\popboxhead{Attention}{poppink}{#1}}}
\newtcolorbox{experience}[1][]{popbase,colframe=popblue,colback=popblueL,before upper={\popboxhead{Expérience}{popblue}{#1}}}
% « Le savais-tu ? » : carte violette pleine (contexte, culture, Cameroun)
\newtcolorbox{savaistu}[1][]{popbase,colback=poppurple,colframe=popink,
  fontupper=\popsemi\fontsize{10.4}{14.2}\selectfont\color{white},
  before upper={\poppill{Le savais-tu\,?}{white}{poppurple}\ifx\relax#1\relax\else\hspace{6pt}{\popxbold\color{white}#1}\fi\par\vspace{7pt}}}
% Exercice résolu : énoncé en haut, \tcblower puis la solution sur fond crème
\newcounter{popexo}\newcounter{popexr}
\newtcolorbox{exoresolu}[1][]{popbase,colframe=poporange,colbacklower=popcream,
  segmentation style={popink,line width=1.2pt,dashed},
  before upper={\stepcounter{popexr}\popboxhead{Exercice résolu \thepopexr}{poporange}{#1}},
  before lower={\poppill{Solution}{popink}{popcream}\par\vspace{6pt}}}
% Exercice (non résolu en place) — la correction est dans la partie Corrigés
\newtcolorbox{exercice}[1][]{popbase,colframe=popink,
  before upper={\stepcounter{popexo}\popboxhead{Exercice \thepopexo}{popink}{#1}}}
% Corrigé
\newtcolorbox{corrige}[1][]{popbase,colframe=popgreen,colback=popgreenL,
  before upper={\popboxhead{Corrigé}{popgreen}{#1}}}
% Objectifs / prérequis / bilan
\newtcolorbox{objectifs}{popbase,colframe=popink,colback=poporangeL,
  before upper={\popboxhead{Objectifs de la leçon}{popink}{}}}
\newtcolorbox{prerequis}{popbase,colframe=popink,colback=popblueL,
  before upper={\popboxhead{Prérequis}{popblue}{}}}
\newtcolorbox{bilan}{popbase,colframe=popink,colback=white,boxrule=2.8pt,
  before upper={\popboxhead{Fiche bilan}{poporange}{L'essentiel à connaître pour l'examen}}}
% Figure encadrée avec légende
\newtcolorbox{popfigure}[1][]{enhanced,breakable=false,colback=white,colframe=popline,boxrule=1.4pt,arc=8pt,
  left=10pt,right=10pt,top=10pt,bottom=8pt,before skip=10pt,after skip=12pt,halign=center,
  lower separated=false,halign lower=center,
  fontlower=\popsemi\fontsize{9}{11.5}\selectfont\color{popmuted},
  #1}
\newcommand{\legende}[1]{\par\vspace{6pt}{\popsemi\fontsize{9}{11.5}\selectfont\color{popmuted}#1}}

% ============================================================
% Signature OnBuch+
% ============================================================
\newcommand{\poplogoinline}[1]{{\poplogo\fontsize{#1}{#1}\selectfont\color{popink}OnBuch\color{poporange}+}}
\newcommand{\signatureonbuch}{%
  \begin{tcolorbox}[enhanced,colback=popink,colframe=popink,boxrule=2.2pt,arc=10pt,
      drop shadow=poporange,left=14pt,right=14pt,top=10pt,bottom=10pt,before skip=0pt,after skip=0pt]
    \tikz[baseline=-0.7ex]\node[circle,fill=popgreen,draw=popcream,line width=1.3pt,minimum size=22pt,inner sep=0]{\popcheck{white}};%
    \hspace{9pt}\begin{minipage}[c]{0.62\linewidth}
      {\popxbold\fontsize{12}{13}\selectfont\color{popcream}Signé\ \textbullet\ {\poplogo OnBuch\color{poporange}+}}\par\vspace{2pt}
      {\raggedright\popsemi\fontsize{8}{10.5}\selectfont\color{popline}\MakeUppercase{Équipe pédagogique OnBuch+}\\\MakeUppercase{Conforme au programme officiel MINESEC}\par}
    \end{minipage}\hfill
    {\poplogo\fontsize{22}{22}\selectfont\color{popcream}OnBuch\color{poporange}+}
  \end{tcolorbox}%
}

% ============================================================
% Page de garde
% #1 titre · #2 matière · #3 niveau · #4 module · #5 n° leçon · #6 durée ·
% #7 description / objectifs (texte ou itemize)
% ============================================================
\newcommand{\pagegarde}[7]{%
  \thispagestyle{empty}
  \newgeometry{a4paper,margin=1.7cm,top=1.6cm,bottom=1.4cm}%
  \begin{tikzpicture}[remember picture,overlay]
    \fill[poporange!18] ([xshift=-3.2cm,yshift=-3.6cm]current page.north east) circle (4.6cm);
    \fill[poppurple!14] ([xshift=2.4cm,yshift=7.6cm]current page.south west) circle (3.8cm);
  \end{tikzpicture}%
  {\poplogo\fontsize{30}{30}\selectfont\color{popink}OnBuch\color{poporange}+}\hfill
  \raisebox{0.6ex}{\poppill{Cours signé \textbullet\ Premium}{popink}{popcream}}\par
  \vspace{1pt}\popsquiggle{poppurple}\par
  \vspace{8pt}
  \begin{tcolorbox}[enhanced,colback=poporange,colframe=popink,boxrule=2.8pt,arc=13pt,
      drop shadow=popink,left=18pt,right=18pt,top=12pt,bottom=13pt,after skip=10pt]
    \poppill{#2 \textbullet\ #3}{popink}{popcream}\hspace{5pt}\poppill{#4}{white}{popink}\par\vspace{8pt}
    {\popdisplay\fontsize{14}{14}\selectfont\color{popink}LEÇON #5}\par\vspace{4pt}
    {\raggedright\hyphenpenalty=10000\exhyphenpenalty=10000\popdisplay\fontsize{25}{27}\selectfont\color{popcream}\MakeUppercase{#1}\par}
    \vspace{8pt}
    {\popxbold\fontsize{9.5}{9.5}\selectfont\color{popink}\MakeUppercase{Durée conseillée : #6}}
  \end{tcolorbox}
  \begin{tcolorbox}[enhanced,colback=white,colframe=popink,boxrule=2.2pt,arc=10pt,
      drop shadow=popink,left=16pt,right=16pt,top=10pt,bottom=10pt,after skip=8pt]
    \poppill{Dans cette leçon}{poppurple}{white}\par\vspace{5pt}
    \popsemi\fontsize{9.3}{12.6}\selectfont\color{popink}#7
  \end{tcolorbox}
  \noindent\begin{minipage}[t]{0.487\linewidth}
    \begin{tcolorbox}[enhanced,colback=popgreen,colframe=popink,boxrule=2.2pt,arc=10pt,
        drop shadow=popink,left=11pt,right=11pt,top=10pt,bottom=10pt,height=3.1cm,valign=center]
      \tcbox[colback=white,colframe=popink,boxrule=1.3pt,arc=5pt,boxsep=0pt,nobeforeafter,
        left=3pt,right=3pt,top=3pt,bottom=3pt,valign=center]{\qrcode[height=1.9cm]{https://play.google.com/store/apps/details?id=com.luvvix.onbuch&hl=fr}}%
      \hspace{8pt}\begin{minipage}[c]{0.5\linewidth}
        {\popdisplay\fontsize{11}{12.5}\selectfont\color{white}\MakeUppercase{Télécharge OnBuch}}\par\vspace{4pt}
        {\raggedright\popsemi\fontsize{8.4}{11}\selectfont\color{white}Gratuit \textbullet\ Google Play\\Tuteur IA, annales, concours, résultats\par}
      \end{minipage}
    \end{tcolorbox}
  \end{minipage}\hfill
  \begin{minipage}[t]{0.487\linewidth}
    \begin{tcolorbox}[enhanced,colback=poppurple,colframe=popink,boxrule=2.2pt,arc=10pt,
        drop shadow=popink,left=11pt,right=11pt,top=10pt,bottom=10pt,height=3.1cm,valign=center]
      \tikz[baseline=-0.7ex]\node[circle,fill=white,draw=popink,line width=1.3pt,minimum size=1.6cm,inner sep=0]{\popdisplay\fontsize{24}{24}\selectfont\color{poppurple}+};%
      \hspace{8pt}\begin{minipage}[c]{0.6\linewidth}
        {\popdisplay\fontsize{11}{12.5}\selectfont\color{white}\MakeUppercase{Passe à OnBuch+}}\par\vspace{4pt}
        {\raggedright\popsemi\fontsize{8.4}{11}\selectfont\color{white}Tous les cours signés, Tuteur IA illimité, fiches PDF — onglet OnBuch+ de l'app\par}
      \end{minipage}
    \end{tcolorbox}
  \end{minipage}
  \vspace{8pt}
  \popcontenu
  \vfill
  \signatureonbuch
  \restoregeometry
  \newpage
}

% Rangée « Ce cours contient » (page de garde)
\newcommand{\poptile}[3]{% #1 couleur #2 gros texte #3 légende
  \begin{tcolorbox}[enhanced,colback=#1,colframe=popink,boxrule=2pt,arc=9pt,shadow={2.5pt}{-2.5pt}{0pt}{popink},
      left=6pt,right=6pt,top=9pt,bottom=9pt,halign=center,valign=center,height=1.75cm,width=\linewidth,nobeforeafter]
    {\popdisplay\fontsize{12.5}{13}\selectfont\color{white}\MakeUppercase{#2}}\par\vspace{3pt}
    {\centering\popsemi\fontsize{7.8}{9.5}\selectfont\color{white}#3\par}
  \end{tcolorbox}}
\newcommand{\popcontenu}{%
  \par\noindent{\popxboldls\fontsize{7.5}{7.5}\selectfont\color{popmuted}CE COURS SIGNÉ CONTIENT}\par\vspace{7pt}
  \noindent\begin{minipage}[t]{0.235\linewidth}\poptile{popblue}{Cours}{complet, illustré, pas à pas}\end{minipage}\hfill
  \begin{minipage}[t]{0.235\linewidth}\poptile{poporange}{Méthodes}{et exercices résolus}\end{minipage}\hfill
  \begin{minipage}[t]{0.235\linewidth}\poptile{poppink}{Exercices}{type examen + corrigés}\end{minipage}\hfill
  \begin{minipage}[t]{0.235\linewidth}\poptile{popgreen}{Fiche bilan}{l'essentiel à retenir}\end{minipage}\par}

% Sommaire
\newtcolorbox{sommaire}{enhanced,colback=white,colframe=popink,boxrule=2.2pt,arc=10pt,
  drop shadow=popink,left=18pt,right=18pt,top=13pt,bottom=13pt,before skip=6pt,after skip=20pt,
  before upper={\poppill{Sommaire}{poppurple}{white}\par\vspace{9pt}\popsemi\fontsize{10.6}{14.5}\selectfont\color{popink}}}

% Signature de fin de cours — poussée tout en bas de la dernière page
\newcommand{\findecours}{%
  \par\vfill\nopagebreak
  \begin{center}\popsquiggle{poporange}\end{center}\vspace{4pt}
  \signatureonbuch
}
\AtBeginDocument{\def\llbracket{\mathopen{[\mkern-3.5mu[}}\def\rrbracket{\mathclose{]\mkern-3.5mu]}}} % intervalles d'entiers (glyphe ⟦ absent de la police)
% Glyphes absents de Fira Math : redéfinis à partir de glyphes disponibles
\AtBeginDocument{%
  \def\setminus{\mathbin{\backslash}}\let\smallsetminus\setminus
  \def\vdots{\mathord{\vbox{\baselineskip3.2pt\lineskiplimit0pt\kern4pt\hbox{.}\hbox{.}\hbox{.}}}}%
  \def\ddots{\mathinner{\mkern1mu\raise6.5pt\hbox{.}\mkern2mu\raise3.6pt\hbox{.}\mkern2mu\raise0.7pt\hbox{.}\mkern1mu}}%
  \def\triangle{\mathord{\scalebox{0.9}{$\symup{\Delta}$}}}%
  \def\nabla{\mathord{\raisebox{\depth}{\scalebox{1}[-1]{$\symup{\Delta}$}}}}%
  \def\checkmark{\popcheck{popdark}}%
  \def\star{\mathbin{\ast}}\def\bigstar{\mathord{\ast}}%
  \def\diamond{\mathbin{\scalebox{0.6}{$\Diamond$}}}%
  \def\bigcup{\mathop{\mathchoice{\raisebox{-0.3em}{\scalebox{1.7}{$\cup$}}}{\scalebox{1.25}{$\cup$}}{\cup}{\cup}}\displaylimits}%
  \def\bigcap{\mathop{\mathchoice{\raisebox{-0.3em}{\scalebox{1.7}{$\cap$}}}{\scalebox{1.25}{$\cap$}}{\cap}{\cap}}\displaylimits}%
  \def\lesseqqgtr{\mathrel{\lessgtr}}\def\bigoplus{\mathop{\oplus}\displaylimits}%
}
\providecommand{\ssi}{si et seulement si} % abréviation fréquente
\AtBeginDocument{\providecommand{\wideparen}{\overparen}} % arc de cercle

\begin{document}

\pagegarde{Leçon 23 — Expression écrite : la démocratie}{Chinois}{Tle A}{Module 3 : Citoyenneté et ouverture au monde}{ZH23}{2 h}{Cette leçon d'expression écrite apprend à rédiger un texte simple en chinois sur le thème de la démocratie : écriture correcte des caractères clés, organisation du texte avec des connecteurs, et traduction dans les deux sens (thème et version). Elle prépare directement aux épreuves écrites du baccalauréat.
\vspace{4pt}
\begin{multicols}{2}\small
\begin{itemize}
\item À la fin de cette leçon, je sais écrire correctement les caractères du thème de la démocratie (traits, radicaux, caractères proches)
\item À la fin de cette leçon, je sais distinguer les caractères proches faciles à confondre (民/氏, 主/王, 权/劝)
\item À la fin de cette leçon, je sais analyser la structure d'un texte argumentatif simple en chinois (introduction, arguments, conclusion)
\item À la fin de cette leçon, je sais utiliser les connecteurs utiles : 首先、其次、然后、所以、但是、因为
\item À la fin de cette leçon, je sais produire un texte simple de 5 à 8 phrases sur la démocratie
\item À la fin de cette leçon, je sais traduire du français vers le chinois des phrases simples sur ce thème (thème)
\end{itemize}\end{multicols}}

\begin{sommaire}
\begin{multicols}{2}
\begin{enumerate}
\item Modèle de texte commenté
\item Lexique et écriture des caractères du thème
\item Structure du texte et connecteurs
\item Thème : français → chinois
\item Version : chinois → français
\item Production guidée et grille d'évaluation
\item Activité d'intégration
\item Méthodes et exercices résolus
\item Exercices
\item Corrigés des exercices
\item Fiche bilan
\end{enumerate}
\end{multicols}
\end{sommaire}

\begin{objectifs}
\begin{itemize}
\item À la fin de cette leçon, je sais écrire correctement les caractères du thème de la démocratie (traits, radicaux, caractères proches)
\item À la fin de cette leçon, je sais distinguer les caractères proches faciles à confondre (民/氏, 主/王, 权/劝)
\item À la fin de cette leçon, je sais analyser la structure d'un texte argumentatif simple en chinois (introduction, arguments, conclusion)
\item À la fin de cette leçon, je sais utiliser les connecteurs utiles : 首先、其次、然后、所以、但是、因为
\item À la fin de cette leçon, je sais produire un texte simple de 5 à 8 phrases sur la démocratie
\item À la fin de cette leçon, je sais traduire du français vers le chinois des phrases simples sur ce thème (thème)
\item À la fin de cette leçon, je sais traduire du chinois vers le français un court texte sur ce thème (version)
\item À la fin de cette leçon, je sais m'auto-évaluer à l'aide d'une grille d'évaluation type examen
\end{itemize}
\end{objectifs}

\begin{prerequis}
\begin{itemize}
\item Structures de phrase de base : sujet + verbe + complément, 是, 有, 很 (Première)
\item Lexique de la citoyenneté vu dans le Module 3 (国家、人民、社会)
\item Connecteurs simples : 和、也、因为……所以…… (Première)
\item Radicaux fréquents et ordre général des traits (de haut en bas, de gauche à droite)
\item Négation 不 / 没 et particule 了
\end{itemize}
\end{prerequis}

\newpage

\coursec{Modèle de texte commenté}

\courssub{Situation d'accroche et problématique}

À Yaoundé, les élèves du club débat préparent une conférence interscolaire avec un lycée partenaire de Pékin sur le thème « la démocratie dans le monde ». Chaque délégation doit remettre un court texte écrit dans la langue de l'autre : les Camerounais rédigeront en chinois, les Chinois en français. Comment présenter clairement, en quelques phrases bien construites, des notions comme le vote, les droits du citoyen ou la liberté d'expression ?

C'est toute la difficulté de l'expression écrite en chinois : il ne suffit pas de connaître des mots, il faut les organiser selon une logique claire, avec une thèse, des arguments reliés par des connecteurs, et une conclusion. Cette leçon vous donne les caractères, les connecteurs et la méthode pour y parvenir. Nous partons d'un texte modèle court, que nous allons lire, observer, décortiquer, puis imiter.

\courssub{Lecture du texte modèle}

Lisez le texte suivant à voix haute \textbf{deux fois}, en respectant les tons, puis recopiez-le soigneusement dans votre cahier. La copie fait partie de l'apprentissage : elle fixe les caractères et l'ordre des mots.

\begin{textebox}[Texte modèle : 民主很重要 (La démocratie est importante)]
民主很重要。首先，民主让人民选择自己的领导人。其次，在民主国家，公民有言论自由，也有投票的权利。但是，民主也需要每个公民负责任。所以，我们应该学习民主的知识，参加社会活动。\\
\emph{Mínzhǔ hěn zhòngyào. Shǒuxiān, mínzhǔ ràng rénmín xuǎnzé zìjǐ de lǐngdǎorén. Qícì, zài mínzhǔ guójiā, gōngmín yǒu yánlùn zìyóu, yě yǒu tóupiào de quánlì. Dànshì, mínzhǔ yě xūyào měi ge gōngmín fù zérèn. Suǒyǐ, wǒmen yīnggāi xuéxí mínzhǔ de zhīshi, cānjiā shèhuì huódòng.}\\
Traduction : La démocratie est très importante. D'abord, la démocratie permet au peuple de choisir ses dirigeants. Ensuite, dans un pays démocratique, les citoyens ont la liberté d'expression et le droit de vote. Mais la démocratie exige aussi que chaque citoyen soit responsable. C'est pourquoi nous devons apprendre les connaissances sur la démocratie et participer aux activités sociales.
\end{textebox}

\textbf{Glossaire du texte modèle}

\begin{center}
\begin{tabularx}{\linewidth}{|X|X|X|X|}
\hline
\rowcolor{popblueL} Caractères & Pinyin & Nature & Traduction \\
\hline
民主 & mínzhǔ & nom / adj. & démocratie ; démocratique \\
\hline
重要 & zhòngyào & adjectif & important \\
\hline
首先 & shǒuxiān & connecteur & d'abord, premièrement \\
\hline
让 & ràng & verbe & faire, permettre, laisser \\
\hline
人民 & rénmín & nom & le peuple \\
\hline
选择 & xuǎnzé & verbe & choisir \\
\hline
自己 & zìjǐ & pronom & soi-même, son propre \\
\hline
领导人 & lǐngdǎorén & nom & dirigeant, leader \\
\hline
其次 & qícì & connecteur & ensuite, deuxièmement \\
\hline
国家 & guójiā & nom & pays, État \\
\hline
公民 & gōngmín & nom & citoyen \\
\hline
言论自由 & yánlùn zìyóu & locution nominale & liberté d'expression \\
\hline
投票 & tóupiào & verbe / nom & voter ; vote \\
\hline
权利 & quánlì & nom & droit (d'un citoyen) \\
\hline
但是 & dànshì & connecteur & mais, cependant \\
\hline
每 & měi & déterminant & chaque \\
\hline
负责任 & fù zérèn & locution verbale & être responsable, assumer ses responsabilités \\
\hline
所以 & suǒyǐ & connecteur & donc, c'est pourquoi \\
\hline
应该 & yīnggāi & verbe modal & devoir (conseil) \\
\hline
知识 & zhīshi & nom & connaissances, savoir \\
\hline
参加 & cānjiā & verbe & participer à \\
\hline
社会活动 & shèhuì huódòng & locution nominale & activités sociales \\
\hline
\end{tabularx}
\end{center}

\textbf{Questions de compréhension} (répondez en français, puis essayez en chinois) :
\begin{enumerate}
\item Quelle est l'idée principale (la thèse) du texte ?
\item Combien d'arguments l'auteur donne-t-il en faveur de la démocratie ?
\item Quelle nuance (quelle limite) est introduite par \zh{但是} ?
\item Que propose l'auteur dans la conclusion ?
\end{enumerate}

\courssub{Repérage de la structure en trois parties}

Un bon petit texte argumentatif chinois fonctionne exactement comme un bon paragraphe français : une \textbf{thèse}, des \textbf{arguments} reliés par des connecteurs, une \textbf{conclusion}. Observons le texte modèle :

\begin{itemize}
\item \textbf{Thèse} (une phrase) : \zh{民主很重要。} — l'auteur annonce son idée générale.
\item \textbf{Arguments} : \zh{首先} (d'abord) introduit l'argument 1 : le peuple choisit ses dirigeants ; \zh{其次} (ensuite) introduit l'argument 2 : liberté d'expression et droit de vote ; \zh{但是} (mais) introduit une \textbf{nuance} : la démocratie exige la responsabilité de chacun.
\item \textbf{Conclusion} : \zh{所以} (donc) introduit la conséquence pratique : apprendre et participer.
\end{itemize}

\begin{popfigure}
\begin{tikzpicture}[node distance=0.45cm and 0.6cm,
  bloc/.style={rectangle, rounded corners=3pt, draw=popblue, fill=popblueL, align=center, font=\small, inner sep=5pt, text width=3.4cm},
  fleche/.style={-{Stealth[length=2.5mm]}, thick, popdark}]
\node[bloc, fill=popgoldL, draw=popgold, align=center] (these){\textbf{Thèse}\\民主很重要。};
\node[bloc, below=of these, align=center] (arg1){\textbf{Argument 1} (首先)\\人民选择领导人};
\node[bloc, below=of arg1, align=center] (arg2){\textbf{Argument 2} (其次)\\言论自由、投票的权利};
\node[bloc, fill=poporangeL, draw=poporange, below=of arg2, align=center] (nuance){\textbf{Nuance} (但是)\\每个公民负责任};
\node[bloc, fill=popgreenL, draw=popgreen, below=of nuance, align=center] (concl){\textbf{Conclusion} (所以)\\学习民主知识、参加社会活动};
\draw[fleche] (these) -- (arg1);
\draw[fleche] (arg1) -- (arg2);
\draw[fleche] (arg2) -- (nuance);
\draw[fleche] (nuance) -- (concl);
\end{tikzpicture}
\legende{Organigramme de la structure argumentative du texte modèle : thèse, arguments, nuance, conclusion.}
\end{popfigure}

\begin{methode}[Analyser un texte argumentatif chinois en 4 étapes]
\begin{enumerate}
\item \textbf{Trouver la phrase de thèse} : c'est souvent la première phrase, courte, avec un adjectif d'opinion comme \zh{重要} (important), \zh{好} (bon), \zh{有用} (utile).
\item \textbf{Entourer les connecteurs} : \zh{首先} (d'abord), \zh{其次} (ensuite), \zh{还有} (en plus), \zh{但是} (mais), \zh{所以} (donc), \zh{因为} (parce que). Ils balisent le raisonnement.
\item \textbf{Numéroter les arguments} : chaque connecteur d'énumération (首先, 其次, 还有…) annonce un argument ; écrivez 1, 2, 3 dans la marge.
\item \textbf{Repérer la conclusion} : elle est presque toujours introduite par \zh{所以} ou \zh{因此} (c'est pourquoi) et contient souvent un conseil avec \zh{应该} (il faut, on devrait).
\end{enumerate}
\end{methode}

\courssub{Écriture des caractères : traits, radicaux et caractères proches}

Maîtriser l'écriture des sept mots-clés du thème est indispensable pour l'examen écrit. Respectez l'ordre des traits (haut → bas, gauche → droite, horizontal avant vertical, trait central avant les ailes) et identifiez le radical (clé de dictionnaire).

\begin{center}
\begin{tabularx}{\linewidth}{|X|X|X|X|X|}
\hline
\rowcolor{popblueL} Caractère & Pinyin & Radical (Clé) & Nb traits & Ordre des traits principaux / Caractères proches \\
\hline
民 & mín & 氏 (shì) / 乙 & 5 & 丿 乀 丿 丶 乙 (haut → bas) \quad Proche : \zh{盲} (máng, aveugle) \\
\hline
主 & zhǔ & 丶 (zhǔ) & 5 & 丶 一 丨 一 一 (le point initial est le radical) \quad Proche : \zh{住} (zhù, habiter), \zh{注} (zhù, verser) \\
\hline
公 & gōng & 八 (bā) & 4 & 八 + 厶 (sī, privé) \quad Proche : \zh{习} (xí, réviser), \zh{功} (gōng, mérite) \\
\hline
权 & quán & 木 (mù) & 6 & Version simplifiée de \zh{權} : 木 + 乛 (forme abrégée) \quad Proche : \zh{棋} (qí, échecs) \\
\hline
利 & lì & 刀 (刂 dāo) & 7 & 禾 (hé, grain) + 刂 (couteau) = « tranchant comme une faucille » \quad Proche : \zh{削} (xiāo, éplucher) \\
\hline
自 & zì & 自 (zì) & 6 & 目 (œil) + 丶 + 丨 + 八 (anciennement « nez ») \quad Proche : \zh{白} (bái, blanc), \zh{目} (mù, œil) \\
\hline
由 & yóu & 田 (tián) & 5 & 田 + 丨 (traversant) \quad Proche : \zh{甲} (jiǎ, carapace), \zh{申} (shēn, étendre) \\
\hline
投 & tóu & 手 (扌 shǒu) & 7 & 扌 + 秀 (xiù, élégant) \quad Proche : \zh{透} (tòu, traverser), \zh{秀} (xiù) \\
\hline
票 & piào & 示 (礻 shì) & 7 & 礻 + 示 (shì) \quad Proche : \zh{祭} (jì, sacrifier), \zh{福} (fú, bonheur) \\
\hline
责 & zé & 贝 (bèi) & 8 & 贝 (coquillage/monnaie) + 朁 (phonétique) \quad Proche : \zh{贷} (dài, prêt), \zh{费} (fèi, frais) \\
\hline
任 & rèn & 人 (亻 rén) & 6 & 亻 + 壬 (rén) \quad Proche : \zh{认} (rèn, reconnaître), \zh{妊} (rèn, grossesse) \\
\hline
\end{tabularx}
\end{center}

\begin{aretenir}[Règles d'or pour l'écriture]
\begin{itemize}
\item \textbf{民} : le trait final \zh{乙} (shù gōu) part du centre et descend vers la gauche.
\item \textbf{主} : le premier trait est le point \zh{丶} (radical), pas un trait horizontal.
\item \textbf{权} : version simplifiée, ne pas écrire l'ancienne forme \zh{權} (22 traits) à l'examen.
\item \textbf{利} : le radical \zh{刂} (dao) s'écrit à droite : verticale puis crochet vertical \zh{㇆}.
\item \textbf{责} : le bas \zh{朁} n'est pas \zh{責} (traditionnel) ; en simplifié, le trait horizontal du bas ne traverse pas.
\end{itemize}
\end{aretenir}

\courssub{Les mots-clés du thème : observation et étymologie}

Soulignez dans le texte les mots-clés du thème de la démocratie : \zh{民主} (démocratie), \zh{人民} (peuple), \zh{公民} (citoyen), \zh{权利} (droit), \zh{自由} (liberté), \zh{投票} (vote), \zh{责任} (responsabilité). Ce sont les sept piliers lexicaux de cette leçon : tout texte que vous rédigerez sur ce thème les utilisera.

\begin{savaistu}[民主, « le peuple maître » : comme en grec !]
\zh{民主} (mínzhǔ) signifie littéralement « le peuple (\zh{民}, mín) maître (\zh{主}, zhǔ) ». C'est exactement la même construction que le mot français « démocratie », formé du grec \emph{dêmos} (peuple) et \emph{kratos} (pouvoir, commandement) ! Le mot chinois a été créé au XIX\textsuperscript{e} siècle pour traduire le concept occidental, en assemblant deux caractères anciens. Le caractère \zh{主} (maître, seigneur) se retrouve aussi dans \zh{主席} (président, litt. « maître de séance ») et \zh{主人} (maître de maison, hôte). Retenez la famille : \zh{民} se lit \emph{mín} (2\textsuperscript{e} ton) et \zh{主} se lit \emph{zhǔ} (3\textsuperscript{e} ton).
\end{savaistu}

\begin{attention}[公民 ou 人民 ? Ne confondez pas « citoyen » et « peuple »]
Les francophones traduisent souvent les deux par « peuple ». Or : \zh{人民} (rénmín) désigne le \textbf{peuple} comme collectivité (« le peuple camerounais »), tandis que \zh{公民} (gōngmín) désigne le \textbf{citoyen} comme individu doté de droits et de devoirs. On dit \zh{每个公民} (chaque citoyen) mais jamais \zh{每个人民}. De même, attention à \zh{权利} (quánlì, droit subjectif : le droit de vote) qu'il ne faut pas confondre avec son homophone \zh{权力} (quánlì, pouvoir, autorité) : mêmes sons, sens différents !
\end{attention}

\courssub{Connecteurs logiques pour l'argumentation}

Pour structurer votre texte, vous devez maîtriser les connecteurs (连词, liáncí). Ils se placent généralement en \textbf{début de phrase} (sujet + connecteur + prédicat) ou \textbf{après le sujet} (sujet + connecteur + prédicat).

\begin{center}
\begin{tabularx}{\linewidth}{|X|X|X|X|}
\hline
\rowcolor{popblueL} Fonction & Chinois & Pinyin & Français \\
\hline
\multicolumn{4}{|c|}{\textbf{Énumération / Ordre}} \\
\hline
Premièrement & 首先 & shǒuxiān & D'abord, en premier lieu \\
\hline
Deuxièmement & 其次 & qícì & Ensuite, en second lieu \\
\hline
De plus / En outre & 还有 / 此外 & hái yǒu / qí wài & Aussi, de plus \\
\hline
Enfin & 最后 & zuìhòu & Enfin, pour finir \\
\hline
\multicolumn{4}{|c|}{\textbf{Opposition / Concession}} \\
\hline
Mais & 但是 & dànshì & Mais, cependant \\
\hline
Cependant (formel) & 然而 & rán'ér & Cependant, toutefois \\
\hline
Même si & 虽然…但是… & suīrán… dànshì… & Bien que… pourtant… \\
\hline
\multicolumn{4}{|c|}{\textbf{Cause / Conséquence}} \\
\hline
Parce que & 因为 & yīnwèi & Parce que (souvent avec 所以) \\
\hline
Donc / C'est pourquoi & 所以 & suǒyǐ & Donc, par conséquent \\
\hline
C'est pourquoi (formel) & 因此 & yīncǐ & C'est pourquoi, par suite \\
\hline
\multicolumn{4}{|c|}{\textbf{But / Condition}} \\
\hline
Afin que & 为了 & wèile & Afin de, pour (but) \\
\hline
Si & 如果 / 要是 & rúguǒ / yàoshi & Si (condition) \\
\hline
\end{tabularx}
\end{center}

\begin{exemplebox}[Position des connecteurs : deux cas]
\textbf{Cas 1 : En tête de phrase (thème) — insiste sur la logique du discours.}\\
\zh{首先，我们要投票。} \emph{Shǒuxiān, wǒmen yào tóupiào.} (D'abord, nous devons voter.)\\[4pt]
\textbf{Cas 2 : Après le sujet — insiste sur l'agent de l'action.}\\
\zh{我们首先要投票。} \emph{Wǒmen shǒuxiān yào tóupiào.} (Nous, d'abord, devons voter.)\\
\cle{Conseil} : Pour l'examen écrit, privilégiez la \textbf{position en tête de phrase} : c'est plus formel et plus clair pour le correcteur.
\end{exemplebox}

\courssub{Observation grammaticale : 让, 应该, 需要}

Reprenons trois phrases du texte et observons-les avant d'énoncer les règles.

\begin{exemplebox}[Observer avant de comprendre]
\zh{民主让人民选择自己的领导人。}\\
\emph{Mínzhǔ ràng rénmín xuǎnzé zìjǐ de lǐngdǎorén.}\\
« La démocratie permet au peuple de choisir ses propres dirigeants. »\\[4pt]
\zh{我们应该学习民主的知识。}\\
\emph{Wǒmen yīnggāi xuéxí mínzhǔ de zhīshi.}\\
« Nous devons apprendre les connaissances sur la démocratie. »\\[4pt]
\zh{民主需要每个公民负责任。}\\
\emph{Mínzhǔ xūyào měi ge gōngmín fù zérèn.}\\
« La démocratie a besoin que chaque citoyen soit responsable. »
\end{exemplebox}

\begin{propriete}[Le verbe 让 (ràng) : « faire / permettre / laisser »]
Structure : \textbf{A + 让 + B + Verbe} = « A fait/permet que B fasse… ».\\
C'est une \emph{phrase à prédicat en série} (兼语句) : \zh{让} introduit un deuxième sujet (B) qui accomplit l'action finale. Pas de mot-outil entre \zh{让} et B.
\begin{itemize}
\item \zh{老师让我们写字。} \emph{Lǎoshī ràng wǒmen xiě zì.} « Le professeur nous fait écrire les caractères. »
\item \zh{投票让公民表达意见。} \emph{Tóupiào ràng gōngmín biǎodá yìjiàn.} « Le vote permet aux citoyens d'exprimer leur opinion. »
\item \zh{妈妈让我学汉语。} \emph{Māma ràng wǒ xué Hànyǔ.} « Maman me fait apprendre le chinois. »
\end{itemize}
En français, on traduit par « permettre à qqn de », « faire que », « laisser » selon le contexte.
\end{propriete}

\begin{propriete}[Le modal 应该 (yīnggāi) : « devoir » au sens du conseil]
Structure : \textbf{Sujet + 应该 + Verbe}. Il exprime ce qui est souhaitable, le devoir moral ou le conseil (« il faut, on devrait »), et non l'obligation stricte (contrairement à \zh{必须} bìxū).
\begin{itemize}
\item \zh{每个公民都应该投票。} \emph{Měi ge gōngmín dōu yīnggāi tóupiào.} « Chaque citoyen devrait voter. »
\item \zh{年轻人应该关心社会。} \emph{Niánqīngrén yīnggāi guānxīn shèhuì.} « Les jeunes devraient s'intéresser à la société. »
\end{itemize}
Négation : \zh{不应该} (ne devrait pas) : \zh{我们不应该忘记自己的责任。} « Nous ne devrions pas oublier nos responsabilités. »
\end{propriete}

\begin{propriete}[Le verbe 需要 (xūyào) : « avoir besoin de / exiger »]
Structure : \textbf{Sujet + 需要 + Nom / Complément}. Le sujet est souvent une chose abstraite (la démocratie, le travail, l'examen).
\begin{itemize}
\item \zh{学习汉语需要时间。} \emph{Xuéxí Hànyǔ xūyào shíjiān.} « Apprendre le chinois prend du temps / a besoin de temps. »
\item \zh{这个工作需要责任心。} \emph{Zhège gōngzuò xūyào zérènxīn.} « Ce travail exige du sens des responsabilités. »
\end{itemize}
\cle{Attention} : \zh{需要} ne s'utilise pas avec un verbe nu directement comme modal (on ne dit pas *\zh{我需要去}). On dit \zh{我需要去} (j'ai besoin d'y aller) mais le sujet est une personne. Si le sujet est une chose, le complément est un nom ou une proposition (avec 的).
\end{propriete}

\begin{attention}[Pièges fréquents des francophones]
\begin{enumerate}
\item \textbf{Pas de « de » après 让} : on écrit \zh{让人民选择} (litt. « faire peuple choisir »), sans mot-outil entre \zh{让} et le verbe. Ne calquez pas « permettre \emph{de} choisir ».
\item \textbf{Place du modal} : \zh{应该} se place \emph{avant} le verbe principal, jamais après : \zh{我们应该学习} et non \zh{我们学习应该}.
\item \textbf{Tons de 让} : \emph{ràng}, 4\textsuperscript{e} ton descendant. Un ton plat le ferait confondre avec d'autres syllabes ; articulez fermement.
\item \textbf{每个 + nom} : après \zh{每} (chaque), il faut le classificateur \zh{个} (ou le classificateur approprié) : \zh{每个公民}, jamais \zh{每公民}. Ajoutez souvent \zh{都} (tous) après le sujet pour distribuer : \zh{每个公民都应该…}.
\item \textbf{负责任} : c'est une locution V + O (verbe + objet). On ne dit pas *\zh{负责} seul pour « être responsable » dans ce contexte (ça veut dire « être en charge de »). La forme complète est \zh{负责任}.
\end{enumerate}
\end{attention}

\begin{center}
\begin{tabularx}{\linewidth}{|X|X|X|}
\hline
\rowcolor{popblueL} Structure & Exemple (caractères + pinyin) & Traduction \\
\hline
A + 让 + B + V & 民主让人民选择领导人。Mínzhǔ ràng rénmín xuǎnzé lǐngdǎorén. & La démocratie permet au peuple de choisir ses dirigeants. \\
\hline
S + 有 + Nom & 公民有言论自由。Gōngmín yǒu yánlùn zìyóu. & Les citoyens ont la liberté d'expression. \\
\hline
S + 应该 + V & 我们应该参加社会活动。Wǒmen yīnggāi cānjiā shèhuì huódòng. & Nous devons participer aux activités sociales. \\
\hline
S + 需要 + Nom/Prop & 民主需要公民负责任。Mínzhǔ xūyào gōngmín fù zérèn. & La démocratie exige que les citoyens soient responsables. \\
\hline
每 + Classif + Nom + 都 + V & 每个学生都喜欢中文。Měi ge xuéshēng dōu xǐhuan Zhōngwén. & Chaque élève aime le chinois. \\
\hline
\end{tabularx}
\end{center}

\begin{exoresolu}[Transformer une phrase avec 让]
Énoncé : traduisez en chinois « La liberté permet aux citoyens de parler », en utilisant \zh{让}, \zh{自由} (zìyóu, liberté), \zh{公民} (gōngmín, citoyen) et \zh{说话} (shuōhuà, parler).
\tcblower
Solution détaillée :
\begin{enumerate}
\item J'identifie la structure : A + 让 + B + Verbe. Ici A = \zh{自由} (la liberté), B = \zh{公民} (les citoyens), Verbe = \zh{说话} (parler).
\item J'assemble dans l'ordre chinois : \zh{自由让公民说话。}
\item Pinyin : \emph{Zìyóu ràng gōngmín shuōhuà.}
\item Vérification : pas de mot-outil entre \zh{让} et \zh{公民}, pas de « de » avant le verbe, ordre Sujet–Verbe–Complément respecté.
\end{enumerate}
Traduction littérale de contrôle : « Liberté fait citoyens parler » = « La liberté permet aux citoyens de parler ». Correct.
\end{exoresolu}

\begin{exercice}[Entraînement : structures clés]
\begin{enumerate}
\item Retrouvez dans le texte modèle la phrase de thèse et recopiez-la avec son pinyin.
\item Entourez les quatre connecteurs du texte et traduisez chacun d'eux.
\item Traduisez en chinois (thème) : « Le vote permet au peuple de choisir ses dirigeants » (utilisez \zh{让}).
\item Complétez avec \zh{应该} et \zh{都} : 每个公民 \_\_\_\_ \_\_\_\_ 负责任。 (Chaque citoyen devrait être responsable.)
\item Traduisez en chinois : « Apprendre la démocratie a besoin de temps. » (Utilisez \zh{需要}).
\end{enumerate}
\end{exercice}

\begin{corrige}[Exercice « Entraînement : structures clés »]
\begin{enumerate}
\item Thèse : \zh{民主很重要。} \emph{Mínzhǔ hěn zhòngyào.} « La démocratie est très importante. »
\item \zh{首先} (d'abord), \zh{其次} (ensuite), \zh{但是} (mais), \zh{所以} (donc).
\item \zh{投票让人民选择自己的领导人。} \emph{Tóupiào ràng rénmín xuǎnzé zìjǐ de lǐngdǎorén.}
\item \zh{每个公民都应该负责任。} \emph{Měi ge gōngmín dōu yīnggāi fù zérèn.}
\item \zh{学习民主需要时间。} \emph{Xuéxí mínzhǔ xūyào shíjiān.}
\end{enumerate}
\end{corrige}

\coursec{Thème : français → chinois}

\courssub{Consignes de traduction (Thème)}
\begin{enumerate}
\item Analysez la phrase française : identifiez le Sujet, le Verbe, les Compléments.
\item Cherchez la structure chinoise correspondante (SVO, 让, 应该, 需要, 比, 把, 被…).
\item Placez les compléments \textbf{avant} le verbe (Lieu, Temps, Manière) : Sujet + Temps + Lieu + Manière + Verbe + Complément.
\item Vérifiez : classificateurs après nombres/\zh{每}, \zh{的} (nom), \zh{得} (verbe), \zh{地} (adverbe), \zh{了} (aspect), tons, ponctuation chinoise.
\end{enumerate}

\begin{exercice}[Thème : phrases isolées (niveau HSK 3-4)]
Traduisez en chinois (caractères + pinyin).
\begin{enumerate}
\item La démocratie est très importante pour le peuple.
\item D'abord, les citoyens ont le droit de vote.
\item Ensuite, la liberté d'expression permet aux gens de parler.
\item Mais chaque citoyen doit être responsable.
\item C'est pourquoi nous devrions apprendre les connaissances sur la démocratie.
\item Le Cameroun est un pays démocratique. (Pays = \zh{喀麦隆} Kāmàilóng)
\item Les jeunes devraient participer aux activités sociales.
\item Choisir un dirigeant est un droit du peuple.
\item Sans responsabilité, la démocratie n'est pas bonne. (Indice : \zh{如果没有}… \zh{就}…)
\item Nous devrions respecter les droits des autres. (Autrui = \zh{别人} biérén)
\end{enumerate}
\end{exercice}

\begin{corrige}[Thème : phrases isolées]
\begin{enumerate}
\item \zh{民主对人民很重要。} \emph{Mínzhǔ duì rénmín hěn zhòngyào.} (Structure : S + 对 + O + Adj)
\item \zh{首先，公民有投票的权利。} \emph{Shǒuxiān, gōngmín yǒu tóupiào de quánlì.}
\item \zh{其次，言论自由让人们说话。} \emph{Qícì, yánlùn zìyóu ràng rénmen shuōhuà.} (\zh{人们} rénmen = les gens/people)
\item \zh{但是，每个公民都应该负责任。} \emph{Dànshì, měi ge gōngmín dōu yīnggāi fù zérèn.}
\item \zh{所以，我们应该学习民主的知识。} \emph{Suǒyǐ, wǒmen yīnggāi xuéxí mínzhǔ de zhīshi.}
\item \zh{喀麦隆是一个民主国家。} \emph{Kāmàilóng shì yí ge mínzhǔ guójiā.}
\item \zh{年轻人应该参加社会活动。} \emph{Niánqīngrén yīnggāi cānjiā shèhuì huódòng.}
\item \zh{选择领导人是人民的权利。} \emph{Xuǎnzé lǐngdǎorén shì rénmín de quánlì.} (Sujet = verbe + 得)
\item \zh{如果没有责任，民主就不好。} \emph{Rúguǒ méiyǒu zérèn, mínzhǔ jiù bù hǎo.} (Structure conditionnelle)
\item \zh{我们应该尊重别人的权利。} \emph{Wǒmen yīnggāi zūnzhòng biérén de quánlì.} (\zh{尊重} zūnzhòng = respecter)
\end{enumerate}
\end{corrige}

\begin{exercice}[Thème : texte court (niveau Bac)]
Traduisez le paragraphe suivant en chinois. Utilisez la structure Thèse - Arguments - Nuance - Conclusion et les connecteurs vus en classe.

« La démocratie est importante. D'abord, elle permet au peuple de choisir ses dirigeants. Ensuite, les citoyens ont la liberté d'expression et le droit de vote. Cependant, la démocratie a aussi besoin que chaque citoyen soit responsable. C'est pourquoi nous devrions étudier la démocratie et participer à la vie sociale. »
\end{exercice}

\begin{corrige}[Thème : texte court]
\zh{民主很重要。首先，民主让人民选择自己的领导人。其次，公民有言论自由，也有投票的权利。然而，民主也需要每个公民都负责任。因此，我们应该学习民主的知识，参加社会生活。}\\
\emph{Mínzhǔ hěn zhòngyào. Shǒuxiān, mínzhǔ ràng rénmín xuǎnzé zìjǐ de lǐngdǎorén. Qícì, gōngmín yǒu yánlùn zìyóu, yě yǒu tóupiào de quánlì. Rán'ér, mínzhǔ yě xūyào měi ge gōngmín dōu fù zérèn. Yīncǐ, wǒmen yīnggāi xuéxí mínzhǔ de zhīshi, cānjiā shèhuì shēnghuó.}

\cle{Notes de correction} :
\begin{itemize}
\item \zh{然而} (rán'ér) utilisé pour « Cependant » (plus formel que \zh{但是}).
\item \zh{都} ajouté après \zh{每个公民} pour la distribution.
\item \zh{因此} (yīncǐ) utilisé pour « C'est pourquoi » (plus écrit que \zh{所以}).
\item \zh{社会生活} (shèhuì shēnghuó) = vie sociale, synonyme de \zh{社会活动} ici.
\end{itemize}
\end{corrige}

\coursec{Version : chinois → français}

\courssub{Consignes de traduction (Version)}
\begin{enumerate}
\item Segmentez la phrase chinoise : identifiez le Sujet (souvent en tête), le Prédicat (verbe, adjectif, modal), les Compléments.
\item Repérez les mots-outils (\zh{的, 得, 地, 了, 过, 着, 把, 被, 让, 被, 连词}).
\item Traduisez littéralement mot à mot, puis reformulez en français correct (syntaxe SVO, conjugaison, accords).
\item Attention aux \textbf{faux amis} : \zh{权利} vs \zh{权力}, \zh{人民} vs \zh{公民}, \zh{民主} (adj/nom).
\end{enumerate}

\begin{exercice}[Version : phrases isolées]
Traduisez en français.
\begin{enumerate}
\item 在民主国家，公民有选举的权利。
\item 言论自由让每个人表达自己的意见。
\item 我们应该负责任地参加投票。
\item 如果没有人民的支持，民主就不能成功。
\item 领导人需要为人民负责任。
\item 每个学生都应该关心国家大事。
\item 选择权利是公民的自由。
\item 民主不仅需要权利，也需要责任。 (\zh{不仅…而且…} = non seulement… mais aussi…)
\end{enumerate}
\end{exercice}

\begin{corrige}[Version : phrases isolées]
\begin{enumerate}
\item Dans un pays démocratique, les citoyens ont le droit d'élire (le droit de vote). (\zh{选举} xuǎnjǔ = élire/élection)
\item La liberté d'expression permet à chacun d'exprimer sa propre opinion. (\zh{意见} yìjiàn = opinion)
\item Nous devrions participer au vote de manière responsable. (\zh{负责任地} fù zérèn de = adv. « responsablement »)
\item S'il n'y a pas le soutien du peuple, la démocratie ne peut pas réussir. (\zh{支持} zhīchí = soutien ; \zh{成功} chénggōng = réussir/réussite)
\item Les dirigeants doivent être responsables envers le peuple. (\zh{为…负责任} wèi… fù zérèn = être responsable envers…)
\item Chaque élève devrait se soucier des grandes affaires de l'État. (\zh{关心} guānxīn = se soucier de ; \zh{国家大事} guójiā dàshì = affaires d'État)
\item Le droit de choisir est la liberté du citoyen. (Attention : \zh{选择权利} = le droit de choisir ; \zh{选择的权利} = le droit de choix).
\item La démocratie a besoin non seulement de droits, mais aussi de responsabilités.
\end{enumerate}
\end{corrige}

\begin{exercice}[Version : texte court]
Traduisez le texte suivant en français.

\zh{我们学校也实行民主。首先，学生可以选择班长。其次，同学们有发表意见的自由。但是，自由不是随便做事。每个学生都应该遵守规则，负责任。所以，我们在班会上讨论问题，共同决定。}
\emph{Wǒmen xuéxiào yě shíxíng mínzhǔ. Shǒuxiān, xuésheng kěyǐ xuǎnzé bānzhǎng. Qícì, tóngxuémen yǒu fābiǎo yìjiàn de zìyóu. Dànshì, zìyóu bùshì suíbiàn zuò shì. Měi ge xuéshēng dōu yīnggāi zūnshǒu guīzé, fù zérèn. Suǒyǐ, wǒmen zài bānhuì shàng tǎolùn wèntí, gòngtóng juédìng.}
\end{exercice}

\begin{corrige}[Version : texte court]
Notre école pratique aussi la démocratie. D'abord, les élèves peuvent choisir le délégué de classe. Ensuite, les camarades ont la liberté d'exprimer leurs opinions. Mais la liberté n'est pas faire n'importe quoi. Chaque élève devrait respecter les règles et être responsable. C'est pourquoi, lors de la réunion de classe, nous discutons les problèmes et décidons ensemble.

\cle{Vocabulaire utile} : \zh{实行} (shíxíng, mettre en œuvre/pratiquer), \zh{班长} (bānzhǎng, délégué de classe), \zh{发表} (fābiǎo, publier/exprimer), \zh{随便} (suíbiàn, n'importe comment/à la légère), \zh{遵守} (zūnshǒu, respecter/observer), \zh{规则} (guīzé, règles), \zh{班会} (bānhuì, réunion de classe/heure de vie de classe), \zh{讨论} (tǎolùn, discuter/débattre), \zh{共同} (gòngtóng, ensemble/conjointement), \zh{决定} (juédìng, décider).
\end{corrige}

\coursec{Production guidée et grille d'évaluation}

\courssub{Sujet de production écrite (Type Bac)}

\begin{experience}[Rédaction : « La démocratie dans mon école »]
\textbf{Consigne} : Rédigez un paragraphe d'environ \textbf{80 à 100 caractères chinois} sur le thème : \zh{我学校的民主} (La démocratie dans mon école).

\textbf{Contraintes obligatoires} :
\begin{enumerate}
\item Structure en 4 étapes : \textbf{Thèse} → \textbf{Arguments} (min. 2 avec \zh{首先} / \zh{其次}) → \textbf{Nuance} (avec \zh{但是} / \zh{然而}) → \textbf{Conclusion} (avec \zh{所以} / \zh{因此} + \zh{应该}).
\item Vocabulaire imposé (au moins 5 sur 7) : \zh{民主}, \zh{学生} (élèves), \zh{选择}, \zh{权利}, \zh{自由}, \zh{责任}, \zh{参加}.
\item Structures imposées : au moins une phrase avec \zh{让}, une avec \zh{应该}, une avec \zh{每个…都…}.
\item Ponctuation chinoise correcte (。 ， ？ ！).
\item Écriture soignée des caractères (ordre des traits respecté).
\end{enumerate}

\textbf{Aide au démarrage (Plan détaillé)} :
\begin{itemize}
\item \textbf{Thèse} : \zh{我们学校的民主很重要。} / \zh{学校里也需要民主。}
\item \textbf{Arg 1} : \zh{首先，学生可以选择自己的班长。} (Utiliser \zh{让} : \zh{民主让学生选择…})
\item \textbf{Arg 2} : \zh{其次，同学们有发表意见的自由。}
\item \textbf{Nuance} : \zh{但是，自由不是随便做事。每个学生都应该遵守规则，负责任。}
\item \textbf{Conclusion} : \zh{所以，我们应该积极参加班会，共同建设学校。} (\zh{积极} jījí = activement ; \zh{建设} jiànshè = construire)
\end{itemize}
\end{experience}

\courssub{Grille d'évaluation (sur 20 points)}

Utilisez cette grille pour vous auto-évaluer ou évaluer un camarade (évaluation par les pairs).

\begin{center}
\begin{tabularx}{\linewidth}{|X|c|c|c|c|}
\hline
\rowcolor{popblueL} \textbf{Critères} & \textbf{4 pts (Excellent)} & \textbf{3 pts (Bon)} & \textbf{2 pts (Moyen)} & \textbf{1 pt (Insuffisant)} \\
\hline
\textbf{Structure argumentative} & 4 parties claires (Thèse, 2 Args, Nuance, Concl), connecteurs variés et justes. & 4 parties présentes, connecteurs corrects mais répétitifs. & Structure incomplète (manque nuance ou conclusion), connecteurs maladroits. & Pas de structure visible, pas de connecteurs. \\
\hline
\textbf{Vocabulaire du thème} & 5+ mots-clés utilisés correctement (\zh{民主, 人民/学生, 公民/同学, 权利, 自由, 投票/选择, 责任}). & 3-4 mots-clés utilisés correctement. & 1-2 mots-clés, ou emplois incorrects (ex: \zh{权力} pour \zh{权利}). & Vocabulaire hors sujet ou absent. \\
\hline
\textbf{Grammaire ciblée} & \zh{让}, \zh{应该}, \zh{需要}, \zh{每个…都…} employés sans faute. & Structures présentes avec 1 erreur mineure (place modal, classificateur). & Erreurs majeures sur structures cibles (ex: *\zh{让 de}, *\zh{应该在后面}). & Structures cibles absentes ou incompréhensibles. \\
\hline
\textbf{Caractères \& Pinyin} & Écriture soignée, ordre traits respecté, pinyin/ton juste si demandé. & 1-2 fautes de traits / tons. & Nombreuses fautes d'écriture / tons, lecture difficile. & Illisible, tons aléatoires. \\
\hline
\textbf{Cohérence \& Logique} & Argumentation fluide, progression logique, ton adapté. & Idées claires mais transitions abruptes. & Idées décousues, contradictions. & Incohérent, hors sujet. \\
\hline
\end{tabularx}
\end{center}

\begin{methode}[Auto-correction avant de rendre la copie]
\begin{enumerate}
\item \textbf{Relisez à voix haute} : entendez-vous les tons ? Les phrases sonnent-elles « chinoises » ?
\item \textbf{Cochez la structure} : Ai-je \zh{首先} ? \zh{其次} ? \zh{但是} ? \zh{所以} ?
\item \textbf{Vérifiez les pièges} : \zh{每个 + 班} ? \zh{都} présent ? \zh{让} sans \zh{de} ? \zh{应该} avant verbe ? \zh{权利} vs \zh{权力} ? \zh{负责任} (pas \zh{负责}) ?
\item \textbf{Comptez les caractères} : suis-je entre 80 et 100 ? (Si trop court : développez les arguments. Si trop long : condensez).
\item \textbf{Relisez l'écriture} : les radicaux sont-ils au bon endroit ? (\zh{扌} à gauche pour \zh{投}, \zh{礻} à gauche pour \zh{票}, \zh{贝} en bas pour \zh{责}, \zh{亻} à gauche pour \zh{任}) ?
\end{enumerate}
\end{methode}

\begin{aretenir}[L'essentiel de la leçon 23 — Expression écrite : La démocratie]
\begin{itemize}
\item \textbf{Structure type} : Thèse (\zh{很重要}) → Arguments (\zh{首先}, \zh{其次}) → Nuance (\zh{但是}, \zh{然而}) → Conclusion (\zh{所以}, \zh{因此} + \zh{应该}).
\item \textbf{7 mots-clés} : \zh{民主} (démocratie), \zh{人民} (peuple), \zh{公民} (citoyen), \zh{权利} (droit), \zh{自由} (liberté), \zh{投票} (vote), \zh{责任} (responsabilité).
\item \textbf{Grammaire} :
  \begin{itemize}
  \item \zh{A 让 B V} = A permet/fait que B fasse V (pas de \zh{de}).
  \item \zh{S 应该 V} = S devrait V (conseil, avant verbe).
  \item \zh{S 需要 N} = S a besoin de N / exige N.
  \item \zh{每 个 N 都 V} = Chaque N V (distribution).
  \end{itemize}
\item \textbf{Écriture} : Maîtriser traits, radicaux (\zh{扌, 礻, 贝, 亻, 木, 刂}) et différencier caractères proches (\zh{主/住}, \zh{权/棋}, \zh{责/贷}, \zh{任/认}).
\item \textbf{Culture} : \zh{民主} = « peuple maître » (\zh{民} + \zh{主}) ≡ grec \emph{dêmos} + \emph{kratos}. Distinction \zh{人民} (collectif) / \zh{公民} (individuel droits/devoirs).
\item \textbf{Examen} : Thème/Version sur phrases + Texte argumenté 80-100 caractères. Grille : Structure (4) + Vocab (4) + Grammaire (4) + Caractères (4) + Cohérence (4) = /20.
\end{itemize}
\end{aretenir}

\coursec{Lexique et écriture des caractères du thème}

Dans la section précédente, nous avons lu et commenté un modèle de texte sur la démocratie. Vous avez remarqué que certains mots revenaient sans cesse : \zh{民主}, \zh{公民}, \zh{权利}, \zh{投票}… Avant de rédiger vous-même, il faut donc \textbf{maîtriser ce lexique} : le comprendre, le prononcer, et surtout \textbf{l'écrire correctement}, caractère par caractère, trait par trait. C'est l'objectif de cette section.

\courssub{Observation : les mots du thème en contexte}

Commençons par observer un court texte qui utilise presque tout le vocabulaire de la leçon. Lisez-le, puis repérez les mots en gras dans le glossaire.

\begin{textebox}[Le vote des jeunes citoyens]
在我们的国家，民主很重要。\\
Zài wǒmen de guójiā, mínzhǔ hěn zhòngyào.\\
« Dans notre pays, la démocratie est très importante. »\\[4pt]
人民是公民，公民有权利，也有责任。\\
Rénmín shì gōngmín, gōngmín yǒu quánlì, yě yǒu zérèn.\\
« Le peuple est fait de citoyens ; les citoyens ont des droits et aussi des responsabilités. »\\[4pt]
每个公民都有投票的权利。选举的时候，我们去投票。\\
Měi ge gōngmín dōu yǒu tóupiào de quánlì. Xuǎnjǔ de shíhou, wǒmen qù tóupiào.\\
« Chaque citoyen a le droit de vote. Au moment des élections, nous allons voter. »\\[4pt]
自由的社会里，人民选择自己的未来。\\
Zìyóu de shèhuì lǐ, rénmín xuǎnzé zìjǐ de wèilái.\\
« Dans une société libre, le peuple choisit son propre avenir. »
\end{textebox}

\textbf{Questions de compréhension :}
\begin{enumerate}
\item Qui a des droits et des responsabilités ? (\emph{Réponse attendue : les citoyens, \zh{公民}.})
\item Que font les citoyens au moment des élections ? (\emph{Réponse : ils vont voter, \zh{去投票}.})
\item Que choisit le peuple dans une société libre ? (\emph{Réponse : son avenir, \zh{自己的未来}.})
\end{enumerate}

\courssub{Glossaire du thème : la démocratie}

Voici le vocabulaire officiel de la leçon. Apprenez-le en prononçant le pinyin à voix haute : en chinois, un mot mal prononcé (mauvais ton) devient un autre mot.

\begin{center}
\begin{tabularx}{\linewidth}{|X|X|X|X|}
\hline
\rowcolor{popblueL} Caractères & Pinyin & Nature & Traduction \\
\hline
民主 & mínzhǔ & n. & démocratie \\
\hline
人民 & rénmín & n. & peuple \\
\hline
公民 & gōngmín & n. & citoyen, citoyenne \\
\hline
权利 & quánlì & n. & droit (d'une personne) \\
\hline
自由 & zìyóu & n. / adj. & liberté ; libre \\
\hline
投票 & tóupiào & v. & voter \\
\hline
选举 & xuǎnjǔ & v. / n. & élire ; élection \\
\hline
责任 & zérèn & n. & responsabilité \\
\hline
国家 & guójiā & n. & pays, État, nation \\
\hline
社会 & shèhuì & n. & société \\
\hline
\end{tabularx}
\end{center}

\begin{definition}[公民 gōngmín : le citoyen]
Un \cle{公民} (gōngmín) est un \textbf{citoyen} : une personne qui appartient à un pays et qui y possède des \textbf{droits} (\zh{权利} quánlì) et des \textbf{devoirs} (\zh{义务} yìwù). Le mot se décompose ainsi : \zh{公} (gōng) = « public, commun » et \zh{民} (mín) = « peuple ». Le citoyen est donc littéralement « une personne du peuple dans l'espace public ». Au Cameroun comme en Chine, être citoyen, c'est avoir le droit de voter et le devoir de respecter les lois.
\end{definition}

\begin{exemplebox}[Le glossaire en phrases]
\zh{公民有投票的权利。}\\
\emph{Gōngmín yǒu tóupiào de quánlì.} « Les citoyens ont le droit de vote. »\\[4pt]
\zh{民主国家的人民选举自己的领导人。}\\
\emph{Mínzhǔ guójiā de rénmín xuǎnjǔ zìjǐ de lǐngdǎorén.} « Dans un pays démocratique, le peuple élit ses dirigeants. »\\[4pt]
\zh{自由和责任都很重要。}\\
\emph{Zìyóu hé zérèn dōu hěn zhòngyào.} « La liberté et la responsabilité sont toutes les deux très importantes. »\\[4pt]
\zh{我们的社会需要负责任的公民。}\\
\emph{Wǒmen de shèhuì xūyào fù zérèn de gōngmín.} « Notre société a besoin de citoyens responsables. »
\end{exemplebox}

\begin{attention}[权利 quánlì ou 权力 quánlì ?]
Deux mots se prononcent exactement pareil (\emph{quánlì}) mais ne s'écrivent pas pareil : \zh{权利} = le \textbf{droit} (ce que le citoyen possède : droit de vote, droit à l'éducation) ; \zh{权力} = le \textbf{pouvoir} (ce que l'État exerce). Dans cette leçon, c'est presque toujours \zh{权利} (avec \zh{利}, « bénéfice ») qu'il faut écrire. Piège classique à l'examen !
\end{attention}

\courssub{Les radicaux : la clé pour comprendre les caractères}

Chaque caractère chinois contient un \cle{radical} (ou « clé », \zh{部首} bùshǒu), qui donne souvent un indice de sens. Connaître le radical, c'est déjà à moitié mémoriser le caractère.

\begin{center}
\begin{tabularx}{\linewidth}{|X|X|X|X|}
\hline
\rowcolor{popblueL} Caractère & Radical & Sens du radical & Idée à retenir \\
\hline
民 mín & 氏 (shì, clan) & clan, lignage & le peuple = les familles \\
\hline
主 zhǔ & 丶 (diǎn, point) / 王 & point ; roi & un point au-dessus du roi = le maître \\
\hline
权 quán & 木 (mù, bois) & bois, arbre & partie gauche = bois \\
\hline
票 piào & 示 (shì) & autel, montrer & bulletin « montré » au-dessus \\
\hline
选 xuǎn & 辶 (chuò) & mouvement, marche & choisir = « avancer vers » \\
\hline
责 zé & 贝 (bèi) & coquillage, argent & rendre des comptes \\
\hline
\end{tabularx}
\end{center}

\begin{savaistu}[Le coquillage qui valait de l'argent]
Le radical \zh{贝} (bèi) représente un \textbf{coquillage}. Dans la Chine ancienne, les coquillages servaient de \textbf{monnaie} ! C'est pourquoi on retrouve \zh{贝} dans presque tous les caractères liés à l'argent : \zh{贵} (guì, cher), \zh{买} (mǎi, acheter, forme ancienne), \zh{财} (cái, richesse). Dans \zh{责} (zé, responsabilité), le coquillage rappelle l'idée de « compte à rendre » : celui qui est responsable doit rendre compte de ce qu'on lui a confié — autrefois, de l'argent en coquillages !
\end{savaistu}

\begin{methode}[Mémoriser un caractère en 5 étapes]
\begin{enumerate}
\item \textbf{Identifier le radical} : quelle partie donne le sens ? (ex. \zh{权} → \zh{木}, le bois, à gauche).
\item \textbf{Décomposer} : quelles sont les autres parties ? (ex. \zh{权} = \zh{木} + \zh{又}).
\item \textbf{Compter les traits} : \zh{权} a 6 traits, \zh{民} en a 5, \zh{票} en a 11.
\item \textbf{Écrire 5 fois} en respectant l'ordre des traits, à voix haute.
\item \textbf{Prononcer le pinyin} à chaque écriture : « quán, quán, quán… ». La main, l'œil et l'oreille mémorisent ensemble.
\end{enumerate}
\end{methode}

\courssub{Ordre des traits : 民, 主, 权, 票}

En chinois, on n'écrit pas un caractère « comme on veut ». Il existe des règles générales : \textbf{de haut en bas}, \textbf{de gauche à droite}, \textbf{l'horizontal avant le vertical}, \textbf{l'extérieur avant l'intérieur}. Respecter l'ordre des traits rend l'écriture plus rapide, plus belle, et permet d'utiliser les dictionnaires électroniques.

\textbf{1. Le caractère 民 (mín, peuple) — 5 traits.}\par
On écrit : d'abord le coin en haut à gauche (trait 1 : horizontal + crochet vertical \zh{(trait brisé)}), puis un trait horizontal (trait 2), puis un trait oblique descendant vers la gauche (trait 3 : \zh{撇}), puis un autre horizontal (trait 4), et enfin le grand trait oblique descendant vers la droite (trait 5 : \zh{捺}). Règle appliquée : de haut en bas.

\textbf{2. Le caractère 主 (zhǔ, maître) — 5 traits.}\par
On écrit d'abord \zh{王} (le roi, 4 traits : trois horizontaux et un vertical qui les traverse, de haut en bas), puis on ajoute le \textbf{point} \zh{丶} tout en haut (trait 5). Attention : le point se pose \emph{en dernier}, comme la cerise sur le gâteau.

\textbf{3. Le caractère 权 (quán, droit/pouvoir) — 6 traits.}\par
D'abord la partie gauche \zh{木} (bois, 4 traits : horizontal, vertical, oblique gauche, oblique droite), ensuite la partie droite \zh{又} (2 traits : horizontal crochu, oblique droite). Règle appliquée : de gauche à droite.

\textbf{4. Le caractère 票 (piào, bulletin de vote) — 11 traits.}\par
On écrit de haut en bas : d'abord la partie supérieure \zh{覀} (6 traits : horizontal, vertical, horizontal, oblique gauche, point, horizontal), puis la partie inférieure \zh{示} (5 traits : horizontal, horizontal, vertical, horizontal, point gauche, point droit — soit 5 traits pour 示 en écriture standard : 一 一 丨 一 丶 丶).

\begin{popfigure}
\begin{center}
\begin{tikzpicture}[scale=0.9, font=\small]
% --- 民 : étapes ---
\node[font=\bfseries, popdark] at (1.5,3.4) {民 (5 traits)};
\foreach \i/\txt in {0/1, 1/2, 2/3, 3/4, 4/5} {
  \draw[popline] (\i*1.1,1.6) rectangle (\i*1.1+0.9,2.5);
  \node[popblue] at (\i*1.1+0.45,2.05) {Trait \txt};
}
% --- 权 : étapes ---
\node[font=\bfseries, popdark] at (1.5,0.9) {权 (6 traits)};
\foreach \i in {0,...,3} {
  \draw[popline] (\i*1.1,-0.9) rectangle (\i*1.1+0.9,0);
}
\node[popgreen] at (0.45,-0.45) {木};
\node[popgreen] at (1.55,-0.45) {木 + 丿};
\node[popgreen] at (2.65,-0.45) {木 + 又};
\node[popgreen, font=\large] at (3.75,-0.45) {权};
\node[align=left, popmuted, font=\small] at (5.6,-0.45) {gauche $\rightarrow$ droite};
\end{tikzpicture}
\end{center}
\legende{Construction progressive de 民 (de haut en bas) et de 权 (de gauche à droite) : chaque case ajoute un ou deux traits.}
\end{popfigure}

\begin{attention}[王 ou 主 ? Le point qui change tout]
Erreur très fréquente chez les francophones : écrire \zh{王} (wáng, \textbf{roi}) à la place de \zh{主} (zhǔ, \textbf{maître}). La seule différence est le petit point \zh{丶} en haut ! Si vous oubliez le point, \zh{民主} (mínzhǔ, « la démocratie », litt. « le peuple maître ») devient \zh{民王} : « le peuple roi » — ce qui ne veut rien dire. Retenez : \textbf{le maître porte un chapeau-point}. Même vigilance pour \zh{玉} (yù, jade), qui a un point en bas à droite : trois caractères, trois positions du point.
\end{attention}

\courssub{Caractères proches : ne pas se tromper}

\begin{center}
\begin{tabularx}{\linewidth}{|X|X|X|}
\hline
\rowcolor{popblueL} Paire & Pinyin et sens & Astuce pour distinguer \\
\hline
民 / 氏 & mín (peuple) / shì (nom de famille) & 民 a un trait horizontal de plus en bas \\
\hline
主 / 王 / 玉 & zhǔ (maître) / wáng (roi) / yù (jade) & point en haut = 主 ; pas de point = 王 ; point en bas = 玉 \\
\hline
权 / 劝 / 欢 & quán (droit) / quàn (conseiller) / huān (joyeux) & tous ont 又 à droite ; regardez la gauche : 木, 又, 欠 \\
\hline
票 / 飘 & piào (bulletin) / piāo (flotter au vent) & 飘 a 风 (vent) à droite : le papier qui s'envole \\
\hline
选 / 远 & xuǎn (choisir) / yuǎn (loin) & même radical 辶 ; l'intérieur diffère : 先 contre 元 \\
\hline
\end{tabularx}
\end{center}

\begin{exoresolu}[Identifier le bon caractère]
Énoncé : complétez chaque phrase avec le bon caractère parmi les propositions.\\
1. 民（主／王）很重要。(La démocratie est importante.)\\
2. 公民有投（票／飘）的权利。(Les citoyens ont le droit de vote.)\\
3. 人民（选／远）举领导人。(Le peuple élit ses dirigeants.)
\tcblower
Solution détaillée :\\
1. \zh{民主} : il faut le point en haut (\zh{主}, maître), pas \zh{王} (roi). Phrase : \zh{民主很重要。} \emph{Mínzhǔ hěn zhòngyào.}\\
2. \zh{投票} : le bulletin de vote \zh{票} (radical \zh{示}), pas \zh{飘} (flotter, avec le vent \zh{风}). Phrase : \zh{公民有投票的权利。} \emph{Gōngmín yǒu tóupiào de quánlì.}\\
3. \zh{选举} : choisir \zh{选} (intérieur \zh{先}), pas \zh{远} (loin, intérieur \zh{元}). Phrase : \zh{人民选举领导人。} \emph{Rénmín xuǎnjǔ lǐngdǎorén.}
\end{exoresolu}

\courssub{Exercice de copie guidée}

\begin{exercice}[À la plume (ou au crayon) !]
Sur votre cahier d'écriture, copiez \textbf{trois lignes} de chaque mot ci-dessous, en respectant l'ordre des traits et en prononçant le pinyin à voix haute à chaque caractère :
\begin{itemize}
\item \zh{民主} (mínzhǔ) — 3 lignes. Vérifiez le point de \zh{主} à chaque fois.
\item \zh{权利} (quánlì) — 3 lignes. Le \zh{木} de \zh{权} s'écrit avant le \zh{又}.
\item \zh{投票} (tóupiào) — 3 lignes. \zh{票} s'écrit de haut en bas : \zh{覀} puis \zh{示}.
\end{itemize}
Auto-correction : relisez votre ligne et entourez tout caractère où le point de \zh{主} manque ou où \zh{票} ressemble à \zh{飘}.
\end{exercice}

\begin{corrige}[Exercice de copie]
Il n'y a pas de réponse unique : le corrigé est votre propre cahier. Critères de réussite : (1) chaque caractère tient dans un carré imaginaire équilibré ; (2) \zh{主} a bien son point au-dessus ; (3) \zh{权} = \zh{木} à gauche + \zh{又} à droite, sans fusion ; (4) \zh{票} a bien la partie \zh{示} en bas avec ses deux petits points. Faites-vous relire par un camarade : échanger les cahiers est le meilleur correcteur.
\end{corrige}

\begin{aretenir}[L'essentiel de la section]
\begin{itemize}
\item Dix mots-clés du thème : \zh{民主} mínzhǔ, \zh{人民} rénmín, \zh{公民} gōngmín, \zh{权利} quánlì, \zh{自由} zìyóu, \zh{投票} tóupiào, \zh{选举} xuǎnjǔ, \zh{责任} zérèn, \zh{国家} guójiā, \zh{社会} shèhuì.
\item Le radical donne le sens : \zh{贝} (argent) dans \zh{责}, \zh{辶} (mouvement) dans \zh{选}, \zh{木} (bois) dans \zh{权}.
\item Ordre des traits : haut → bas, gauche → droite ; le point de \zh{主} s'écrit en dernier.
\item Pièges : \zh{主}/\zh{王}/\zh{玉}, \zh{权利}/\zh{权力}, \zh{票}/\zh{飘}, \zh{选}/\zh{远}.
\item Méthode : radical → décomposition → comptage des traits → 5 écritures en prononçant.
\end{itemize}
\end{aretenir}

Maintenant que le lexique est acquis et que votre main sait tracer ces caractères, nous pouvons passer à la section suivante : la structure du texte argumentatif et les connecteurs qui relieront vos idées.

\coursec{Structure du texte et connecteurs}

Dans la section précédente, nous avons appris à écrire correctement les caractères du thème de la démocratie : \zh{民主} (mínzhǔ), \zh{权利} (quánlì), \zh{选举} (xuǎnjǔ), \zh{自由} (zìyóu), \zh{责任} (zérèn). Mais posséder des mots ne suffit pas : pour écrire un texte argumentatif simple, il faut savoir \cle{organiser} ses idées et les \cle{relier} entre elles. C'est exactement le rôle des connecteurs, les « petites charnières » du discours. Cette section vous donne le plan type du texte argumentatif et la boîte à outils complète des connecteurs chinois.

\courssub{Observer d'abord : un texte argumentatif modèle}

Lisons d'abord ce court texte d'un élève camerounais. Repérez les mots soulignés par votre œil : ce sont les connecteurs.

\begin{textebox}[民主很重要吗？ — La démocratie est-elle importante ?]
我认为民主很重要。\zh{首先}，因为民主保护公民的权利，所以每个人可以自由地说话。\\
\zh{其次}，民主让人民选择自己的领导人。在喀麦隆，我们应该参加选举，因为选举是公民的责任。\\
\zh{但是}，民主不只是权利，\zh{而且}也需要责任：我们应该尊重别人的意见，也应该遵守法律。\\
\zh{最后}，我认为民主能让社会更和平、更进步。\zh{所以}，我们年轻人应该学习民主，也应该参加国家的生活。\\[0.5em]
\emph{Wǒ rènwéi mínzhǔ hěn zhòngyào. Shǒuxiān, yīnwèi mínzhǔ bǎohù gōngmín de quánlì, suǒyǐ měi gè rén kěyǐ zìyóu de shuōhuà. Qícì, mínzhǔ ràng rénmín xuǎnzé zìjǐ de lǐngdǎorén. Zài Kāmàilóng, wǒmen yīnggāi cānjiā xuǎnjǔ, yīnwèi xuǎnjǔ shì gōngmín de zérèn. Dànshì, mínzhǔ bù zhǐ shì quánlì, érqiě yě xūyào zérèn： wǒmen yīnggāi zūnzhòng biérén de yìjiàn, yě yīnggāi zūnshǒu fǎlǜ. Zuìhòu, wǒ rènwéi mínzhǔ néng ràng shèhuì gèng hépíng, gèng jìnbù. Suǒyǐ, wǒmen niánqīngrén yīnggāi xuéxí mínzhǔ, yě yīnggāi cānjiā guójiā de shēnghuó.}\\[0.5em]
« Je pense que la démocratie est très importante. D'abord, parce que la démocratie protège les droits des citoyens, chacun peut parler librement. Ensuite, la démocratie permet au peuple de choisir ses dirigeants. Au Cameroun, nous devons participer aux élections, car voter est un devoir de citoyen. Mais la démocratie n'est pas seulement des droits, elle exige aussi des responsabilités : nous devons respecter l'opinion des autres et respecter la loi. Enfin, je pense que la démocratie rend la société plus pacifique et plus progressiste. C'est pourquoi nous, les jeunes, devons apprendre la démocratie et participer à la vie du pays. »
\end{textebox}

\textbf{Questions de compréhension}
\begin{enumerate}
\item Combien d'arguments l'auteur donne-t-il pour sa thèse ? Lesquels ?
\item Où se trouve la nuance (la limite) dans ce texte ?
\item Quel connecteur introduit la conclusion finale ?
\end{enumerate}

\textbf{Glossaire du texte}

\begin{center}\begin{tabularx}{\linewidth}{|X|X|X|X|}\hline
\rowcolor{popblueL} Caractères & Pinyin & Nature & Traduction \\ \hline
认为 & rènwéi & verbe & penser, estimer \\ \hline
保护 & bǎohù & verbe & protéger \\ \hline
公民 & gōngmín & nom & citoyen \\ \hline
领导人 & lǐngdǎorén & nom & dirigeant \\ \hline
尊重 & zūnzhòng & verbe & respecter (une personne, une opinion) \\ \hline
遵守 & zūnshǒu & verbe & respecter (une règle, une loi) \\ \hline
法律 & fǎlǜ & nom & loi, législation \\ \hline
和平 & hépíng & nom / adj. & paix ; pacifique \\ \hline
进步 & jìnbù & nom / adj. & progrès ; progressiste \\ \hline
年轻人 & niánqīngrén & nom & jeune(s), la jeunesse \\ \hline
国家 & guójiā & nom & pays, État, nation \\ \hline
生活 & shēnghuó & nom / verbe & vie, vivre \\ \hline
参加 & cānjiā & verbe & participer à \\ \hline
选择 & xuǎnzé & verbe & choisir \\ \hline
自己 & zìjǐ & pronom & soi-même, propre \\ \hline
因为 & yīnwèi & conjonction & parce que \\ \hline
所以 & suǒyǐ & conjonction & donc, c'est pourquoi \\ \hline
但是 & dànshì & conjonction & mais \\ \hline
而且 & érqiě & conjonction & de plus, et en outre \\ \hline
最后 & zuìhòu & adv. / conj. & enfin, finalement \\ \hline
也 & yě & adv. & aussi \\ \hline
应该 & yīnggāi & auxiliaire & devoir (obligation morale) \\ \hline
可以 & kěyǐ & auxiliaire & pouvoir (permission) \\ \hline
需要 & xūyào & verbe / aux. & avoir besoin de, exiger \\ \hline
让 & ràng & verbe & permettre, faire (causatif) \\ \hline
\end{tabularx}\end{center}

\courssub{Le plan type du texte argumentatif simple}

En Terminale, on vous demande de produire un texte simple et clair, pas un essai de philosophe. Le plan à retenir comporte \cle{cinq blocs} logiques :

\begin{definition}[Plan type en cinq blocs]
\begin{enumerate}
\item \textbf{La thèse} : une phrase qui annonce votre opinion. Formule utile : \zh{我认为……} (Wǒ rènwéi... = je pense que...).
\item \textbf{L'argument 1} : introduit par \zh{首先} (d'abord), développé avec une cause (\zh{因为}…) et une conséquence (\zh{所以}…).
\item \textbf{L'argument 2} : introduit par \zh{其次} (ensuite) ou \zh{然后} (puis), éventuellement renforcé par \zh{而且} (de plus).
\item \textbf{La nuance (concession)} : une phrase avec \zh{但是} ou \zh{可是} qui montre que vous voyez les limites ou l'autre face de la médaille.
\item \textbf{La conclusion} : une phrase avec \zh{所以} (c'est pourquoi) ou \zh{最后} (enfin) qui résume votre position et ouvre sur l'action.
\end{enumerate}
\end{definition}

\begin{popfigure}
\begin{tikzpicture}[mindmap, grow cyclic, every node/.style={concept}, concept color=popblueL, text=popdark,
level 1/.append style={level distance=3.2cm, sibling angle=90},
level 2/.append style={level distance=2.4cm, sibling angle=45}]
\node[concept, concept color=popblue, text=white, align=center]{民主\\ mínzhǔ}
  child[concept color=popgreenL] { node[concept, align=center]{Thèse\\ 我认为…} }
  child[concept color=poporangeL] { node[concept, align=center]{Argument 1\\ 首先\\ 因为…所以…} }
  child[concept color=poporangeL] { node[concept, align=center]{Argument 2\\ 其次\\ 而且…} }
  child[concept color=poppinkL] { node[concept, align=center]{Nuance\\ 但是\\ 可是} }
  child[concept color=popgoldL] { node[concept, align=center]{Conclusion\\ 所以\\ 最后} };
\end{tikzpicture}
\legende{Carte mentale : la structure d'un texte argumentatif simple sur la démocratie (5 blocs).}
\end{popfigure}

\courssub{Les connecteurs d'énumération : ordonner ses arguments}

Pour présenter plusieurs arguments, le chinois utilise une série de connecteurs d'ordre, comme « d'abord, ensuite, puis, enfin » en français.

\begin{center}\begin{tabularx}{\linewidth}{|X|X|X|}\hline
\rowcolor{popblueL} Connecteur & Pinyin & Traduction \\ \hline
首先 & shǒuxiān & d'abord, premièrement \\ \hline
其次 & qícì & ensuite, deuxièmement \\ \hline
然后 & ránhòu & puis, après (souvent pour actions successives) \\ \hline
最后 & zuìhòu & enfin, finalement \\ \hline
\end{tabularx}\end{center}

\begin{exemplebox}[Énumération en action]
\zh{首先，民主保护我们的权利；其次，民主让人民选择领导人；最后，民主让社会更进步。}\\
\emph{Shǒuxiān, mínzhǔ bǎohù wǒmen de quánlì； qícì, mínzhǔ ràng rénmín xuǎnzé lǐngdǎorén； zuìhòu, mínzhǔ ràng shèhuì gèng jìnbù.}\\
« D'abord, la démocratie protège nos droits ; ensuite, elle permet au peuple de choisir ses dirigeants ; enfin, elle fait progresser la société. »
\end{exemplebox}

\begin{attention}[首先 / 其次 / 最后 vs 然后 : deux registres]
\zh{首先}、\zh{其次}、\zh{最后} servent à énumérer des \textbf{arguments logiques} (registre écrit, argumentation, dissertation).\\
\zh{然后} sert surtout à énumérer des \textbf{actions chronologiques} dans le récit : \zh{我先吃饭，然后去学校。} (Wǒ xiān chīfàn, ránhòu qù xuéxiào. = Je mange d'abord, puis je vais à l'école.) Dans un texte argumentatif, préférez la série \zh{首先}…\zh{其次}…\zh{最后}.
\end{attention}

\courssub{Cause et conséquence : 因为……所以……}

Observons deux phrases :

\begin{exemplebox}[Observation]
\zh{因为民主保护公民的权利，所以民主很重要。}\\
\emph{Yīnwèi mínzhǔ bǎohù gōngmín de quánlì, suǒyǐ mínzhǔ hěn zhòngyào.}\\
« Parce que la démocratie protège les droits des citoyens, elle est très importante. »\\[0.4em]
\zh{民主保护公民的权利，所以它很重要。}\\
\emph{Mínzhǔ bǎohù gōngmín de quánlì, suǒyǐ tā hěn zhòngyào.}\\
« La démocratie protège les droits des citoyens, donc elle est très importante. »
\end{exemplebox}

\begin{propriete}[La règle de 因为……所以……]
Structure : \textbf{因为 + cause，所以 + conséquence。}
\begin{itemize}
\item Contrairement au français, le chinois \textbf{utilise couramment les deux connecteurs dans la même phrase} : \zh{因为}…\zh{所以}… (litt. « parce que... donc... »).
\item On peut aussi n'utiliser que \zh{因为} seul (la conséquence est implicite), ou \zh{所以} seul (la cause est connue du contexte).
\item \zh{所以} seul est le connecteur idéal pour ouvrir la \textbf{conclusion} d'un texte : \zh{所以，我们应该学习民主。} (Suǒyǐ, wǒmen yīnggāi xuéxí mínzhǔ. = C'est pourquoi nous devons apprendre la démocratie.)
\end{itemize}
\end{propriete}

\begin{attention}[Piège majeur en version (chinois → français)]
En français, « parce que... donc... » est une \textbf{faute de syntaxe} (pléonasme). Quand vous traduisez une phrase chinoise contenant \zh{因为}…\zh{所以}…, vous devez \textbf{choisir un seul} connecteur en français : soit « parce que / comme / puisque » (cause), soit « donc / c'est pourquoi / par conséquent » (conséquence).\\
\zh{因为选举很重要，所以我们都应该参加。} → « \textbf{Parce que} les élections sont importantes, nous devons tous y participer » \textbf{OU} « Les élections sont importantes, \textbf{c'est pourquoi} nous devons tous y participer. » (Jamais « Parce que... donc... »).
\end{attention}

\courssub{Opposition et addition : nuancer et enrichir}

\textbf{L'opposition} se dit \zh{但是} (dànshì, mais) ou \zh{可是} (kěshì, cependant, un peu plus oral). Le connecteur se place en tête de la deuxième proposition (soujet + 但是 + verbe...).

\begin{exemplebox}[La nuance avec 但是]
\zh{民主给我们自由，但是我们也需要责任。}\\
\emph{Mínzhǔ gěi wǒmen zìyóu, dànshì wǒmen yě xūyào zérèn.}\\
« La démocratie nous donne la liberté, mais nous avons aussi besoin de responsabilité. »
\end{exemplebox}

\textbf{L'addition} se dit \zh{也} (yě, aussi) — adverbe placé \textbf{avant le verbe} — ou \zh{而且} (érqiě, de plus, et en outre) qui relie deux propositions complètes.

\begin{exemplebox}[Addition : 也 vs 而且]
\zh{民主也需要责任。}\\
\emph{Mínzhǔ yě xūyào zérèn.} « La démocratie exige \textbf{aussi} la responsabilité. » (焦点 sur le verbe \zh{需要})\\[0.4em]
\zh{我们应该参加选举，而且应该尊重别人的意见。}\\
\emph{Wǒmen yīnggāi cānjiā xuǎnjǔ, érqiě yīnggāi zūnzhòng biérén de yìjiàn.}\\
« Nous devons participer aux élections \textbf{et de plus} respecter l'opinion des autres. » (liaison de deux propositions)
\end{exemplebox}

\begin{attention}[Place de 也 : avant le verbe !]
\zh{也} se place \textbf{toujours immédiatement avant le verbe} (ou l'auxiliaire), jamais en fin de phrase comme « aussi » en français.\\
Correct : \zh{我也喜欢民主。} (Wǒ \textbf{yě} xǐhuan mínzhǔ. = Moi aussi, j'aime la démocratie.)\\
Incorrect : \zh{我喜欢民主也。}
\end{attention}

\courssub{Les verbes auxiliaires indispensables pour argumenter}

Pour exprimer le devoir, la permission, le besoin ou la causalité, quatre auxiliaires sont incontournables. Ils se placent \textbf{avant le verbe principal} (Sujet + Auxiliaire + Verbe + Complément).

\begin{center}\begin{tabularx}{\linewidth}{|X|X|X|X|}\hline
\rowcolor{popblueL} Auxiliaire & Pinyin & Sens & Exemple (Caractères + Pinyin + Traduction) \\ \hline
应该 & yīnggāi & devoir (moral, conseil) & \zh{我们应该参加选举。} (Wǒmen yīnggāi cānjiā xuǎnjǔ.) Nous devons participer aux élections. \\ \hline
可以 & kěyǐ & pouvoir (permission, possibilité) & \zh{公民可以自由地说话。} (Gōngmín kěyǐ zìyóu de shuōhuà.) Les citoyens peuvent parler librement. \\ \hline
需要 & xūyào & avoir besoin de, exiger & \zh{民主需要责任。} (Mínzhǔ xūyào zérèn.) La démocratie exige de la responsabilité. \\ \hline
让 & ràng & permettre, faire (causatif) & \zh{民主让社会更和平。} (Mínzhǔ ràng shèhuì gèng hépíng.) La démocratie rend la société plus pacifique. \\ \hline
\end{tabularx}\end{center}

\begin{exemplebox}[Argument complet avec auxiliaire]
\zh{公民应该参加选举，因为选举让人民选择领导人。}\\
\emph{Gōngmín yīnggāi cānjiā xuǎnjǔ, yīnwèi xuǎnjǔ ràng rénmín xuǎnzé lǐngdǎorén.}\\
« Les citoyens doivent participer aux élections, car les élections permettent au peuple de choisir ses dirigeants. »
\end{exemplebox}

\begin{attention}[应该 n'est pas un verbe conjugué]
En chinois, l'auxiliaire \zh{应该} (comme \zh{可以}, \zh{需要}, \zh{让}) \textbf{ne change jamais de forme} (pas de conjugaison, pas d'accord) et le verbe qui suit reste à la forme de base (pas de suffixe).\\
\zh{他应该学习} (Tā yīnggāi xuéxí = Il doit étudier) — pas de \zh{了}, pas de \zh{着}, pas de marque de personne.
\end{attention}

\courssub{Tableau récapitulatif des connecteurs et marqueurs d'argumentation}

\begin{center}\begin{tabularx}{\linewidth}{|X|X|X|X|}\hline
\rowcolor{popblueL} Fonction & Chinois & Pinyin & Français \\ \hline
Thèse (opinion) & 我认为… & wǒ rènwéi & je pense que… / j'estime que… \\ \hline
Énumération (arguments) & 首先 / 其次 / 最后 & shǒuxiān / qícì / zuìhòu & d'abord / ensuite / enfin \\ \hline
Cause & 因为 & yīnwèi & parce que / comme \\ \hline
Conséquence / Conclusion & 所以 & suǒyǐ & donc / c'est pourquoi / par conséquent \\ \hline
Opposition / Nuance & 但是 / 可是 & dànshì / kěshì & mais / cependant \\ \hline
Addition (même proposition) & 也 & yě & aussi (devant verbe) \\ \hline
Addition (deux propositions) & 而且 & érqiě & de plus / et en outre \\ \hline
\end{tabularx}\end{center}

\courssub{Méthode : construire son texte étape par étape}

\begin{methode}[Du plan français au texte chinois]
\begin{enumerate}
\item \textbf{Rédiger le plan en français} : 1 phrase de thèse, 2 arguments distincts, 1 nuance, 1 conclusion. Maximum 5 phrases.
\item \textbf{Insérer les connecteurs logiques français} : \emph{D'abord}, \emph{Ensuite}, \emph{Mais}, \emph{Donc / C'est pourquoi}.
\item \textbf{Traduire phrase par phrase} en chinois, en substituant les connecteurs : \emph{D'abord} → \zh{首先}, \emph{Ensuite} → \zh{其次}, \emph{Mais} → \zh{但是}, \emph{Donc} → \zh{所以}.
\item \textbf{Vérifier l'ordre des mots (SOV strict)} : Sujet + (Temps) + (Lieu) + Auxiliaire (\zh{应该}/\zh{可以}…) + \zh{也} + Verbe + Complément. Cause (\zh{因为}…) avant Conséquence (\zh{所以}…).
\item \textbf{Relire et valider} : Chaque phrase argumentative contient-elle un connecteur ou un auxiliaire modal (\zh{应该}, \zh{可以}, \zh{需要}, \zh{让}) ? La ponctuation chinoise (， 。 ；) est-elle correcte ?
\end{enumerate}
\end{methode}

\begin{exoresolu}[Relier les phrases avec le bon connecteur]
Énoncé : Complétez avec \zh{首先}, \zh{但是}, \zh{所以} ou \zh{而且}.\\
1. \zh{……，我认为民主很重要。} \\
2. \zh{民主给我们自由，……我们也需要责任。} \\
3. \zh{民主保护权利，……它让社会进步。} \\
4. \zh{……，年轻人应该学习民主。}
\tcblower
Solution détaillée :\\
1. \zh{首先，我认为民主很重要。} (Shǒuxiān, wǒ rènwéi mínzhǔ hěn zhòngyào.) — On ouvre l'argumentation : « D'abord, je pense que la démocratie est importante. »\\
2. \zh{民主给我们自由，但是我们也需要责任。} — Opposition entre liberté (droit) et responsabilité (devoir) : « mais ».\\
3. \zh{民主保护权利，而且它让社会进步。} — On ajoute un second effet positif (addition) : « de plus ». \zh{但是} serait illogique (pas d'opposition).\\
4. \zh{所以，年轻人应该学习民主。} — Conclusion logique tirée des arguments précédents : « c'est pourquoi ».
\end{exoresolu}

\begin{exercice}[À vous de jouer]
\begin{enumerate}
\item Traduisez en chinois : « D'abord, la démocratie protège nos droits ; ensuite, elle nous permet de choisir nos dirigeants ; c'est pourquoi elle est importante. »
\item Transformez en une seule phrase avec \zh{因为}…\zh{所以}… : \zh{选举很重要。我们都应该参加。}
\item Ajoutez une nuance avec \zh{但是} à la phrase : \zh{民主给我们自由。}
\end{enumerate}
\end{exercice}

\begin{corrige}[Exercice « À vous de jouer »]
\begin{enumerate}
\item \zh{首先，民主保护我们的权利；其次，它让我们选择自己的领导人；所以，民主很重要。} (Shǒuxiān, mínzhǔ bǎohù wǒmen de quánlì； qícì, tā ràng wǒmen xuǎnzé zìjǐ de lǐngdǎorén； suǒyǐ, mínzhǔ hěn zhòngyào.)
\item \zh{因为选举很重要，所以我们都应该参加。} (Yīnwèi xuǎnjǔ hěn zhòngyào, suǒyǐ wǒmen dōu yīnggāi cānjiā.)
\item \zh{民主给我们自由，但是我们也应该尊重别人。} (Mínzhǔ gěi wǒmen zìyóu, dànshì wǒmen yě yīnggāi zūnzhòng biérén.) — Toute nuance logique (responsabilité, loi, respect d'autrui) est acceptée.
\end{enumerate}
\end{corrige}

\begin{aretenir}[L'essentiel de la section]
Un bon texte argumentatif chinois tient en cinq blocs : thèse (\zh{我认为}…) + argument 1 (\zh{首先} + \zh{因为}…\zh{所以}…) + argument 2 (\zh{其次} / \zh{而且}) + nuance (\zh{但是}) + conclusion (\zh{所以} / \zh{最后}). Différence clé avec le français : \zh{因为} et \zh{所以} coexistent dans la phrase chinoise, mais en français on choisit \textbf{soit} la cause \textbf{soit} la conséquence. Ordre immuable : \zh{也} + verbe ; Auxiliaire (\zh{应该}, \zh{可以}, \zh{需要}, \zh{让}) + verbe. Ponctuation chinoise pleine chasse (， 。 ； ：) obligatoire.
\end{aretenir}

\coursec{Thème : français → chinois}

Dans la section précédente, nous avons vu comment structurer un texte sur la démocratie et quels connecteurs employer (\zh{因为……所以……}, \zh{首先}, \zh{但是}…). Nous allons maintenant utiliser ces acquis dans l'exercice roi de l'expression écrite au programme de Terminale : le \cle{thème}, c'est-à-dire la traduction du français vers le chinois. Traduire vers le chinois n'est pas un simple remplacement de mots : c'est un travail en quatre étapes — comprendre, découper, réassembler, relire. C'est cette méthode que nous allons construire ensemble, phrase par phrase.

\begin{center}
\begin{tabularx}{\linewidth}{|X|X|X|X|}
\hline
\rowcolor{popblueL} Caractères & Pinyin & Nature & Traduction \\
\hline
民主 & mínzhǔ & n. & démocratie \\
\hline
公民 & gōngmín & n. & citoyen \\
\hline
权利 & quánlì & n. & droit, privilège \\
\hline
选举 & xuǎnjǔ & n./v. & élection / élire \\
\hline
法律 & fǎlǜ & n. & loi \\
\hline
领导人 & lǐngdǎorén & n. & dirigeant, leader \\
\hline
保护 & bǎohù & v. & protéger \\
\hline
自由 & zìyóu & n./adj. & liberté / libre \\
\hline
投票 & tóupiào & v./n. & voter / vote \\
\hline
选择 & xuǎnzé & v. & choisir \\
\hline
应该 & yīnggāi & modal & devoir, devrait \\
\hline
因为 & yīnwèi & conj. & parce que \\
\hline
所以 & suǒyǐ & conj. & donc, par conséquent \\
\hline
很 & hěn & adv. & très (marque de l'adjectif prédicat) \\
\hline
非常 & fēicháng & adv. & très, extrêmement \\
\hline
也 & yě & adv. & aussi \\
\hline
都 & dōu & adv. & tous, toutes \\
\hline
自己 & zìjǐ & pron. & soi-même, propre \\
\hline
的 & de & part. & particule de détermination \\
\hline
重要 & zhòngyào & adj. & important \\
\hline
活力 & huólì & n. & vitalité, dynamisme \\
\hline
尊重 & zūnzhòng & v. & respecter \\
\hline
爱 & ài & v. & aimer \\
\hline
人民 & rénmín & n. & le peuple \\
\hline
参加 & cānjiā & v. & participer \\
\hline
\end{tabularx}
\end{center}

\begin{definition}[Écriture et analyse des caractères clés du thème]
Voici les points de vigilance pour l'écriture correcte des caractères fréquents de cette leçon (ordre des traits : haut → bas, gauche → droite, horizontal avant vertical, trait central avant les ailes) :
\begin{itemize}
\item \textbf{民 (mín)} : 5 traits. Clé \zh{氏} (shì). \textbf{Piège :} ne pas confondre avec \zh{氏} (shì, nom de famille, clé \zh{氏} aussi mais structure différente : \zh{民} a un trait horizontal de plus en bas).
\item \textbf{主 (zhǔ)} : 5 traits. Clé \zh{丶} (zhǔ). Le point initial est distinct du trait horizontal de \zh{王} (wáng).
\item \textbf{权 (quán)} : 6 traits (simplifié). Clé \zh{木} (mù). \textbf{Piège :} ne pas confondre avec \zh{劝} (quàn, persuader, clé \zh{力} lì) ni avec \zh{卷} (juǎn, rouleau).
\item \textbf{选 (xuǎn)} : 9 traits. Clé \zh{辶} (chuò, marche). Écrire le corps \zh{巽} (xùn) puis l'enveloppant \zh{辶}.
\item \textbf{领 (lǐng)} : 11 traits. Clé \zh{页} (yè, tête/page). Le haut est \zh{令} (lìng).
\item \textbf{导 (dǎo)} : 15 traits. Clé \zh{寸} (cùn). Structure complexe : \zh{巳} + \zh{寸} en bas à droite.
\item \textbf{护 (hù)} : 7 traits. Clé \zh{扌} (shǒu, main). Droite : \zh{戌} (xū). Ne pas ajouter de trait horizontal en bas (différent de \zh{護} traditionnel).
\item \textbf{法 (fǎ)} : 8 traits. Clé \zh{氵} (shuǐ). Droite : \zh{去} (qù).
\item \textbf{律 (lǜ)} : 9 traits. Clé \zh{彳} (chì, double pas). Droite : \zh{聿} (yù).
\end{itemize}
\end{definition}

\courssub{Observer avant de traduire : quatre phrases modèles}

Lisons d'abord les quatre phrases du jour, déjà traduites. Observez ce qui change entre le français et le chinois.

\begin{exemplebox}[Les quatre phrases du thème]
\textbf{Phrase 1.} « La démocratie est importante. »\\
\zh{民主很重要。} \emph{Mínzhǔ hěn zhòngyào.}\\
\textbf{Phrase 2.} « Les citoyens ont le droit de vote. »\\
\zh{公民有投票的权利。} \emph{Gōngmín yǒu tóupiào de quánlì.}\\
\textbf{Phrase 3.} « Nous devons choisir nos dirigeants. »\\
\zh{我们应该选择自己的领导人。} \emph{Wǒmen yīnggāi xuǎnzé zìjǐ de lǐngdǎorén.}\\
\textbf{Phrase 4.} « Parce que la démocratie protège la liberté, elle est importante. »\\
\zh{因为民主保护自由，所以民主很重要。} \emph{Yīnwèi mínzhǔ bǎohù zìyóu, suǒyǐ mínzhǔ hěn zhòngyào.}
\end{exemplebox}

Trois observations s'imposent :

\begin{itemize}
\item Dans la phrase 1, le verbe français « est » a \textbf{disparu} : l'adjectif \zh{重要} fait office de verbe à lui tout seul (adjectif prédicat).
\item Dans la phrase 3, « devons » (\zh{应该}, verbe modal) se place \textbf{avant} le verbe principal \zh{选择}, exactement comme en français.
\item Dans la phrase 4, le chinois garde \textbf{les deux} connecteurs \zh{因为} et \zh{所以}, alors que le français n'emploie souvent que « parce que » (le « donc » étant implicite ou omis).
\end{itemize}

\courssub{La méthode du thème en quatre étapes}

\begin{methode}[Traduire une phrase du français vers le chinois]
\begin{enumerate}
\item \textbf{Souligner} dans la phrase française les mots qui figurent dans le glossaire ou le tableau de vocabulaire de la leçon (\zh{民主}, \zh{公民}, \zh{权利}, \zh{选择}…). Ce sont eux qui doivent apparaître dans la traduction.
\item \textbf{Traduire mot à mot} chaque élément, sans se soucier encore de l'ordre : poser les « briques » chinoises sur la table.
\item \textbf{Réordonner} selon le schéma chinois canonique : \textbf{Sujet + (Adverbe) + (Modal) + Verbe + Complément}.\\
Vérifier spécifiquement :
\begin{itemize}
\item Place des adverbes (\zh{也}, \zh{都}, \zh{很}, \zh{不}) : \textbf{avant} le verbe (ou le modal s'il y en a un).
\item Place des compléments de temps et de lieu : \textbf{avant} le verbe.
\item Particule \zh{的} : obligatoire entre un déterminant (nom, pronom, verbe, phrase) et le nom qu'il détermine (\zh{投票的权利}, \zh{自己的领导人}).
\end{itemize}
\item \textbf{Relire} : tons corrects sur chaque syllabe (attention au \emph{sandhi} du 3e ton : \zh{很} hěn + \zh{重要} zhòngyào $\rightarrow$ hén zhòngyào à l'oral, mais on écrit hěn), caractères exacts (attention aux paires proches : \zh{民}/\zh{氏}, \zh{权}/\zh{劝}, \zh{选}/\zh{癣}, \zh{领}/\zh{令}), ponctuation chinoise (， 。).
\end{enumerate}
\end{methode}

\begin{popfigure}
\begin{tikzpicture}[node distance=0.55cm, every node/.style={font=\small, align=center}]
\node[draw=popblue, fill=popblueL, rounded corners, text width=4.0cm] (fr) {Phrase française\\ « Nous devons choisir nos dirigeants. »};
\node[draw=poporange, fill=poporangeL, rounded corners, text width=4.4cm, right=1.0cm of fr] (dec) {Découpage\\ S : nous \zh{我们}\\ Modal : devons \zh{应该}\\ V : choisir \zh{选择}\\ C : nos dirigeants \zh{自己的领导人}};
\node[draw=poppurple, fill=poppurpleL, rounded corners, text width=4.2cm, right=1.0cm of dec] (briques) {Briques chinoises\\ 我们 / 应该 / 选择 / 自己的领导人};
\node[draw=popgreen, fill=popgreenL, rounded corners, text width=6.0cm, below=1.0cm of briques] (zh) {Réassemblage S + Modal + V + C\\ \zh{我们应该选择自己的领导人。}};
\draw[-{Stealth[length=3mm]}, thick, popdark] (fr) -- (dec);
\draw[-{Stealth[length=3mm]}, thick, popdark] (dec) -- (briques);
\draw[-{Stealth[length=3mm]}, thick, popdark] (briques) -- (zh);
\end{tikzpicture}
\legende{Schéma de transformation : de la phrase française à la phrase chinoise en quatre temps.}
\end{popfigure}

\courssub{Le piège n° 1 : « est + adjectif » ne se traduit jamais par 是}

Comparons les deux langues. En français, l'adjectif attribut exige le verbe « être » : « la démocratie \emph{est} importante ». Le francophone traduit donc spontanément « être » par \zh{是}. C'est l'erreur la plus fréquente — et la plus grave — du thème.

\begin{propriete}[L'adjectif prédicat : l'adjectif chinois fait office de verbe]
En chinois, un adjectif (verbe d'état) peut fonctionner seul comme prédicat, \textbf{sans} \zh{是}. On ajoute presque toujours un adverbe de degré devant, le plus souvent \zh{很} (qui perd alors son sens lexical de « très » pour ne garder qu'une fonction prosodique d'équilibre de la phrase).\\
Structure : \textbf{Sujet + 很 / 非常 / 真 / 不 + Adjectif}\\
\zh{自由很重要。} \emph{Zìyóu hěn zhòngyào.} « La liberté est importante. »\\
\zh{民主非常好。} \emph{Mínzhǔ fēicháng hǎo.} « La démocratie est excellente. »\\
\zh{这个问题不简单。} \emph{Zhège wèntí bù jiǎndān.} « Ce problème n'est pas simple. »\\
Exception : la construction \zh{是……的} existe, mais elle sert à \textbf{insister} sur un fait accompli ou à confirmer (« c'est bien… », « il est vrai que… »), jamais à traduire le simple verbe « être » attributif.
\end{propriete}

\begin{attention}[Ne traduisez pas « est » par 是 devant un adjectif]
Phrase fautive typique du francophone : \zh{民主是重要。} (incorrect) ou \zh{民主是重要的。} (signifie « la démocratie, \emph{elle}, est vraiment importante » — insistance contrastive, souvent en réponse à une objection).\\
\textbf{Règle d'or :} « être + adjectif » $\rightarrow$ \zh{很 + adjectif} (ou \zh{非常}, \zh{真}, \zh{不}).\\
En revanche, « être + nom » garde bien \zh{是} : \zh{他是公民。} \emph{Tā shì gōngmín.} « Il est citoyen. »
\end{attention}

\courssub{Analyse guidée des quatre phrases}

\textbf{Phrase 1.} « La démocratie est importante. »\\
Découpage : S = la démocratie (\zh{民主}) ; prédicat = est importante (adjectif \zh{重要}). Application de la règle : pas de \zh{是}, on place \zh{很} devant l'adjectif. Résultat : \zh{民主很重要。} \emph{Mínzhǔ hěn zhòngyào.}

\textbf{Phrase 2.} « Les citoyens ont le droit de vote. »\\
Découpage : S = les citoyens \zh{公民} ; V = avoir \zh{有} ; C = le droit de vote. « Le droit de vote » est un nom (\zh{权利}) complété par un verbe (\zh{投票}) : le chinois exige la particule \zh{的} entre les deux : \zh{投票的权利} \emph{tóupiào de quánlì}, littéralement « le droit de voter ». Résultat : \zh{公民有投票的权利。} — ordre SVO identique au français, seule la particule \zh{的} est nouvelle. \textit{Note : le mot composé \zh{选举权} (xuǎnjǔquán, droit de vote/suffrage) est aussi très usité.}

\textbf{Phrase 3.} « Nous devons choisir nos dirigeants. »\\
Découpage : S = nous \zh{我们} ; modal = devoir \zh{应该} ; V = choisir \zh{选择} ; C = nos dirigeants. Deux points de vigilance : le verbe modal se place \textbf{avant} le verbe principal (\zh{应该选择}), et « nos » devant un nom de personne se rend élégamment par \zh{自己的} (« ses propres ») : \zh{自己的领导人}. Résultat : \zh{我们应该选择自己的领导人。}

\textbf{Phrase 4.} « Parce que la démocratie protège la liberté, elle est importante. »\\
Le français n'emploie qu'un seul connecteur explicite ; le chinois préfère la paire corrélative \zh{因为……所以……}. Attention au pronom « elle » : le chinois répète volontiers le nom \zh{民主} plutôt que d'utiliser le pronom \zh{它} dans ce type de phrase formelle. Et dans la seconde proposition, « est importante » redevient \zh{很重要} (règle de l'adjectif prédicat). Résultat : \zh{因为民主保护自由，所以民主很重要。}

\begin{center}
\begin{tabularx}{\linewidth}{|X|X|X|}
\hline
\rowcolor{popblueL} Structure française & Exemple chinois + pinyin & Traduction \\
\hline
être + adjectif & 民主很重要。Mínzhǔ hěn zhòngyào. & La démocratie est importante. \\
\hline
être + nom & 他是公民。Tā shì gōngmín. & Il est citoyen. \\
\hline
avoir + V + 的 + nom & 公民有投票的权利。Gōngmín yǒu tóupiào de quánlì. & Les citoyens ont le droit de vote. \\
\hline
modal + verbe & 我们应该选择领导人。Wǒmen yīnggāi xuǎnzé lǐngdǎorén. & Nous devons choisir les dirigeants. \\
\hline
parce que… (donc)… & 因为民主保护自由，所以民主很重要。Yīnwèi mínzhǔ bǎohù zìyóu, suǒyǐ mínzhǔ hěn zhòngyào. & Parce que la démocratie protège la liberté, elle est importante. \\
\hline
aussi / tous + verbe & 人民也都参加选举。Rénmín yě dōu cānjiā xuǎnjǔ. & Le peuple aussi participe tous aux élections. \\
\hline
\end{tabularx}
\end{center}

\courssub{Le piège n° 2 : la place de 也 et de 都}

En français, « aussi » et « tous » se placent après le verbe ou en fin de phrase : « les citoyens votent \emph{aussi} », « ils votent \emph{tous} ». En chinois, \zh{也} et \zh{都} sont des adverbes : ils se placent \textbf{toujours avant le verbe} (et avant le verbe modal s'il y en a un), jamais en fin de phrase. L'ordre entre \zh{也} et \zh{都} est généralement \zh{也都} (scope : « aussi tous ») ou \zh{都也} (« tous aussi ») selon la nuance voulue.

\begin{exemplebox}[也 et 都 avant le verbe (ou le modal)]
\zh{年轻人也关心民主。} \emph{Niánqīngrén yě guānxīn mínzhǔ.} « Les jeunes s'intéressent aussi à la démocratie. »\\
\zh{我们都应该投票。} \emph{Wǒmen dōu yīnggāi tóupiào.} « Nous devons tous voter. » (\zh{都} avant le modal)\\
\zh{公民也都有权利。} \emph{Gōngmín yě dōu yǒu quánlì.} « Les citoyens aussi ont tous des droits. »
\end{exemplebox}

\begin{attention}[Jamais 也 / 都 en fin de phrase]
Faute typique : \zh{我们投票都。} pour « nous votons tous ». En chinois, l'adverbe précède le verbe : \zh{我们都投票。} De même, « moi aussi » se dit \zh{我也} (et non \zh{我…也} en fin de phrase) : \zh{我也认为民主很重要。} \emph{Wǒ yě rènwéi mínzhǔ hěn zhòngyào.} « Moi aussi, je pense que la démocratie est importante. »
\end{attention}

\courssub{Exercice résolu et entraînement}

\begin{exoresolu}[Thème guidé : trois phrases]
Traduisez en chinois : 1) « La liberté est très importante. » 2) « Les citoyens doivent aussi respecter la loi. » 3) « Parce que nous aimons la liberté, nous votons tous. »
\tcblower
\textbf{1)} Prédicat adjectival : pas de \zh{是}. « Très » = \zh{非常} (plus fort que \zh{很}) ou \zh{很}. $\rightarrow$ \zh{自由非常重要。} \emph{Zìyóu fēicháng zhòngyào.}\\
\textbf{2)} Modal \zh{应该} avant le verbe \zh{尊重} ; adverbe \zh{也} avant le modal : \zh{也应该}. $\rightarrow$ \zh{公民也应该尊重法律。} \emph{Gōngmín yě yīnggāi zūnzhòng fǎlǜ.}\\
\textbf{3)} Paire \zh{因为……所以……} ; \zh{都} avant le verbe \zh{投票}. $\rightarrow$ \zh{因为我们爱自由，所以我们都投票。} \emph{Yīnwèi wǒmen ài zìyóu, suǒyǐ wǒmen dōu tóupiào.}
\end{exoresolu}

\begin{exercice}[À vous de traduire]
Traduisez en chinois (caractères + pinyin) :
\begin{enumerate}
\item « Les élections sont importantes. »
\item « Le peuple a le droit de choisir ses dirigeants. »
\item « Nous devons tous protéger la démocratie. »
\item « Parce que les citoyens votent, la démocratie est vivante. »
\end{enumerate}
\end{exercice}

\begin{corrige}[Exercice : thème]
\textbf{1)} \zh{选举很重要。} \emph{Xuǎnjǔ hěn zhòngyào.} (adjectif prédicat, sans \zh{是})\\
\textbf{2)} \zh{人民有选择自己\textbf{的}领导人的权利。} \emph{Rénmín yǒu xuǎnzé zìjǐ \textbf{de} lǐngdǎorén de quánlì.} (\zh{的} obligatoire entre le complément verbal \zh{选择自己\textbf{的}领导人} et le nom \zh{权利})\\
\textbf{3)} \zh{我们都应该保护民主。} \emph{Wǒmen dōu yīnggāi bǎohù mínzhǔ.} (\zh{都} avant le modal \zh{应该})\\
\textbf{4)} \zh{因为公民投票，所以民主很有活力。} \emph{Yīnwèi gōngmín tóupiào, suǒyǐ mínzhǔ hěn yǒu huólì.} (paire de connecteurs ; « vivante » rendu par \zh{有活力} « avoir de la vitalité »)
\end{corrige}

\begin{aretenir}[Les cinq réflexes du thème]
\begin{enumerate}
\item « être + adjectif » $\rightarrow$ \zh{很 + adjectif} (ou \zh{非常}, \zh{真}, \zh{不}), \textbf{jamais} \zh{是}.
\item « être + nom » $\rightarrow$ \zh{是} + nom.
\item Verbe modal (\zh{应该}, \zh{能}, \zh{会}, \zh{想}, \zh{必须}) : toujours \textbf{avant} le verbe principal.
\item \zh{也} et \zh{都} (et \zh{很}, \zh{不}, \zh{只}, \zh{就}…) : toujours \textbf{avant} le verbe (ou le modal).
\item « parce que… donc… » : paire complète \zh{因为……所以……}, et répéter le nom (sujet) plutôt que d'utiliser un pronom dans le style formel.
\end{enumerate}
\end{aretenir}

Ces réflexes acquis, nous ferons le chemin inverse dans la section suivante : la version, du chinois vers le français.

\coursec{Version : chinois → français}

Dans la section précédente, nous avons appris à passer du français au chinois (le \cle{thème}) : partir d'une phrase française, choisir les bons mots chinois, respecter l'ordre des mots et les connecteurs. Nous allons maintenant faire le chemin inverse : partir d'un texte \textbf{chinois} et produire un texte \textbf{français} correct et naturel. C'est la \cle{version}. L'exercice semble plus facile — on comprend le chinois, on connaît le français — mais il cache des pièges redoutables : le mot à mot, les contresens et les calques de l'ordre chinois. Cette section vous donne une méthode en cinq étapes, un texte d'entraînement complet et les outils pour éviter les erreurs classiques des francophones.

\courssub{Observer d'abord : un texte à traduire}

Lisez le texte suivant \textbf{en entier}, sans dictionnaire, avant toute traduction. C'est le texte de version de cette leçon.

\begin{textebox}[Texte de version — 民主与社会]
在中国，人民也关心民主和社会问题。每个国家有自己的制度。我们应该互相学习，互相尊重。\\
\emph{Zài Zhōngguó, rénmín yě guānxīn mínzhǔ hé shèhuì wèntí. Měi ge guójiā yǒu zìjǐ de zhìdù. Wǒmen yīnggāi hùxiāng xuéxí, hùxiāng zūnzhòng.}\\
« En Chine, le peuple s'intéresse aussi à la démocratie et aux questions sociales. Chaque pays a son propre système. Nous devons apprendre les uns des autres et nous respecter mutuellement. »
\end{textebox}

\textbf{Questions de compréhension (avant de traduire) :}
\begin{enumerate}
\item De quel pays parle la première phrase ?
\item Qui est le sujet de la deuxième phrase ? Que possède-t-il ?
\item Quels sont les deux verbes accompagnés de \zh{互相} dans la dernière phrase ?
\end{enumerate}

\begin{center}
\begin{tabularx}{\linewidth}{|X|X|X|X|}
\hline
\rowcolor{popblueL} Caractères & Pinyin & Nature & Traduction \\
\hline
关心 & guānxīn & verbe & s'intéresser à, se soucier de \\
\hline
制度 & zhìdù & nom & système, régime, institution \\
\hline
互相 & hùxiāng & adverbe & mutuellement, l'un l'autre \\
\hline
尊重 & zūnzhòng & verbe & respecter \\
\hline
人民 & rénmín & nom & le peuple \\
\hline
社会 & shèhuì & nom & la société \\
\hline
问题 & wèntí & nom & question, problème \\
\hline
每个 & měi ge & dét. + classif. & chaque \\
\hline
自己 & zìjǐ & pronom & soi-même, son propre \\
\hline
应该 & yīnggāi & verbe modal & devoir (il faut) \\
\hline
\end{tabularx}
\end{center}

\courssub{La méthode de la version en cinq étapes}

\begin{methode}[Comment réussir une version]
\begin{enumerate}
\item \textbf{Lecture globale.} Lisez le texte chinois entier deux fois. Identifiez le thème général (ici : la démocratie et le respect entre les pays). Ne traduisez encore rien.
\item \textbf{Repérage.} Soulignez les connecteurs (\zh{也}, \zh{和}, \zh{应该}…), les adverbes (\zh{互相}), les pronoms (\zh{自己}) et les mots nouveaux. Devinez le sens des mots inconnus grâce au contexte.
\item \textbf{Traduction phrase par phrase.} Traduisez chaque phrase séparément, en identifiant d'abord le sujet, puis le verbe, puis les compléments. À ce stade, un français un peu « brut » est acceptable.
\item \textbf{Reformulation.} Réécrivez chaque phrase en français naturel : supprimez les mots à mot, replacez les adverbes à la française, choisissez des tournures idiomatiques.
\item \textbf{Relecture.} Vérifiez que rien du texte chinois n'a été oublié ni ajouté, et que le français est correct (orthographe, accords, ponctuation).
\end{enumerate}
\end{methode}

\begin{popfigure}
\begin{tikzpicture}[node distance=0.35cm]
\node[draw=popblue, fill=popblueL, rounded corners, minimum width=3.1cm, minimum height=0.9cm, align=center, font=\small] (a) {\textbf{1. Lecture globale}\\ comprendre le thème};
\node[draw=poporange, fill=poporangeL, rounded corners, minimum width=3.1cm, minimum height=0.9cm, align=center, font=\small, right=of a] (b) {\textbf{2. Repérage}\\ connecteurs, mots-clés};
\node[draw=poppurple, fill=poppurpleL, rounded corners, minimum width=3.1cm, minimum height=0.9cm, align=center, font=\small, right=of b] (c) {\textbf{3. Traduction}\\ phrase par phrase};
\node[draw=popgreen, fill=popgreenL, rounded corners, minimum width=3.1cm, minimum height=0.9cm, align=center, font=\small, below=0.5cm of c] (d) {\textbf{4. Reformulation}\\ français naturel};
\node[draw=poppink, fill=poppinkL, rounded corners, minimum width=3.1cm, minimum height=0.9cm, align=center, font=\small, left=of d] (e) {\textbf{5. Relecture}\\ vérification finale};
\draw[-{Stealth}, thick, popblue] (a) -- (b);
\draw[-{Stealth}, thick, poporange] (b) -- (c);
\draw[-{Stealth}, thick, poppurple] (c) -- (d);
\draw[-{Stealth}, thick, popgreen] (d) -- (e);
\end{tikzpicture}
\legende{Frise des cinq étapes de la version : du chinois vers un français naturel.}
\end{popfigure}

\begin{attention}[Le piège du mot isolé]
Ne traduisez \textbf{jamais} un mot chinois isolé sans avoir lu la phrase entière. Le sens d'un mot dépend du contexte : \zh{问题} peut être « question » ou « problème » ; \zh{制度} peut être « système », « régime » ou « institution ». Dans notre texte, \zh{社会问题} se traduit naturellement par « questions sociales » (et non « problèmes de la société », qui est un calque lourd). Lisez d'abord, traduisez ensuite.
\end{attention}

\courssub{Traduction commentée, phrase par phrase}

\textbf{Phrase 1 :} \zh{在中国，人民也关心民主和社会问题。}

\emph{Zài Zhōngguó, rénmín yě guānxīn mínzhǔ hé shèhuì wèntí.}

\begin{itemize}
\item \zh{在中国} : complément de lieu en tête de phrase → « En Chine » (en français aussi, on peut le mettre en tête : parfait).
\item \zh{人民} : « le peuple » (et non « les peuples » : en chinois, pas de marque de pluriel).
\item \zh{也} : « aussi ». En français, il se place devant le verbe : « s'intéresse \textbf{aussi} ».
\item \zh{关心} : « s'intéresser à ». Attention : ce verbe chinois est \emph{actif}, mais son équivalent français est \emph{pronominal}. Traduction brute « le peuple aussi intéresse la démocratie » → contresens ! Reformulation : « le peuple s'intéresse à la démocratie ».
\end{itemize}

Traduction finale : « En Chine, le peuple s'intéresse aussi à la démocratie et aux questions sociales. »

\textbf{Phrase 2 :} \zh{每个国家有自己的制度。}

\emph{Měi ge guójiā yǒu zìjǐ de zhìdù.}

\begin{itemize}
\item \zh{每个国家} : « chaque pays ».
\item \zh{有} : « a » (possession).
\item \zh{自己的制度} : « son propre système ».
\end{itemize}

Traduction finale : « Chaque pays a son propre système. »

\begin{attention}[自己 ne veut pas toujours dire « soi-même »]
Dans \zh{每个国家有自己的制度}, le pronom \zh{自己} renvoie au \textbf{sujet de la phrase}, c'est-à-dire \zh{每个国家} (« chaque pays »), et non à une personne. Il sert simplement à insister sur la possession : « son \textbf{propre} système ». Ne traduisez donc pas « chaque pays a le système de soi-même », ce qui serait incompréhensible. Règle pratique : \zh{sujet + 自己的 + nom} = « son/sa propre + nom » (le « propre » renvoyant au sujet).
\end{attention}

\textbf{Phrase 3 :} \zh{我们应该互相学习，互相尊重。}

\emph{Wǒmen yīnggāi hùxiāng xuéxí, hùxiāng zūnzhòng.}

\begin{itemize}
\item \zh{我们} : « nous » ; \zh{应该} : « devons ».
\item \zh{互相学习} : littéralement « mutuellement apprendre » → reformulation : « apprendre les uns des autres ».
\item \zh{互相尊重} : littéralement « mutuellement respecter » → « nous respecter mutuellement ».
\end{itemize}

Traduction finale : « Nous devons apprendre les uns des autres et nous respecter mutuellement. »

\begin{propriete}[Place et emploi de 互相]
\zh{互相} est un \textbf{adverbe} : il se place \textbf{toujours avant le verbe} qu'il modifie, jamais après.

Structure : \textbf{Sujet + 互相 + Verbe}

\begin{center}
\begin{tabularx}{\linewidth}{|X|X|X|}
\hline
\rowcolor{popblueL} Structure & Exemple (caractères + pinyin) & Traduction \\
\hline
Sujet + 互相 + V & 我们互相学习。Wǒmen hùxiāng xuéxí. & Nous apprenons les uns des autres. \\
\hline
Sujet + 互相 + V & 他们互相帮助。Tāmen hùxiāng bāngzhù. & Ils s'entraident. \\
\hline
Sujet + modal + 互相 + V & 我们应该互相尊重。Wǒmen yīnggāi hùxiāng zūnzhòng. & Nous devons nous respecter mutuellement. \\
\hline
Sujet + 互相 + V (réciproque) & 两国人民互相了解。Liǎng guó rénmín hùxiāng liǎojiě. & Les peuples des deux pays se comprennent mutuellement. \\
\hline
\end{tabularx}
\end{center}

En français, la « mutualité » s'exprime de plusieurs façons : pronom réciproque (« \textbf{nous} nous respectons »), tournure « les uns des autres », ou adverbe « mutuellement » placé \emph{après} le verbe. L'ordre chinois (互相 + verbe) est donc l'inverse de l'ordre français le plus naturel.
\end{propriete}

\begin{exemplebox}[Exemples traduits avec 互相]
\begin{itemize}
\item \zh{我们应该互相尊重。} \emph{Wǒmen yīnggāi hùxiāng zūnzhòng.} « Nous devons nous respecter mutuellement. »
\item \zh{同学们互相帮助。} \emph{Tóngxuémen hùxiāng bāngzhù.} « Les camarades s'entraident. »
\item \zh{中国和喀麦隆互相学习。} \emph{Zhōngguó hé Kāmàilóng hùxiāng xuéxí.} « La Chine et le Cameroun apprennent l'un de l'autre. »
\item \zh{我们要互相了解。} \emph{Wǒmen yào hùxiāng liǎojiě.} « Nous devons apprendre à nous connaître mutuellement. »
\end{itemize}
\end{exemplebox}

\begin{savaistu}[互相, un mot d'examen]
\zh{互相} (hùxiāng) est un adverbe figé très fréquent dans les épreuves de version et de thème. Il se combine presque toujours avec les mêmes verbes : \zh{互相学习} (apprendre l'un de l'autre), \zh{互相帮助} (s'entraider), \zh{互相尊重} (se respecter mutuellement), \zh{互相了解} (se connaître mutuellement). Mémorisez ces quatre combinaisons comme des blocs : elles reviennent dans presque tous les textes sur la coopération, la citoyenneté et les relations entre les pays.
\end{savaistu}

\courssub{Du mot à mot au français naturel : la reformulation}

La quatrième étape est celle qui fait la différence entre une copie moyenne et une excellente copie. Voici le contraste entre traduction littérale (à proscrire) et traduction reformulée (attendue).

\begin{center}
\begin{tabularx}{\linewidth}{|X|X|X|}
\hline
\rowcolor{popblueL} Phrase chinoise & Mot à mot (mauvais) & Français naturel (bon) \\
\hline
人民也关心民主 & le peuple aussi intéresse la démocratie & le peuple s'intéresse aussi à la démocratie \\
\hline
每个国家有自己的制度 & chaque pays a le système de soi-même & chaque pays a son propre système \\
\hline
互相学习 & mutuellement apprendre & apprendre les uns des autres \\
\hline
互相尊重 & mutuellement respecter & se respecter mutuellement \\
\hline
社会问题 & problèmes de la société & questions sociales \\
\hline
\end{tabularx}
\end{center}

\begin{aretenir}[Les trois réflexes de la reformulation]
\begin{enumerate}
\item \textbf{Verbe chinois actif → verbe français pronominal} quand le sens l'exige : \zh{关心} → « s'intéresser à ».
\item \textbf{Adverbe avant le verbe en chinois → après le verbe en français} : \zh{互相尊重} → « se respecter mutuellement ».
\item \textbf{Nom + 的 + nom} : ne pas garder « de » partout ; préférer l'adjectif français quand il existe : \zh{社会问题} → « questions \emph{sociales} ».
\end{enumerate}
\end{aretenir}

\courssub{Deux connecteurs fréquents en version : 虽然……但是…… et 只有……才……}

Ces deux paires de connecteurs reviennent souvent dans les textes de version du thème « citoyenneté ». Leur piège est le même : le chinois utilise \textbf{deux} mots là où le français n'en garde qu'\textbf{un}.

\begin{propriete}[虽然……但是…… : « bien que »]
Structure : \textbf{虽然 + proposition 1，但是 + proposition 2}

En chinois, les deux connecteurs se combinent dans la même phrase ; en français, « bien que » et « mais » sont \textbf{incompatibles} : on ne traduit que \zh{虽然} par « bien que » (suivi du subjonctif) et on supprime \zh{但是}.

\zh{虽然我们的制度不一样，但是我们应该互相学习。} \emph{Suīrán wǒmen de zhìdù bù yíyàng, dànshì wǒmen yīnggāi hùxiāng xuéxí.} « Bien que nos systèmes soient différents, nous devons apprendre les uns des autres. » (et non « bien que… mais… »)
\end{propriete}

\begin{propriete}[只有……才…… : « ce n'est que… que… »]
Structure : \textbf{只有 + condition，才 + résultat}

Cette paire exprime une condition \textbf{nécessaire} : le résultat n'arrive que si la condition est remplie. En français, on la rend par « ce n'est que… que… », « seulement si… », ou « il faut… pour que… ».

\zh{只有互相尊重，世界才会更美好。} \emph{Zhǐyǒu hùxiāng zūnzhòng, shìjiè cái huì gèng měihǎo.} « Ce n'est qu'en nous respectant mutuellement que le monde deviendra plus beau. »
\end{propriete}

\courssub{Entraînement : un texte de version complet}

\begin{textebox}[Pour s'entraîner — 不同的国家]
世界上有很多国家，每个国家有自己的历史、文化和制度。中国是一个大国，喀麦隆也是一个美丽的国家。两国的人民都关心民主和社会问题。虽然我们的制度不一样，但是我们应该互相学习，互相帮助。中国学生学习法语，喀麦隆学生学习汉语。我们一起上课，一起踢球，互相了解，互相尊重。老师说：不同的国家可以成为好朋友。我觉得这句话很对。只有互相尊重，世界才会更美好。\\
\emph{Shìjiè shang yǒu hěn duō guójiā, měi ge guójiā yǒu zìjǐ de lìshǐ, wénhuà hé zhìdù. Zhōngguó shì yí ge dà guó, Kāmàilóng yě shì yí ge měilì de guójiā. Liǎng guó de rénmín dōu guānxīn mínzhǔ hé shèhuì wèntí. Suīrán wǒmen de zhìdù bù yíyàng, dànshì wǒmen yīnggāi hùxiāng xuéxí, hùxiāng bāngzhù. Zhōngguó xuésheng xuéxí Fǎyǔ, Kāmàilóng xuésheng xuéxí Hànyǔ. Wǒmen yìqǐ shàng kè, yìqǐ tī qiú, hùxiāng liǎojiě, hùxiāng zūnzhòng. Lǎoshī shuō : bùtóng de guójiā kěyǐ chéngwéi hǎo péngyou. Wǒ juéde zhè jù huà hěn duì. Zhǐyǒu hùxiāng zūnzhòng, shìjiè cái huì gèng měihǎo.}\\
« Il existe de nombreux pays dans le monde, et chaque pays a sa propre histoire, sa propre culture et son propre système. La Chine est un grand pays, le Cameroun est aussi un beau pays. Les peuples des deux pays s'intéressent tous deux à la démocratie et aux questions sociales. Bien que nos systèmes soient différents, nous devons apprendre les uns des autres et nous entraider. Les élèves chinois apprennent le français, les élèves camerounais apprennent le chinois. Nous allons en classe ensemble, nous jouons au football ensemble, nous apprenons à nous connaître et nous nous respectons mutuellement. Le professeur dit : des pays différents peuvent devenir de bons amis. Je trouve que cette phrase est très juste. Ce n'est qu'en nous respectant mutuellement que le monde deviendra plus beau. »
\end{textebox}

\textbf{Glossaire :} \zh{历史} lìshǐ (histoire) ; \zh{文化} wénhuà (culture) ; \zh{虽然……但是……} suīrán… dànshì… (bien que…) ; \zh{不一样} bù yíyàng (différent) ; \zh{一起} yìqǐ (ensemble) ; \zh{成为} chéngwéi (devenir) ; \zh{觉得} juéde (trouver, penser) ; \zh{只有……才……} zhǐyǒu… cái… (ce n'est que… que…) ; \zh{美好} měihǎo (beau, magnifique).

\textbf{Questions de compréhension :}
\begin{enumerate}
\item Quels sont les deux pays comparés dans ce texte ?
\item Quelle langue les élèves camerounais apprennent-ils ?
\item Relevez toutes les occurrences de \zh{互相} et traduisez chaque groupe verbal.
\item Que faut-il faire, selon la dernière phrase, pour que le monde devienne plus beau ?
\end{enumerate}

\begin{exoresolu}[Version guidée]
Énoncé : traduisez en français naturel la phrase suivante, en appliquant les cinq étapes de la méthode.

\zh{虽然中国和喀麦隆的制度不一样，但是两国人民应该互相学习，互相尊重。}

\emph{Suīrán Zhōngguó hé Kāmàilóng de zhìdù bù yíyàng, dànshì liǎng guó rénmín yīnggāi hùxiāng xuéxí, hùxiāng zūnzhòng.}
\tcblower
Solution détaillée :
\begin{enumerate}
\item \textbf{Lecture globale} : la phrase compare deux pays et exprime une obligation morale.
\item \textbf{Repérage} : connecteur \zh{虽然……但是……} (bien que…), adverbe \zh{互相} (deux fois), modal \zh{应该} (devoir).
\item \textbf{Traduction brute} : « Bien que la Chine et le Cameroun leurs systèmes pas pareils, mais les peuples des deux pays doivent mutuellement apprendre, mutuellement respecter. »
\item \textbf{Reformulation} : en français, on ne dit pas « bien que… mais… » : on garde seulement « bien que » (avec le subjonctif). « Mutuellement apprendre » devient « apprendre les uns des autres » ; « mutuellement respecter » devient « se respecter mutuellement ».
\item \textbf{Traduction finale} : « Bien que la Chine et le Cameroun aient des systèmes différents, les peuples des deux pays doivent apprendre les uns des autres et se respecter mutuellement. »
\end{enumerate}
Remarque : \zh{制度不一样} se reformule élégamment par « des systèmes différents » plutôt que « les systèmes ne sont pas pareils ».
\end{exoresolu}

\begin{exercice}[À vous de jouer]
Traduisez en français naturel, en rédigeant vos cinq étapes sur le brouillon :
\begin{enumerate}
\item \zh{每个国家有自己的文化。}\emph{Měi ge guójiā yǒu zìjǐ de wénhuà.}
\item \zh{我们应该互相帮助。}\emph{Wǒmen yīnggāi hùxiāng bāngzhù.}
\item \zh{喀麦隆人民也关心社会问题。}\emph{Kāmàilóng rénmín yě guānxīn shèhuì wèntí.}
\item \zh{只有互相了解，我们才能成为好朋友。}\emph{Zhǐyǒu hùxiāng liǎojiě, wǒmen cái néng chéngwéi hǎo péngyou.}
\end{enumerate}
\end{exercice}

\begin{corrige}[Exercice « À vous de jouer »]
\begin{enumerate}
\item « Chaque pays a sa propre culture. » (\zh{自己} renvoie au sujet \zh{每个国家} : « sa \textbf{propre} culture », pas « la culture de soi-même ».)
\item « Nous devons nous entraider. » ou « Nous devons nous aider mutuellement. » (\zh{互相帮助} : les deux reformulations sont correctes ; « s'entraider » est la plus naturelle.)
\item « Le peuple camerounais s'intéresse aussi aux questions sociales. » (\zh{关心} → verbe pronominal « s'intéresser à » ; \zh{也} → « aussi » devant le verbe.)
\item « Ce n'est qu'en apprenant à nous connaître mutuellement que nous pourrons devenir de bons amis. » (\zh{只有……才……} → « ce n'est que… que… » ; \zh{能} → « pouvoir », au futur ou au conditionnel selon la tournure choisie.)
\end{enumerate}
\end{corrige}

\begin{attention}[Les trois fautes qui coûtent des points en version]
\begin{enumerate}
\item \textbf{Le mot à mot} : « mutuellement apprendre » n'existe pas en français. Reformulez toujours.
\item \textbf{Le contresens sur 关心} : \zh{人民关心民主} ne veut pas dire « le peuple inquiète la démocratie » mais « le peuple \textbf{s'intéresse à} la démocratie ».
\item \textbf{L'oubli d'un élément} : chaque mot chinois doit avoir sa trace dans le français (et rien ne doit être ajouté). Relisez en comparant les deux textes, phrase par phrase.
\end{enumerate}
\end{attention}

\begin{aretenir}[Bilan de la section]
La version se réussit en cinq étapes : lecture globale, repérage, traduction phrase par phrase, reformulation, relecture. Quatre points de grammaire à retenir : \zh{互相} se place \textbf{avant le verbe} en chinois mais se reformule \textbf{après} en français ; \zh{自己的} renvoie au sujet (« son propre ») ; \zh{关心} se traduit par le verbe pronominal « s'intéresser à » ; \zh{虽然……但是……} ne se traduit que par « bien que » (jamais « bien que… mais… »), et \zh{只有……才……} par « ce n'est que… que… ». Une bonne version n'est pas un mot à mot : c'est un texte français naturel qui dit exactement ce que dit le texte chinois, ni plus, ni moins.
\end{aretenir}

\coursec{Production guidée et grille d'évaluation}

Après avoir traduit du chinois vers le français dans la section précédente, nous franchissons la dernière étape du parcours : produire \emph{nous-mêmes} un texte en chinois. C'est l'épreuve reine de l'expression écrite au Baccalauréat : on vous donne un sujet, une trame, et vous devez rédiger un texte court, correct et cohérent. Cette section vous donne la méthode complète, une banque de phrases prêtes à l'emploi, un modèle commenté et la grille d'évaluation officielle.

\courssub{Le sujet et la trame imposée}

\begin{definition}[Le sujet de production]
\textbf{Sujet :} \zh{民主和年轻人} (\emph{Mínzhǔ hé niánqīngrén}) — « La démocratie et les jeunes ».\\
\textbf{Longueur :} 5 à 8 phrases.\\
\textbf{Trame imposée :} thèse + \zh{首先} + \zh{其次} + \zh{但是} + \zh{所以}.
\end{definition}

La \cle{trame} est le squelette du texte : elle vous indique l'ordre des idées et les connecteurs obligatoires. La respecter, c'est déjà gagner des points.

\begin{center}
\begin{tabularx}{\linewidth}{|X|X|X|}
\hline
\rowcolor{popblueL} Élément de la trame & Rôle dans le texte & Équivalent français \\
\hline
Thèse (1 phrase) & Annoncer le sujet et votre idée principale & « La démocratie est importante. » \\
\hline
\zh{首先} shǒuxiān & Premier argument & « d'abord » \\
\hline
\zh{其次} qícì & Deuxième argument & « ensuite » \\
\hline
\zh{但是} dànshì & Nuance, difficulté, concession & « mais » \\
\hline
\zh{所以} suǒyǐ & Conclusion, conséquence logique & « donc » \\
\hline
\end{tabularx}
\end{center}

\courssub{Banque de phrases à réutiliser}

Au Bac, on ne part jamais de zéro : on \emph{réutilise} des phrases modèles apprises par cœur, en changeant un ou deux mots. Voici la banque officielle de cette leçon, enrichie d'un exemple contextualisé au Cameroun.

\begin{center}
\begin{tabularx}{\linewidth}{|X|X|X|}
\hline
\rowcolor{popblueL} Phrase modèle & Pinyin & Traduction \\
\hline
\zh{民主很重要。} & Mínzhǔ hěn zhòngyào. & La démocratie est très importante. \\
\hline
\zh{年轻人应该关心社会。} & Niánqīngrén yīnggāi guānxīn shèhuì. & Les jeunes doivent s'intéresser à la société. \\
\hline
\zh{我们有选举的权利。} & Wǒmen yǒu xuǎnjǔ de quánlì. & Nous avons le droit de voter. \\
\hline
\zh{我们有说话的自由。} & Wǒmen yǒu shuōhuà de zìyóu. & Nous avons la liberté de parler. \\
\hline
\zh{所以，我们应该学习民主。} & Suǒyǐ, wǒmen yīnggāi xuéxí mínzhǔ. & Donc, nous devons apprendre la démocratie. \\
\hline
\zh{年轻人也应该参加选举。} & Niánqīngrén yě yīnggāi cānjiā xuǎnjǔ. & Les jeunes doivent aussi participer aux élections. \\
\hline
\zh{喀麦隆的年轻人关心选举。} & Kāmàilóng de niánqīngrén guānxīn xuǎnjǔ. & Les jeunes du Cameroun s'intéressent aux élections. \\
\hline
\end{tabularx}
\end{center}

\begin{exemplebox}[Exemples traduits à imiter]
\zh{年轻人应该关心社会。} \emph{(Niánqīngrén yīnggāi guānxīn shèhuì.)} « Les jeunes doivent s'intéresser à la société. »\\
\zh{首先，民主很重要。} \emph{(Shǒuxiān, mínzhǔ hěn zhòngyào.)} « D'abord, la démocratie est très importante. »\\
\zh{但是，有些年轻人不关心政治。} \emph{(Dànshì, yǒuxiē niánqīngrén bù guānxīn zhèngzhì.)} « Mais certains jeunes ne s'intéressent pas à la politique. »\\
\zh{所以，我们应该帮助年轻人了解民主。} \emph{(Suǒyǐ, wǒmen yīnggāi bāngzhù niánqīngrén liǎojiě mínzhǔ.)} « Donc, nous devons aider les jeunes à comprendre la démocratie. »
\end{exemplebox}

\courssub{Modèle de texte rédigé et commenté}

\begin{textebox}[Modèle : 民主和年轻人]
\zh{民主和年轻人}\\
\zh{民主很重要，年轻人也很重要。}\\
\zh{首先，年轻人有选举的权利，也有说话的自由。}\\
\zh{其次，年轻人应该关心社会，学习民主。}\\
\zh{但是，有些年轻人不喜欢政治，也不参加选举。}\\
\zh{所以，学校和家庭应该帮助年轻人了解民主。}\\
\zh{民主的未来在年轻人的手里。}\\[0.5em]
\emph{Mínzhǔ hé niánqīngrén}\\
\emph{Mínzhǔ hěn zhòngyào, niánqīngrén yě hěn zhòngyào.}\\
\emph{Shǒuxiān, niánqīngrén yǒu xuǎnjǔ de quánlì, yě yǒu shuōhuà de zìyóu.}\\
\emph{Qícì, niánqīngrén yīnggāi guānxīn shèhuì, xuéxí mínzhǔ.}\\
\emph{Dànshì, yǒuxiē niánqīngrén bù xǐhuan zhèngzhì, yě bù cānjiā xuǎnjǔ.}\\
\emph{Suǒyǐ, xuéxiào hé jiātíng yīnggāi bāngzhù niánqīngrén liǎojiě mínzhǔ.}\\
\emph{Mínzhǔ de wèilái zài niánqīngrén de shǒuli.}\\[0.5em]
« La démocratie et les jeunes. La démocratie est très importante, et les jeunes sont très importants aussi. D'abord, les jeunes ont le droit de voter et la liberté de parole. Ensuite, les jeunes doivent s'intéresser à la société et apprendre la démocratie. Mais certains jeunes n'aiment pas la politique et ne participent pas aux élections. Donc, l'école et la famille doivent aider les jeunes à comprendre la démocratie. L'avenir de la démocratie est entre les mains des jeunes. »
\end{textebox}

Observons comment ce modèle respecte la trame : phrase 1 = thèse ; \zh{首先} introduit le premier argument (les droits) ; \zh{其次} le deuxième (les devoirs) ; \zh{但是} la nuance ; \zh{所以} la conclusion ; la dernière phrase est une \cle{chute} courte qui donne du relief au texte. Sept phrases au total : dans la fourchette demandée (5 à 8).

\courssub{Les trois étapes de la production en classe}

\begin{methode}[Réussir sa production écrite en trois étapes]
\begin{enumerate}
\item \textbf{Étape 1 — Plan en français (2 minutes).} Sur le brouillon, notez une idée par ligne : thèse / argument 1 / argument 2 / nuance / conclusion. Ne cherchez pas d'idées originales : des idées simples bien exprimées valent mieux que des idées brillantes intraduisibles.
\item \textbf{Étape 2 — Rédaction (15 minutes).} Rédigez en phrases courtes : Sujet + verbe + complément. Réutilisez la banque de phrases. Un connecteur de la trame par phrase clé : \zh{首先}、\zh{其次}、\zh{但是}、\zh{所以}.
\item \textbf{Étape 3 — Relecture avec la grille (5 minutes).} Relisez en vous mettant à la place de l'examinateur : caractères, grammaire, connecteurs, longueur, présentation. Réservez au moins 3 minutes aux seuls caractères.
\end{enumerate}
\end{methode}

\begin{propriete}[Ce que l'examinateur évalue vraiment]
Quatre critères guident la correction au Bac :
\begin{itemize}
\item l'\cle{exactitude des caractères} : aucun trait manquant, aucune confusion (主/王, 民/氏) ;
\item la \cle{correction grammaticale} : ordre des mots, adjectif sans \zh{是}, place de \zh{也} et \zh{应该} avant le verbe ;
\item la \cle{cohérence du texte} : les connecteurs de la trame sont présents et logiques ;
\item la \cle{richesse raisonnable du lexique} : le vocabulaire du thème (民主、权利、自由、选举、社会) est utilisé, sans répétition excessive.
\end{itemize}
\end{propriete}

\begin{attention}[Chaque caractère mal écrit coûte des points]
Un trait manquant transforme le caractère en un autre mot : \zh{主} (zhǔ, « principal ») sans son point devient \zh{王} (wáng, « roi ») ; \zh{民} mal fermé (dernier trait horizontal manquant) devient illisible ou confondu avec d'autres caractères. De même, ne dites jamais \zh{民主是很重要} : l'adjectif seul fait office de verbe, on écrit \zh{民主很重要}. Et rappelez-vous : \zh{也} et \zh{应该} se placent \emph{avant} le verbe, jamais en fin de phrase. Réservez toujours 3 minutes à la relecture des caractères, trait par trait.
\end{attention}

\courssub{Grille d'évaluation sur 20 points}

\begin{popfigure}
\begin{tikzpicture}[x=1cm,y=1cm]
\fill[popblueL] (0,0) rectangle (10.5,1);
\node[anchor=west,font=\bfseries] at (0.2,0.5) {Critère};
\node[anchor=west,font=\bfseries] at (7.2,0.5) {Barème};
\fill[popcream] (0,-1) rectangle (10.5,0);
\node[anchor=west] at (0.2,-0.5) {Caractères corrects (traits, radicaux, 主/王)};
\node[anchor=west] at (7.2,-0.5) {/ 4};
\fill[white] (0,-2) rectangle (10.5,-1);
\node[anchor=west] at (0.2,-1.5) {Grammaire et ordre des mots (也, 应该, adjectif)};
\node[anchor=west] at (7.2,-1.5) {/ 4};
\fill[popcream] (0,-3) rectangle (10.5,-2);
\node[anchor=west] at (0.2,-2.5) {Connecteurs de la trame (首先、其次、但是、所以)};
\node[anchor=west] at (7.2,-2.5) {/ 4};
\fill[white] (0,-4) rectangle (10.5,-3);
\node[anchor=west] at (0.2,-3.5) {Respect du sujet et de la longueur (5 à 8 phrases)};
\node[anchor=west] at (7.2,-3.5) {/ 4};
\fill[popcream] (0,-5) rectangle (10.5,-4);
\node[anchor=west] at (0.2,-4.5) {Présentation (netteté, ponctuation chinoise)};
\node[anchor=west] at (7.2,-4.5) {/ 4};
\fill[poporangeL] (0,-6) rectangle (10.5,-5);
\node[anchor=west,font=\bfseries] at (0.2,-5.5) {Total};
\node[anchor=west,font=\bfseries] at (7.2,-5.5) {/ 20};
\draw[popline,thick] (0,0) rectangle (10.5,-6);
\draw[popline] (0,-1)--(10.5,-1) (0,-2)--(10.5,-2) (0,-3)--(10.5,-3) (0,-4)--(10.5,-4) (0,-5)--(10.5,-5);
\draw[popline] (7,-6)--(7,1);
\end{tikzpicture}
\legende{Grille d'évaluation de la production écrite : cinq critères, quatre points chacun, total sur 20.}
\end{popfigure}

\courssub{Auto-correction : la check-list de relecture}

\begin{exoresolu}[Auto-corriger un texte d'élève]
Énoncé. Voici trois phrases écrites par un élève. Repérez et corrigez les erreurs à l'aide de la check-list.\\
1. \zh{民主是重要。} 2. \zh{年轻人关心应该社会。} 3. \zh{我们也参加选举也。}
\tcblower
Solution détaillée.\\
1. \zh{民主是重要。} : l'adjectif \zh{重要} ne prend pas \zh{是}. Correction : \zh{民主很重要。} « La démocratie est très importante. »\\
2. \zh{年轻人关心应该社会。} : \zh{应该} doit précéder le verbe \zh{关心}. Correction : \zh{年轻人应该关心社会。} « Les jeunes doivent s'intéresser à la société. »\\
3. \zh{我们也参加选举也。} : \zh{也} ne se place jamais en fin de phrase ; un seul \zh{也}, avant le verbe. Correction : \zh{我们也参加选举。} « Nous participons aussi aux élections. »
\end{exoresolu}

\begin{methode}[Check-list d'auto-correction avant de rendre la copie]
\begin{enumerate}
\item \textbf{Caractères :} \zh{主} a bien son point (et non \zh{王}) ; \zh{民} est bien fermé (dernier trait horizontal) ; \zh{选举} et \zh{权利} sont complets.
\item \textbf{Adjectif sans 是 :} \zh{民主很重要}, jamais \zh{民主是很重要}.
\item \textbf{Place de 也 et 应该 :} toujours avant le verbe (\zh{也应该关心}).
\item \textbf{Connecteurs :} \zh{首先}、\zh{其次}、\zh{但是}、\zh{所以} sont tous présents, suivis d'une virgule chinoise \zh{，}.
\item \textbf{Longueur :} comptez vos phrases : entre 5 et 8.
\end{enumerate}
\end{methode}

\begin{experience}[Production guidée en classe entière]
Par groupes de deux, rédigez votre texte « \zh{民主和年轻人} » en suivant les trois étapes : 2 minutes de plan en français, 15 minutes de rédaction en chinois, 5 minutes de relecture croisée (chaque binôme corrige la copie de l'autre avec la grille sur 20). Le professeur ramasse les grilles et compare les notes attribuées : c'est un excellent entraînement à l'œil de l'examinateur.
\end{experience}

\begin{exercice}[À vous de produire]
Rédigez votre propre texte « \zh{民主和年轻人} » (5 à 8 phrases) en respectant la trame : thèse + \zh{首先} + \zh{其次} + \zh{但是} + \zh{所以}. Réutilisez au moins trois phrases de la banque. Puis évaluez-vous avec la grille sur 20 et notez votre score pour chaque critère.
\end{exercice}

\begin{corrige}[Exercice : production « 民主和年轻人 »]
Exemple de copie notée 18/20 :\\
\zh{民主很重要。首先，我们有选举的权利，也有说话的自由。其次，年轻人应该关心社会，也应该学习民主。但是，有些年轻人不喜欢政治。所以，我们应该帮助年轻人了解民主。民主的未来在我们手里。}\\
Caractères : 4/4 (aucune confusion 主/王). Grammaire : 4/4 (adjectif sans 是, 也 et 应该 bien placés). Connecteurs : 4/4 (trame complète). Sujet et longueur : 4/4 (6 phrases, sur le sujet). Présentation : 2/4 si la ponctuation française est utilisée à la place de \zh{，。} — pensez à la ponctuation chinoise pleine chasse.
\end{corrige}

\begin{aretenir}[L'essentiel de la production écrite]
Un plan simple en français, des phrases courtes et correctes, la trame respectée (\zh{首先}、\zh{其次}、\zh{但是}、\zh{所以}), trois minutes de relecture des caractères : c'est la recette d'une copie réussie. La banque de phrases du thème est votre trésor : apprenez-la par cœur et réutilisez-la sans hésiter.
\end{aretenir}

\newpage

\coursec{Activité d'intégration}

\coursec{Leçon 23 — Expression écrite : la démocratie}

\courssub{Objectifs de la leçon}
À l'issue de cette leçon, tu seras capable de :
\begin{itemize}
\item produire une \cle{lettre} argumentative courte (6 à 8 phrases) en chinois sur le thème de la \cle{démocratie}, respectant les conventions épistolaires ;
\item structurer un raisonnement simple : \cle{thèse} (définition), deux \cle{arguments} (\zh{首先}、\zh{其次}), une \cle{nuance} (\zh{但是}), une \cle{conclusion} (\zh{所以}) ;
\item mobiliser au minimum 6 mots du lexique thématique et 3 connecteurs argumentatifs ;
\item écrire correctement les caractères du thème (ordre des traits, radicaux, discrimination des caractères proches) ;
\item réaliser des exercices de \cle{thème} (français $\rightarrow$ chinois) et de \cle{version} (chinois $\rightarrow$ français) sur le vocabulaire et les structures vus ;
\item évaluer la production d'un pair à l'aide d'une grille critériée et réviser ton propre texte.
\end{itemize}

\courssub{Situation déclenchante : Projet « Jeunes citoyens du monde »}
Ton lycée de Yaoundé participe à un échange avec un lycée de Pékin. Le professeur chinois, M. Wang, a adressé la lettre ci-dessous à la classe. Ta mission : répondre à ton correspondant \zh{李明} (Lǐ Míng) en chinois.

\begin{textebox}[Lettre de M. Wang — 北京]
亲爱的喀麦隆朋友们：\\
Qīn'ài de Kāmàilóng péngyoumen\\
« Chers amis camerounais »\\[4pt]
你们好！我是北京一所高中的汉语老师。\\
Nǐmen hǎo! Wǒ shì Běijīng yì suǒ gāozhōng de Hànyǔ lǎoshī.\\
« Bonjour ! Je suis professeur de chinois dans un lycée de Pékin. »\\[4pt]
我的学生们很想了解喀麦隆。请你们写一封信，谈谈你们国家的民主：\\
Wǒ de xuéshengmen hěn xiǎng liǎojiě Kāmàilóng. Qǐng nǐmen xiě yì fēng xìn, tántan nǐmen guójiā de mínzhǔ:\\
« Mes élèves veulent beaucoup découvrir le Cameroun. Écrivez une lettre pour parler de la démocratie dans votre pays : »\\[4pt]
什么是民主？人民怎么参加选举？年轻人有什么看法？\\
Shénme shì mínzhǔ? Rénmín zěnme cānjiā xuǎnjǔ? Niánqīngrén yǒu shénme kànfǎ?\\
« Qu'est-ce que la démocratie ? Comment le peuple participe-t-il aux élections ? Quelle est l'opinion des jeunes ? »\\[4pt]
期待你们的来信！\\
Qīdài nǐmen de láixìn!\\
« Nous attendons vos lettres avec impatience ! »
\end{textebox}

\courssub{Lexique thématique : la démocratie et la citoyenneté}

\begin{center}
\begin{tabularx}{\linewidth}{|X|X|X|X|}
\hline
\rowcolor{popblueL} Caractères & Pinyin & Nature & Traduction \\
\hline
民主 & mínzhǔ & nom / adj. & démocratie ; démocratique \\
\hline
人民 & rénmín & nom & peuple \\
\hline
选举 & xuǎnjǔ & nom / verbe & élection ; élire \\
\hline
投票 & tóupiào & verbe & voter (litt. « jeter le bulletin ») \\
\hline
权利 & quánlì & nom & droit \\
\hline
自由 & zìyóu & nom / adj. & liberté ; libre \\
\hline
国家 & guójiā & nom & pays, État \\
\hline
看法 & kànfǎ & nom & opinion, point de vue \\
\hline
参加 & cānjiā & verbe & participer à \\
\hline
重要 & zhòngyào & adj. & important \\
\hline
言论自由 & yánlùn zìyóu & nom & liberté d'expression \\
\hline
总统 & zǒngtǒng & nom & président (de la République) \\
\hline
领导人 & lǐngdǎorén & nom & dirigeant, leader \\
\hline
选择 & xuǎnzé & verbe & choisir \\
\hline
政治 & zhèngzhì & nom & politique \\
\hline
完美 & wánměi & adj. & parfait \\
\hline
了解 & liǎojiě & verbe & comprendre, connaître \\
\hline
\end{tabularx}
\end{center}

\begin{propriete}[Analyse des caractères clés et pièges graphiques]
\textbf{1. 民 (mín, « peuple ») vs 氏 (shì, « nom de clan »)} \\
\zh{民} : radical \zh{乙} (yǐ) en haut + \zh{氏} en bas. Ordre : haut $\rightarrow$ bas. \emph{Sens : le peuple, la population.} \\
\zh{氏} : radical \zh{氏} (shì) lui-même. \emph{Sens : nom de famille, clan.} \\
\textit{Astuce} : \zh{民} a un trait horizontal de plus en haut (le peuple est « au-dessus »).

\textbf{2. 选 (xuǎn, « choisir ») vs 远 (yuǎn, « loin »)} \\
\zh{选} : radical \zh{辶} (chuò, « marche ») + \zh{巛} (quán) + \zh{走} (zǒu) contracté. Ordre : partie gauche (\zh{巛}) $\rightarrow$ radical \zh{辶} (point, horizontal, point). \emph{Sens : choisir, élire.} \\
\zh{远} : radical \zh{辶} + \zh{元} (yuán). \emph{Sens : loin, distant.} \\
\textit{Astuce} : pour \zh{choisir} (\zh{选}), il faut « marcher » (\zh{辶}) vers les « options » (\zh{巛}) ; pour \zh{loin} (\zh{远}), on part de l'\zh{origine} (\zh{元}).

\textbf{3. 权 (quán, « droit/pouvoir ») vs 欢 (huān, « joie »)} \\
\zh{权} : radical \zh{木} (mù, « bois ») + \zh{刂} (dāo, « couteau » à droite). Ordre : \zh{木} $\rightarrow$ \zh{刂}. \emph{Sens : pouvoir, droit, autorité (le bois taillé par le couteau = mesure/autorité).} \\
\zh{欢} : radical \zh{欠} (qiàn, « souffle/manque ») + \zh{雚} (guàn) simplifié. \emph{Sens : joie, réjouissance.} \\
\textit{Astuce} : \zh{权} a un couteau (\zh{刂}) pour trancher (le pouvoir) ; \zh{欢} a le souffle (\zh{欠}) de la joie.
\end{propriete}

\courssub{Connecteurs argumentatifs et structures syntaxiques}

\begin{center}
\begin{tabularx}{\linewidth}{|X|X|X|}
\hline
\rowcolor{popblueL} Structure & Exemple (caractères + pinyin) & Traduction \\
\hline
\zh{首先}…, \zh{其次}… & \zh{首先，人民可以投票。其次，我们有言论自由。} Shǒuxiān, rénmín kěyǐ tóupiào. Qícì, wǒmen yǒu yánlùn zìyóu. & D'abord…, ensuite… (énumération d'arguments) \\
\hline
\zh{但是} + phrase & \zh{但是，民主还不完美。} Dànshì, mínzhǔ hái bù wánměi. & Mais… (concession / nuance) \\
\hline
\zh{所以} + conclusion & \zh{所以，我们应该参加。} Suǒyǐ, wǒmen yīnggāi cānjiā. & Donc / C'est pourquoi… (conséquence) \\
\hline
\zh{我认为} / \zh{我觉得} + proposition & \zh{我认为民主很重要。} Wǒ rènwéi mínzhǔ hěn zhòngyào. & Je pense que… / Je trouve que… (opinion) \\
\hline
Sujet + \zh{可以} / \zh{应该} + V & \zh{年轻人应该投票。} Niánqīngrén yīnggāi tóupiào. & Les jeunes peuvent / doivent voter (capacité / devoir moral) \\
\hline
\zh{在} + Lieu + Sujet + V & \zh{在喀麦隆，人民参加选举。} Zài Kāmàilóng, rénmín cānjiā xuǎnjǔ. & Au Cameroun, le peuple participe aux élections (complément de lieu \textbf{avant} le verbe) \\
\hline
\zh{对}…\zh{感兴趣} & \zh{年轻人对政治不感兴趣。} Niánqīngrén duì zhèngzhì bù gǎn xìngqù. & Les jeunes ne s'intéressent pas à la politique (structure figée) \\
\hline
\end{tabularx}
\end{center}

\begin{propriete}[Règles grammaticales essentielles pour la lettre]
\textbf{1. Position du complément de lieu (\zh{在} + Lieu)} \\
En chinois, le lieu se place \textbf{avant} le verbe (et non après comme en français). \\
Schéma : \cle{Sujet + 在 + Lieu + Verbe + Complément} \\
Ex. : \zh{我们在喀麦隆投票。} (Wǒmen zài Kāmàilóng tóupiào.) — « Nous votons au Cameroun. »

\textbf{2. Verbes modaux \zh{可以} (kěyǐ) vs \zh{应该} (yīnggāi)} \\
\zh{可以} = permission, possibilité objective (« pouvoir », « avoir le droit de »). \\
\zh{应该} = obligation morale, conseil (« devoir », « devrait »). \\
Dans une lettre d'opinion, \zh{应该} exprime l'engagement citoyen.

\textbf{3. Structure \zh{就是} (jiùshì) pour définir} \\
\zh{民主就是人民当家作主。} (Mínzhǔ jiùshì rénmín dāngjiāzuòzhǔ.) — « La démocratie, \textbf{c'est} le peuple qui est maître chez lui. » \\
\zh{就是} insiste sur l'identité / l'explication (plus fort que \zh{是} seul).

\textbf{4. Négation de \zh{感兴趣} (gǎn xìngqù)} \\
On nie \zh{感} : \zh{不感兴趣} (bù gǎn xìngqù). Pas \zh{没感兴趣}.
\end{propriete}

\begin{attention}[Erreurs fréquentes des francophones — \textbf{À éviter absolument}]
\begin{itemize}
\item \textbf{« Chez moi » $\neq$ \zh{在我家} (wǒ jiā = ma maison).} $\rightarrow$ Dire \zh{在我们国家} (zài wǒmen guójiā) ou \zh{在喀麦隆}.
\item \textbf{Lieu après le verbe :} \zh{*人民参加选举在喀麦隆} $\rightarrow$ \textbf{Faux}. Correct : \zh{在喀麦隆，人民参加选举}.
\item \textbf{Confusion \zh{但是} / \zh{而且} / \zh{其次}} : \zh{但是} = \emph{mais} (opposition/nuance). Pour ajouter un argument : \zh{而且} (de plus) ou \zh{其次} (ensuite).
\item \textbf{Tons :} \zh{民} (mín, 2\textsuperscript{e} ton montant) $\neq$ \zh{明} (míng, 2\textsuperscript{e} ton montant mais voyelle \emph{ing}). \zh{主} (zhǔ, 3\textsuperscript{e} ton bas-descendant-remontant) $\neq$ \zh{住} (zhù, 4\textsuperscript{e} ton descendant).
\item \textbf{Classificateur oublié :} \zh{一封信} (yì fēng xìn), \zh{一个看法} (yí ge kànfǎ), \zh{一些年轻人} (yìxiē niánqīngrén).
\end{itemize}
\end{attention}

\courssub{Méthode : Rédiger une lettre argumentative en 5 étapes}

\begin{methode}[De la consigne à la version finale]
\begin{enumerate}
\item \textbf{Analyser la commande} : Identifie le destinataire (李明), le format (lettre), les 3 questions à traiter (définition, élections, avis jeunes), la longueur (6--8 phrases).
\item \textbf{Planifier sur brouillon} : Remplis le squelette : 
\begin{itemize}
    \item Thèse : \zh{我认为，民主就是……}
    \item Arg. 1 : \zh{首先，在喀麦隆……}
    \item Arg. 2 : \zh{其次，我们……}
    \item Nuance : \zh{但是，……}
    \item Conclusion : \zh{所以，我觉得……}
\end{itemize}
\item \textbf{Rédiger en chinois} : Utilise le lexique du tableau, les connecteurs, les structures \zh{在+Lieu+V}, \zh{应该/可以}. Vérifie l'ordre des mots à chaque phrase.
\item \textbf{Soigner l'écriture} : Écris les caractères au propre (traits, radicaux). Relis pour repérer \zh{民/氏}, \zh{选/远}, \zh{权/欢}.
\item \textbf{Relire et corriger} : Grille d'auto-évaluation (caractères, grammaire, lexique, connecteurs, politesse). Échange avec un pair pour correction croisée.
\end{enumerate}
\end{methode}

\courssub{Production modèle commentée : « La démocratie chez moi »}

\begin{exoresolu}[Lettre type à 李明 — 7 phrases]
Rédige la réponse à 李明 en respectant le plan thèse / arguments / nuance / conclusion.
\tcblower
\textbf{Texte chinois (caractères + pinyin) :}\\[4pt]
亲爱的李明：\\
Qīn'ài de Lǐ Míng:\\
\textit{Cher Li Ming,}\\[4pt]
你好！我想跟你谈谈我们国家的民主。\\
Nǐ hǎo! Wǒ xiǎng gēn nǐ tántan wǒmen guójiā de mínzhǔ.\\
\textit{Bonjour ! Je voudrais te parler de la démocratie dans notre pays.}\\[4pt]
我认为，民主就是人民可以选择自己的领导人。\\
Wǒ rènwéi, mínzhǔ jiùshì rénmín kěyǐ xuǎnzé zìjǐ de lǐngdǎorén.\\
\textit{Je pense que la démocratie, c'est le fait que le peuple peut choisir ses propres dirigeants.}\\[4pt]
首先，在喀麦隆，人民可以参加选举，投票选出总统。\\
Shǒuxiān, zài Kāmàilóng, rénmín kěyǐ cānjiā xuǎnjǔ, tóupiào xuǎnchū zǒngtǒng.\\
\textit{D'abord, au Cameroun, le peuple peut participer aux élections et voter pour élire le président.}\\[4pt]
其次，我们有言论自由，每个人都可以说自己的看法。\\
Qícì, wǒmen yǒu yánlùn zìyóu, měi ge rén dōu kěyǐ shuō zìjǐ de kànfǎ.\\
\textit{Ensuite, nous avons la liberté d'expression, chacun peut exprimer son opinion.}\\[4pt]
但是，民主还不完美：有些年轻人对政治不感兴趣。\\
Dànshì, mínzhǔ hái bù wánměi: yǒuxiē niánqīngrén duì zhèngzhì bù gǎn xìngqù.\\
\textit{Mais la démocratie n'est pas encore parfaite : certains jeunes ne s'intéressent pas à la politique.}\\[4pt]
所以，我觉得我们应该多了解民主，多参加国家的生活。\\
Suǒyǐ, wǒ juéde wǒmen yīnggāi duō liǎojiě mínzhǔ, duō cānjiā guójiā de shēnghuó.\\
\textit{C'est pourquoi je pense que nous devrions mieux connaître la démocratie et participer davantage à la vie du pays.}\\[4pt]
祝好！\\
Zhù hǎo!\\
\textit{Amicalement / Meilleurs vœux!}\\[4pt]
你的朋友，恩东\\
Nǐ de péngyou, Ēndōng\\
\textit{Ton ami, Ndong}\\[6pt]
\textbf{Commentaire linguistique détaillé :}
\begin{itemize}
\item \textbf{Formules épistolaires} : \zh{亲爱的李明} (appel) / \zh{祝好！} (clôture standard, équivalent « Amicalement ») / \zh{你的朋友，恩东} (signature).
\item \textbf{Thèse (phrase 3)} : \zh{我认为} + \zh{就是} pour une définition nette. \zh{选择} (choisir) + \zh{自己} (soi-même) + \zh{领导人} (dirigeants) : vocabulaire précis HSK 3-4.
\item \textbf{Argument 1 (phrase 4)} : \zh{首先} + \zh{在喀麦隆} (lieu antéposé) + \zh{参加选举} (collocation) + \zh{投票选出} (série verbale : voter $\rightarrow$ élire). \zh{选出} (xuǎnchū) = complément de résultat « élire avec succès ».
\item \textbf{Argument 2 (phrase 5)} : \zh{其次} + \zh{有言论自由} (avoir la liberté d'expression) + \zh{每个人都可以} (structure \zh{都} pour l'universel) + \zh{说看法} (exprimer son avis).
\item \textbf{Nuance (phrase 6)} : \zh{但是} + \zh{还不完美} (pas encore parfait) + \zh{对……不感兴趣} (structure figée). Le deux-points \zh{：} introduit l'explication de la nuance.
\item \textbf{Conclusion (phrase 7)} : \zh{所以} + \zh{应该} (devoir moral) + \zh{多+V} (faire davantage V). \zh{参加国家的生活} (participer à la vie nationale) : belle image civique.
\item \textbf{Bilan lexical} : 11 mots du glossaire utilisés (\zh{民主、人民、选举、投票、自由、看法、国家、参加、重要(impliqué), 权利(impliqué), 选择}). 4 connecteurs (\zh{首先、其次、但是、所以}).
\end{itemize}
\end{exoresolu}

\courssub{Entraînement guidé : Thème et Version}

\begin{exercice}[Phrases à traduire — Travail individuel 15 min]
\textbf{A. Thème (Français $\rightarrow$ Chinois). Traduis en caractères + pinyin.}
\begin{enumerate}
\item La démocratie est importante pour notre pays.
\item Au Cameroun, les jeunes peuvent voter aux élections.
\item Je pense que le peuple a le droit de choisir.
\item Mais la liberté d'expression n'est pas encore parfaite.
\item Donc, nous devons participer à la vie politique.
\end{enumerate}

\textbf{B. Version (Chinois $\rightarrow$ Français). Traduis en français naturel.}
\begin{enumerate}
\item 我认为民主就是人民有权利。
\item 首先，在喀麦隆，我们可以投票。
\item 其次，年轻人有言论自由。
\item 但是，有些人不参加选举。
\item 所以，我们应该了解政治。
\end{enumerate}
\end{exercice}

\begin{corrige}[Exercice Thème / Version]
\textbf{A. Thème (propositions) :}
\begin{enumerate}
\item \zh{民主对我们国家很重要。} (Mínzhǔ duì wǒmen guójiā hěn zhòngyào.)
\item \zh{在喀麦隆，年轻人可以参加选举投票。} (Zài Kāmàilóng, niánqīngrén kěyǐ cānjiā xuǎnjǔ tóupiào.)
\item \zh{我认为人民有权利选择。} (Wǒ rènwéi rénmín yǒu quánlì xuǎnzé.)
\item \zh{但是，言论自由还不完美。} (Dànshì, yánlùn zìyóu hái bù wánměi.)
\item \zh{所以，我们应该参加政治生活。} (Suǒyǐ, wǒmen yīnggāi cānjiā zhèngzhì shēnghuó.)
\end{enumerate}

\textbf{B. Version :}
\begin{enumerate}
\item Je pense que la démocratie, c'est que le peuple a des droits. / La démocratie signifie que le peuple a des droits.
\item D'abord, au Cameroun, nous pouvons voter.
\item Ensuite, les jeunes ont la liberté d'expression.
\item Mais, certaines personnes ne participent pas aux élections.
\item Donc, nous devrions nous intéresser à la politique / participer à la vie politique.
\end{enumerate}
\end{corrige}

\courssub{Activité finale : Écrire à 李明 (Évaluation formative)}

\begin{experience}[Rédaction, correction par les pairs, version définitive]
\textbf{Consignes :}
\begin{enumerate}
\item \textbf{Rédige} ta lettre sur ton cahier (brouillon $\rightarrow$ propre). Respecte le modèle : appel, 6--8 phrases, thèse/arguments/nuance/conclusion, formules de politesse, signature.
\item \textbf{Échange} ta copie avec ton voisin. Utilise la \textbf{grille ci-dessous} pour corriger. Note le score sur 20 et écris \textbf{2 points forts} + \textbf{2 conseils}.
\item \textbf{Réécris} ta version finale en tenant compte des corrections. Les meilleures lettres seront affichées et numérisées pour l'envoi à Pékin.
\end{enumerate}
\end{experience}

\courssub{Grille d'évaluation par les pairs (20 points)}

\begin{methode}[Grille critériée — À photocopier ou projeter]
\begin{center}
\begin{tabularx}{\linewidth}{|X|c|X|}
\hline
\rowcolor{popblueL} \textbf{Critère} & \textbf{Points} & \textbf{Indicateurs de réussite} \\
\hline
\textbf{Caractères \& Écriture} & 5 & Traits corrects, proportions équilibrées, pas de confusion \zh{民/氏}, \zh{选/远}, \zh{权/欢}. Radicaux visibles. \\
\hline
\textbf{Grammaire \& Syntaxe} & 5 & Ordre SVO respecté. Lieu \zh{在...} avant verbe. Modaux \zh{可以/应该} bien employés. Négation correcte (\zh{不}). \\
\hline
\textbf{Lexique Thématique} & 4 & $\ge$ 6 mots du glossaire utilisés à bon escient. Pas de répétition excessive. \\
\hline
\textbf{Connecteurs \& Organisation} & 4 & 4 connecteurs (\zh{首先、其次、但是、所以}) présents et logiques. Structure thèse/args/nuance/concl. claire. \\
\hline
\textbf{Conventions Épistolaires} & 2 & Appel \zh{亲爱的李明}, clôture \zh{祝好！}, signature \zh{你的朋友, [Prénom]}. \\
\hline
\textbf{TOTAL} & \textbf{20} &  \\
\hline
\end{tabularx}
\end{center}
\textbf{Retour pair :} \quad Points forts : 1) \ldots \quad 2) \ldots \quad Conseils : 1) \ldots \quad 2) \ldots
\end{methode}

\courssub{Bilan et ouverture socioculturelle}

\begin{aretenir}[Ce qu'il faut retenir de la leçon 23]
\begin{itemize}
\item \textbf{Structure type d'un texte argumentatif simple} : \zh{我认为…} (Thèse) $\rightarrow$ \zh{首先… 其次…} (Arguments) $\rightarrow$ \zh{但是…} (Nuance) $\rightarrow$ \zh{所以…} (Conclusion).
\item \textbf{Mots-outils indispensables} : \zh{在} (lieu antéposé), \zh{就是} (définition), \zh{可以/应该} (modalité), \zh{对…感兴趣} (intérêt).
\item \textbf{Caractères à maîtriser} : \zh{民、主、人、选、举、投、票、权、利、自、由、国、看、法、参、加、重、要}.
\item \textbf{Conventions de la lettre} : Appel (\zh{亲爱的+Prénom}), Corps, Formule de politesse (\zh{祝好/祝一切顺利}), Signature (\zh{你的朋友, + Prénom}).
\end{itemize}
\end{aretenir}

\begin{savaistu}[Jeunesse, citoyenneté : regards croisés Cameroun — Chine]
\textbf{Au Cameroun} : Majorité électorale à 20 ans. Vote au scrutin uninominal majoritaire à un tour (présidentielle) ou à deux tours (législatives). Fort engagement associatif des jeunes (clubs UNESCO, conseils des jeunes), mais abstention notable chez les 20--35 ans (désaffection, problèmes logistiques).
\textbf{En Chine} : Pas d'élections nationales au suffrage universel direct pour le président (élu par l'Assemblée populaire nationale). Élections locales directes au niveau des villages et arrondissements (\zh{村委会}, \zh{居委会}) dès 18 ans. Participation encadrée par le Parti communiste ; engagement civique via la Ligue de la jeunesse communiste (\zh{共青团}), bénévolat (\zh{志愿者}), crédit social.
\textbf{Point commun} : Dans les deux pays, l'école (lycée) est le lieu premier d'éducation à la citoyenneté (cours d'éducation civique, simulations d'élections, débats). Les jeunes camerounais et chinois partagent le souci de l'emploi, de l'environnement et de la stabilité.
\textbf{Pour aller plus loin} : Regarde le documentaire \emph{« Please Vote for Me »} (élection de délégués en primaire à Wuhan) et compare avec l'élection des délégués de classe dans ton lycée.
\end{savaistu}

\begin{popfigure}
\begin{tikzpicture}[
  node distance=1.2cm and 1.5cm,
  box/.style={rectangle, draw=popblue, thick, rounded corners, fill=popblueL, text width=3.2cm, align=center, font=\small},
  arrow/.style={-Stealth, thick, popdark}
]
\node[box, align=center] (start){Analyser la lettre\\de M. Wang\\ (3 questions)};
\node[box, below=of start, align=center] (plan){Planifier :\\ Thèse / Arg1 / Arg2 / Nuance / Concl.};
\node[box, below=of plan, align=center] (write){Rédiger en chinois\\ (lexique + connecteurs)};
\node[box, below=of write, align=center] (check){Vérifier :\\ Caractères, Grammaire, Politesse};
\node[box, below=of check, align=center] (peer){Correction croisée\\ (Grille 20 pts)};
\node[box, below=of peer, align=center] (final){Version finale\\ (Affichage / Envoi)};
\draw[arrow] (start) -- (plan);
\draw[arrow] (plan) -- (write);
\draw[arrow] (write) -- (check);
\draw[arrow] (check) -- (peer);
\draw[arrow] (peer) -- (final);
\end{tikzpicture}
\legende{Processus d'écriture de la lettre argumentative — Leçon 23}
\end{popfigure}

\newpage

\coursec{Méthodes et exercices résolus}

\courssub{Savoir-faire 1 : Produire un texte simple sur la démocratie}

\begin{textebox}[Modèle de texte : 民主和我们 (La démocratie et nous)]
\zh{民主对每个国家都很重要。} Mínzhǔ duì měi ge guójiā dōu hěn zhòngyào. « La démocratie est très importante pour chaque pays. »\\
\zh{首先，民主给公民自由和权利。} Shǒuxiān, mínzhǔ gěi gōngmín zìyóu hé quánlì. « D'abord, la démocratie donne aux citoyens la liberté et des droits. »\\
\zh{然后，每个公民都可以投票，选举自己喜欢的人。} Ránhòu, měi ge gōngmín dōu kěyǐ tóupiào, xuǎnjǔ zìjǐ xǐhuan de rén. « Ensuite, chaque citoyen peut voter et élire les personnes qu'il aime. »\\
\zh{但是，有些年轻人觉得投票没有用。} Dànshì, yǒuxiē niánqīngrén juéde tóupiào méiyǒu yòng. « Mais certains jeunes trouvent que voter ne sert à rien. »\\
\zh{我认为这个想法不对，因为投票是我们的权利，也是我们的责任。} Wǒ rènwéi zhè ge xiǎngfa bú duì, yīnwèi tóupiào shì wǒmen de quánlì, yě shì wǒmen de zérèn. « Je pense que cette idée est fausse, parce que voter est notre droit et aussi notre responsabilité. »\\
\zh{总之，没有公民的参与，就没有真正的民主。} Zǒngzhī, méiyǒu gōngmín de cānyù, jiù méiyǒu zhēnzhèng de mínzhǔ. « En somme, sans la participation des citoyens, il n'y a pas de vraie démocratie. »
\end{textebox}

\textbf{Lecture commentée du modèle}\par
Ce texte suit le plan en trois parties attendu à l'examen : une \textbf{introduction} (phrase 1 : définition ou importance du sujet), un \textbf{développement} (phrases 2 à 4 : arguments et exemples, avec les connecteurs \zh{首先}, \zh{然后}, \zh{但是}), et une \textbf{conclusion} (phrases 5 et 6 : avis personnel avec \zh{我认为} et bilan avec \zh{总之}). Chaque phrase ne contient qu'une seule idée, avec la structure Sujet + Verbe + Complément.

\begin{methode}[Rédiger un texte simple sur la démocratie]
Pour produire un texte court et correct sur le thème de la démocratie, suivez cette démarche en cinq étapes :
\begin{enumerate}
\item \textbf{Comprendre la consigne} : repérez le sujet (parler de la démocratie, donner son avis, décrire une situation), la longueur demandée (généralement 5 à 8 phrases) et le destinataire.
\item \textbf{Lister le vocabulaire utile} avant d'écrire : \zh{民主} (mínzhǔ, démocratie), \zh{选举} (xuǎnjǔ, élection), \zh{投票} (tóupiào, voter), \zh{权利} (quánlì, droit), \zh{自由} (zìyóu, liberté), \zh{公民} (gōngmín, citoyen), \zh{国家} (guójiā, pays).
\item \textbf{Construire un plan en trois parties} : introduction (qu'est-ce que la démocratie ?), développement (exemples : les élections, les droits des citoyens), conclusion (mon avis avec \zh{我认为}).
\item \textbf{Rédiger avec des phrases simples} : Sujet + Verbe + Complément. Une idée par phrase. Utilisez les connecteurs : \zh{首先} (d'abord), \zh{然后} (ensuite), \zh{但是} (mais), \zh{所以} (donc), \zh{因为} (parce que).
\item \textbf{Relire} : vérifiez les tons du pinyin, l'ordre des mots (le complément de temps et de lieu se place avant le verbe), et l'emploi des classificateurs.
\end{enumerate}
Structures d'opinion à retenir : \zh{我认为} + phrase (wǒ rènwéi, je pense que) ; \zh{每个人都应该} + verbe (měi ge rén dōu yīnggāi, chacun devrait) ; \zh{我觉得} + phrase (wǒ juéde, je trouve que).
\end{methode}

\begin{exoresolu}[Production : un court texte sur le vote]
Énoncé : Écrivez en chinois un texte de 5 à 6 phrases intitulé « Le vote est important » (\zh{投票很重要}). Utilisez au moins deux connecteurs et le mot \zh{权利}.
\tcblower
Solution détaillée :
\begin{enumerate}
\item Analyse de la consigne : il faut donner une opinion (\zh{投票很重要}), donc la structure \zh{我认为} est appropriée. Le mot \zh{权利} doit apparaître : on peut écrire \zh{投票是我们的权利} (voter est notre droit).
\item Plan : phrase 1 (introduction : la démocratie), phrases 2-4 (le vote, droit des citoyens, exemple), phrase 5 (conclusion avec avis personnel).
\item Rédaction :\\
\zh{每个国家都需要民主。} Měi ge guójiā dōu xūyào mínzhǔ. « Chaque pays a besoin de la démocratie. »\\
\zh{首先，投票是每个公民的权利。} Shǒuxiān, tóupiào shì měi ge gōngmín de quánlì. « D'abord, voter est le droit de chaque citoyen. »\\
\zh{然后，我们可以选举我们喜欢的人。} Ránhòu, wǒmen kěyǐ xuǎnjǔ wǒmen xǐhuan de rén. « Ensuite, nous pouvons élire les personnes que nous aimons. »\\
\zh{但是，有些年轻人不去投票。} Dànshì, yǒuxiē niánqīngrén bú qù tóupiào. « Mais certains jeunes ne vont pas voter. »\\
\zh{我认为投票很重要，因为这关系到我们国家的未来。} Wǒ rènwéi tóupiào hěn zhòngyào, yīnwèi zhè guānxì dào wǒmen guójiā de wèilái. « Je pense que voter est très important, parce que cela concerne l'avenir de notre pays. »
\item Vérification : \zh{每个} + nom sans classificateur supplémentaire ; \zh{因为} introduit bien la cause ; les connecteurs \zh{首先} et \zh{然后} sont en tête de phrase, suivis de la virgule chinoise （，）.
\end{enumerate}
Conclusion : le texte compte 5 phrases, respecte le plan, emploie deux connecteurs et le mot \zh{权利} : la consigne est entièrement remplie.
\end{exoresolu}

\courssub{Savoir-faire 2 : Lexique et écriture des caractères du thème}

\begin{center}
\begin{tabularx}{\linewidth}{|X|X|X|X|}
\hline
\rowcolor{popblueL} Caractères & Pinyin & Nature & Traduction \\
\hline
\zh{民主} & mínzhǔ & nom & démocratie \\
\hline
\zh{选举} & xuǎnjǔ & nom / verbe & élection ; élire \\
\hline
\zh{投票} & tóupiào & verbe & voter \\
\hline
\zh{权利} & quánlì & nom & droit \\
\hline
\zh{自由} & zìyóu & nom / adjectif & liberté ; libre \\
\hline
\zh{公民} & gōngmín & nom & citoyen \\
\hline
\zh{国家} & guójiā & nom & pays, État \\
\hline
\zh{法律} & fǎlǜ & nom & loi \\
\hline
\zh{责任} & zérèn & nom & responsabilité \\
\hline
\zh{参与} & cānyù & verbe / nom & participer ; participation \\
\hline
\zh{重要} & zhòngyào & adjectif & important \\
\hline
\zh{未来} & wèilái & nom & avenir \\
\hline
\end{tabularx}
\end{center}

\begin{methode}[Maîtriser les caractères du thème « démocratie »]
Pour écrire sans faute les caractères du thème, appliquez ces réflexes :
\begin{enumerate}
\item \textbf{Identifier le radical (clé)} de chaque caractère : il donne le sens général. Exemples : \zh{民} (mín, peuple) est lui-même une clé ; \zh{权} (quán, droit) a la clé du bois \zh{木} à gauche ; \zh{法} (fǎ, loi) a la clé de l'eau \zh{氵} à gauche ; \zh{选} (xuǎn, choisir) a la clé du mouvement \zh{辶}.
\item \textbf{Respecter l'ordre des traits} : de haut en bas, de gauche à droite, l'horizontal avant le vertical, l'extérieur avant l'intérieur, et on ferme la boîte en dernier (ex. \zh{国} : on écrit le contour, puis \zh{玉} à l'intérieur, puis le trait horizontal du bas qui ferme).
\item \textbf{Distinguer les caractères proches} : \zh{民} (mín, peuple) et \zh{氏} (shì, nom de clan) ; \zh{主} (zhǔ, maître, principal) et \zh{王} (wáng, roi) — \zh{主} a un point en plus au-dessus ; \zh{投} (tóu, lancer, voter) et \zh{没} (méi, ne pas avoir) — même partie droite, clés différentes (la main \zh{扌} contre l'eau \zh{氵}).
\item \textbf{Mémoriser les mots complets}, pas les caractères isolés : \zh{民主} (démocratie), \zh{选举} (élection), \zh{投票} (voter), \zh{自由} (liberté), \zh{法律} (loi).
\item \textbf{S'entraîner par dictée} : écrire chaque mot cinq fois en prononçant le pinyin avec le ton à voix haute.
\end{enumerate}
\end{methode}

\begin{exoresolu}[Dictée et correction de caractères]
Énoncé : a) Écrivez en caractères les mots suivants : mínzhǔ, xuǎnjǔ, tóupiào, quánlì. b) Le caractère \zh{主} est-il correct dans le mot « démocratie » ? Justifiez. c) Combien de traits compte le caractère \zh{民} ?
\tcblower
Solution détaillée :
\begin{enumerate}
\item a) mínzhǔ = \zh{民主} ; xuǎnjǔ = \zh{选举} ; tóupiào = \zh{投票} ; quánlì = \zh{权利}. Vérifiez : \zh{选} s'écrit avec \zh{先} à l'intérieur et la clé du mouvement \zh{辶} tracée en dernier ; \zh{票} a la composante \zh{示} en bas, sans oublier les deux petits traits finaux.
\item b) Oui, \zh{主} est correct : \zh{民主} signifie littéralement « le peuple (\zh{民}) est maître (\zh{主}) ». Attention à ne pas écrire \zh{王} (wáng, roi) : \zh{主} possède un point supplémentaire au-dessus. Écrire \zh{民王} serait une faute grave, car ce mot n'existe pas.
\item c) \zh{民} compte 5 traits : d'abord le pli (horizontal puis vertical vers le bas), puis le petit trait horizontal, ensuite le relèvement vertical, puis le trait horizontal, et enfin le grand crochet oblique qui descend vers la droite. Ordre général : de haut en bas.
\end{enumerate}
Conclusion : la maîtrise des radicaux (\zh{木} dans \zh{权}, \zh{氵} dans \zh{法}, \zh{辶} dans \zh{选}) et de l'ordre des traits permet d'éviter les confusions entre caractères proches, faute la plus sanctionnée à l'écrit.
\end{exoresolu}

\courssub{Savoir-faire 3 : Structurer le texte avec les connecteurs}

\begin{methode}[Utiliser les connecteurs logiques]
Un texte argumentatif simple repose sur des connecteurs bien placés :
\begin{enumerate}
\item \textbf{Énumération} : \zh{首先} (shǒuxiān, d'abord), \zh{然后} (ránhòu, ensuite), \zh{最后} (zuìhòu, enfin). Ils se placent en tête de phrase, suivis d'une virgule.
\item \textbf{Cause et conséquence} : \zh{因为……所以……} (yīnwèi... suǒyǐ..., parce que... donc...). En chinois, on peut utiliser les deux ensemble, contrairement au français : \zh{因为民主很重要，所以每个人都应该投票。}
\item \textbf{Opposition} : \zh{但是} (dànshì, mais) se place en tête de la seconde proposition, après la virgule ; \zh{可是} (kěshì) est plus oral.
\item \textbf{Condition} : \zh{如果……就……} (rúguǒ... jiù..., si... alors...) : \zh{如果公民不投票，就没有民主。}
\item \textbf{Opinion et conclusion} : \zh{我认为} (je pense que), \zh{我觉得} (je trouve que), \zh{总之} (zǒngzhī, en somme).
\end{enumerate}
Réflexe : un seul connecteur en tête de phrase ; ne jamais placer \zh{但是} au milieu de la phrase comme le « mais » français.
\end{methode}

\begin{center}
\begin{tabularx}{\linewidth}{|X|X|X|}
\hline
\rowcolor{popblueL} Structure & Exemple (caractères + pinyin) & Traduction \\
\hline
\zh{首先}……\zh{然后}…… & \zh{首先，投票是权利；然后，投票是责任。} Shǒuxiān, tóupiào shì quánlì ; ránhòu, tóupiào shì zérèn. & D'abord, voter est un droit ; ensuite, voter est une responsabilité. \\
\hline
\zh{因为}……\zh{所以}…… & \zh{因为民主很重要，所以我们都应该投票。} Yīnwèi mínzhǔ hěn zhòngyào, suǒyǐ wǒmen dōu yīnggāi tóupiào. & Parce que la démocratie est importante, nous devrions tous voter. \\
\hline
\zh{如果}……\zh{就}…… & \zh{如果公民不投票，就没有民主。} Rúguǒ gōngmín bù tóupiào, jiù méiyǒu mínzhǔ. & Si les citoyens ne votent pas, il n'y a pas de démocratie. \\
\hline
\end{tabularx}
\end{center}

\begin{exoresolu}[Transformation : relier des phrases]
Énoncé : Reliez les paires de phrases avec le connecteur indiqué. a) \zh{民主很重要。每个人都应该投票。} (\zh{因为……所以……}) b) \zh{投票是权利。有些年轻人不投票。} (\zh{但是}) c) \zh{公民不投票。没有民主。} (\zh{如果……就……})
\tcblower
Solution détaillée :
\begin{enumerate}
\item a) La cause est « la démocratie est importante », la conséquence est « chacun doit voter ». On obtient : \zh{因为民主很重要，所以每个人都应该投票。} Yīnwèi mínzhǔ hěn zhòngyào, suǒyǐ měi ge rén dōu yīnggāi tóupiào. « Parce que la démocratie est importante, chacun devrait voter. » En chinois, \zh{因为} et \zh{所以} peuvent coexister dans la même phrase : ce n'est pas une faute, contrairement au français.
\item b) L'opposition porte sur le contraste droit / comportement : \zh{投票是权利，但是有些年轻人不投票。} Tóupiào shì quánlì, dànshì yǒuxiē niánqīngrén bù tóupiào. « Voter est un droit, mais certains jeunes ne votent pas. » \zh{但是} se place après la virgule, en tête de la seconde proposition.
\item c) La condition : \zh{如果公民不投票，就没有民主。} Rúguǒ gōngmín bù tóupiào, jiù méiyǒu mínzhǔ. « Si les citoyens ne votent pas, il n'y a pas de démocratie. » \zh{就} marque la conséquence et se place juste avant le verbe de la seconde proposition.
\end{enumerate}
Conclusion : les paires \zh{因为……所以……} et \zh{如果……就……} fonctionnent en binôme ; supprimer le second élément rend la phrase bancale.
\end{exoresolu}

\courssub{Savoir-faire 4 : Thème (français → chinois)}

\begin{methode}[Traduire du français vers le chinois]
\begin{enumerate}
\item \textbf{Repérer le squelette} : sujet, verbe, complément d'objet. Traduisez d'abord ce noyau.
\item \textbf{Placer les compléments} : en chinois, le temps et le lieu se mettent avant le verbe (Sujet + Temps + Lieu + Verbe). « Les citoyens votent à Yaoundé demain » → \zh{公民明天在雅温得投票。} Gōngmín míngtiān zài Yǎwēndé tóupiào.
\item \textbf{Choisir les mots-outils} : \zh{的} pour la possession et la qualification (\zh{我们的国家}, notre pays), \zh{了} pour l'action accomplie, \zh{应该} pour « devoir ».
\item \textbf{Classificateurs} : « chaque citoyen » = \zh{每个公民} ; « un pays » = \zh{一个国家} ; « une élection » = \zh{一次选举}.
\item \textbf{Vérifier les tons} en relisant le pinyin : mínzhǔ (2-3), tóupiào (2-4), zìyóu (4-2).
\end{enumerate}
\end{methode}

\begin{exoresolu}[Thème guidé]
Énoncé : Traduisez en chinois : a) « La démocratie est très importante. » b) « Chaque citoyen a le droit de voter. » c) « Hier, nous avons parlé des élections en classe. »
\tcblower
Solution détaillée :
\begin{enumerate}
\item a) Noyau : démocratie + être + importante. Devant un adjectif, le verbe « être » ne se traduit pas par \zh{是} : on utilise \zh{很}. Réponse : \zh{民主很重要。} Mínzhǔ hěn zhòngyào. Ne dites jamais \zh{民主是很重要} : devant un adjectif prédicat, \zh{是} est superflu.
\item b) « Avoir le droit de » = \zh{有……的权利} : \zh{每个公民都有投票的权利。} Měi ge gōngmín dōu yǒu tóupiào de quánlì. Remarquez \zh{都} qui accompagne \zh{每} (chaque... tous) et \zh{的} qui relie le verbe \zh{投票} au nom \zh{权利}.
\item c) Le temps \zh{昨天} se place en tête ou après le sujet, avant le verbe ; le lieu \zh{在班上} se place avant le verbe ; \zh{了} marque l'action accomplie : \zh{昨天我们在班上讨论了选举。} Zuótiān wǒmen zài bān shang tǎolùnle xuǎnjǔ. « Hier, nous avons discuté des élections en classe. » Ordre : Temps + Sujet + Lieu + Verbe + Complément.
\end{enumerate}
Conclusion : le piège principal du thème est l'ordre des mots : toujours placer temps et lieu avant le verbe, et ne jamais calquer la structure française.
\end{exoresolu}

\courssub{Savoir-faire 5 : Version (chinois → français)}

\begin{methode}[Traduire du chinois vers le français]
\begin{enumerate}
\item \textbf{Lire la phrase entière} avant de traduire : repérez le sujet, le verbe principal et les compléments.
\item \textbf{Traduire les blocs} dans l'ordre français : le chinois met temps et lieu avant le verbe, le français les met souvent après. \zh{我们明天在学校开会。} → « Nous aurons une réunion à l'école demain. »
\item \textbf{Rendre les mots-outils} : \zh{的} = possession ou « de » ; \zh{了} = passé composé ; \zh{都} = « tous » ; \zh{应该} = « devoir ».
\item \textbf{Restituer les nuances des connecteurs} : \zh{因为……所以……} = « parce que » (un seul mot en français), \zh{如果……就……} = « si... alors ».
\item \textbf{Relire en français} : la phrase doit être naturelle et grammaticale, pas un mot-à-mot.
\end{enumerate}
\end{methode}

\begin{exoresolu}[Version guidée]
Énoncé : Traduisez en français : a) \zh{民主给每个公民自由。} b) \zh{因为选举很重要，所以年轻人都应该去投票。} c) \zh{如果没有法律，国家就没有民主。}
\tcblower
Solution détaillée :
\begin{enumerate}
\item a) Sujet \zh{民主} (la démocratie), verbe \zh{给} (donner), compléments \zh{每个公民} (à chaque citoyen) et \zh{自由} (la liberté). Traduction : « La démocratie donne la liberté à chaque citoyen. » Structure chinoise : donner + personne + chose ; en français : donner + chose + à + personne.
\item b) \zh{因为……所以……} se rend par un seul « parce que » en français : « Parce que les élections sont importantes, les jeunes devraient tous aller voter. » \zh{都} renforce « tous » ; \zh{应该} = « devraient » ; \zh{去投票} = « aller voter » (verbe de déplacement + but).
\item c) \zh{如果……就……} = « si... (alors) » ; \zh{没有法律} = « il n'y a pas de loi » ; \zh{国家就没有民主} = « le pays n'a pas de démocratie ». Traduction : « S'il n'y a pas de loi, le pays n'a pas de démocratie. » On peut aussi écrire : « Sans loi, pas de démocratie dans un pays. »
\end{enumerate}
Conclusion : en version, le mot-à-mot est interdit : il faut réorganiser les compléments selon la syntaxe française et ne garder qu'un seul connecteur de cause.
\end{exoresolu}

\courssub{Savoir-faire 6 : Production guidée et grille d'évaluation}

\begin{experience}[Atelier d'écriture : « Les jeunes et le vote au Cameroun »]
Par groupes de trois, rédigez un texte de 6 à 8 phrases sur le sujet : \zh{年轻人应该投票吗？} (Niánqīngrén yīnggāi tóupiào ma ?, Les jeunes devraient-ils voter ?).
\begin{enumerate}
\item Répartissez les rôles : un élève propose le plan, un élève rédige les phrases, un élève vérifie les caractères et le pinyin.
\item Imposez-vous d'utiliser : au moins trois connecteurs (\zh{首先}, \zh{但是}, \zh{所以}...), la structure \zh{我认为}, et quatre mots du tableau de vocabulaire.
\item Échangez votre texte avec un autre groupe et corrigez-le à l'aide de la grille ci-dessous.
\item Un représentant de chaque groupe lit le texte final à voix haute ; la classe repère les connecteurs employés.
\end{enumerate}
\end{experience}

\begin{aretenir}[Grille d'évaluation d'un texte simple]
\begin{enumerate}
\item \textbf{Respect de la consigne} : sujet traité, longueur respectée (5 à 8 phrases), avis personnel exprimé.
\item \textbf{Structure} : introduction, développement, conclusion ; au moins deux connecteurs bien placés.
\item \textbf{Grammaire} : ordre des mots (temps et lieu avant le verbe), pas de \zh{是} devant un adjectif, mots-outils \zh{的}, \zh{了}, \zh{都} correctement employés.
\item \textbf{Caractères} : aucun caractère proche confondu (\zh{主}/\zh{王}, \zh{民}/\zh{氏}, \zh{投}/\zh{没}), traits complets.
\item \textbf{Vocabulaire} : au moins quatre mots du thème « démocratie » employés à bon escient.
\end{enumerate}
Barème indicatif sur 20 : consigne 4, structure 4, grammaire 6, caractères 3, vocabulaire 3.
\end{aretenir}

\begin{attention}[Les 5 erreurs qui coûtent des points au Bac]
\begin{enumerate}
\item \textbf{Confondre les caractères proches} : \zh{主} (zhǔ, maître) et \zh{王} (wáng, roi) ; \zh{民} (mín, peuple) et \zh{氏} (shì) ; \zh{投} (tóu, voter) et \zh{没} (méi). Un point ou une clé oubliée change le mot : \zh{民王} n'existe pas.
\item \textbf{Calquer l'ordre français} : « Nous votons demain à l'école » ne se traduit pas mot à mot. En chinois : Sujet + Temps + Lieu + Verbe → \zh{我们明天在学校投票。} Le temps et le lieu passent AVANT le verbe.
\item \textbf{Mettre \zh{是} devant un adjectif} : « la démocratie est importante » = \zh{民主很重要}, jamais \zh{民主是很重要}. \zh{是} ne s'emploie qu'entre deux noms (\zh{我是学生}, je suis élève).
\item \textbf{Oublier les classificateurs et \zh{的}} : « un pays » = \zh{一个国家} (pas \zh{一国}) ; « le droit de voter » = \zh{投票的权利}, avec \zh{的} obligatoire entre le verbe et le nom.
\item \textbf{Faux amis et connecteurs} : ne placez pas \zh{但是} au milieu de la phrase comme le « mais » français ; rappelez-vous que \zh{因为……所以……} s'utilisent ensemble en chinois, alors qu'en français un seul suffit. En version, n'écrivez jamais « parce que... donc... ».
\end{enumerate}
\end{attention}

\newpage

\coursec{Exercices}

\courssub{Je vérifie mes connaissances}

\begin{exercice}[QCM : vocabulaire de la démocratie]
Pour chaque question, entoure la bonne réponse.

\begin{enumerate}
\item \zh{民主} (mínzhǔ) signifie :
\begin{enumerate}
\item la liberté \item la démocratie \item la responsabilité
\end{enumerate}
\item « le droit de vote » se dit :
\begin{enumerate}
\item \zh{权利} \item \zh{自由} \item \zh{投票权}
\end{enumerate}
\item \zh{选举} (xuǎnjǔ) signifie :
\begin{enumerate}
\item l'élection \item le gouvernement \item le citoyen
\end{enumerate}
\item « la liberté d'expression » se dit :
\begin{enumerate}
\item \zh{言论自由} \item \zh{互相尊重} \item \zh{参加选举}
\end{enumerate}
\end{enumerate}
\end{exercice}

\begin{exercice}[Vrai ou faux — justifie ta réponse]
Dis si chaque phrase est vraie ou fausse, puis justifie en une phrase en français.

\begin{enumerate}
\item \zh{公民} (gōngmín) signifie « citoyen ».
\item Dans la phrase \zh{我们应该互相尊重。}, le mot \zh{互相} signifie « toujours ».
\item \zh{责任} (zérèn) signifie « responsabilité, devoir ».
\item En chinois, le complément de cause introduit par \zh{因为} se place toujours après la phrase principale.
\end{enumerate}
\end{exercice}

\begin{exercice}[Vocabulaire : complète le tableau]
Complète le tableau en écrivant les caractères chinois correspondant à chaque mot français (aide-toi de la leçon).

\begin{center}
\begin{tabularx}{\linewidth}{|X|X|X|}\hline
\rowcolor{popblueL} Français & Pinyin & Caractères \\ \hline
la démocratie & mínzhǔ & \\ \hline
le droit & quánlì & \\ \hline
voter & tóupiào & \\ \hline
l'élection & xuǎnjǔ & \\ \hline
la responsabilité & zérèn & \\ \hline
respecter & zūnzhòng & \\ \hline
\end{tabularx}
\end{center}
\end{exercice}

\begin{exercice}[Associe les caractères au pinyin]
Relie chaque mot en caractères à son pinyin.

\begin{enumerate}
\item \zh{民主} \hspace{2cm} a. quánlì
\item \zh{权利} \hspace{2cm} b. xuǎnjǔ
\item \zh{投票} \hspace{2cm} c. mínzhǔ
\item \zh{选举} \hspace{2cm} d. zérèn
\item \zh{责任} \hspace{2cm} e. tóupiào
\end{enumerate}
\end{exercice}

\courssub{Je m'entraîne}

\begin{exercice}[Traduis en français]
Traduis les phrases suivantes en français.

\begin{enumerate}
\item \zh{公民有投票的权利。} (Gōngmín yǒu tóupiào de quánlì.)
\item \zh{我们应该参加选举。} (Wǒmen yīnggāi cānjiā xuǎnjǔ.)
\item \zh{民主国家里，人民有言论自由。} (Mínzhǔ guójiā li, rénmín yǒu yánlùn zìyóu.)
\item \zh{年轻人也有责任。} (Niánqīngrén yě yǒu zérèn.)
\end{enumerate}
\end{exercice}

\begin{exercice}[Texte à trous : les connecteurs]
Complète le texte suivant avec les connecteurs : \zh{首先}、\zh{其次}、\zh{但是}、\zh{所以}、\zh{也} (shǒuxiān, qícì, dànshì, suǒyǐ, yě).

\begin{textebox}[Texte à compléter]
民主很重要。\_\_\_\_，人民可以选择领导人。\_\_\_\_，公民有自由。\_\_\_\_，民主也需要责任。\_\_\_\_，我们应该学习民主。\\
Mínzhǔ hěn zhòngyào. \_\_\_\_, rénmín kěyǐ xuǎnzé lǐngdǎorén. \_\_\_\_, gōngmín yǒu zìyóu. \_\_\_\_, mínzhǔ yě xūyào zérèn. \_\_\_\_, wǒmen yīnggāi xuéxí mínzhǔ.
\end{textebox}
\end{exercice}

\begin{exercice}[Remets les mots dans l'ordre]
Remets les mots dans le bon ordre pour former des phrases correctes, puis traduis-les en français.

\begin{enumerate}
\item \zh{有 / 公民 / 言论自由 / 。}
\item \zh{应该 / 我们 / 尊重 / 互相 / 。}
\item \zh{参加 / 年轻人 / 选举 / 应该 / 。}
\item \zh{因为 / 重要 / 很 / 民主 / ，所以 / 我们 / 学习 / 要 / 。}
\end{enumerate}
\end{exercice}

\begin{exercice}[Discrimination des caractères]
Entoure le bon caractère dans chaque paire, puis recopie le mot complet trois fois en respectant l'ordre des traits (haut → bas, gauche → droite). Ensuite, dicte ces mots à un camarade.

\begin{enumerate}
\item 民（主 / 王）
\item （权 / 劝）利
\item 投（票 / 飘）
\item （选 / 远）举
\item 责（任 / 住）
\end{enumerate}
\end{exercice}

\courssub{Je me prépare au Bac}

\begin{exercice}[Compréhension de texte et questions de langue]
Lis le texte suivant, puis réponds aux questions en français.

\begin{textebox}[La démocratie et les citoyens]
民主国家里，公民有言论自由，也有责任。我们应该互相尊重，互相学习。每年，人民可以参加选举，选择自己的领导人。年轻人也应该学习民主，因为民主对大家很重要。\\
Mínzhǔ guójiā li, gōngmín yǒu yánlùn zìyóu, yě yǒu zérèn. Wǒmen yīnggāi hùxiāng zūnzhòng, hùxiāng xuéxí. Měi nián, rénmín kěyǐ cānjiā xuǎnjǔ, xuǎnzé zìjǐ de lǐngdǎorén. Niánqīngrén yě yīnggāi xuéxí mínzhǔ, yīnwèi mínzhǔ duì dàjiā hěn zhòngyào.\\
« Dans un pays démocratique, les citoyens ont la liberté d'expression et aussi des responsabilités. Nous devons nous respecter mutuellement et apprendre les uns des autres. Chaque année, le peuple peut participer aux élections et choisir ses dirigeants. Les jeunes doivent aussi apprendre la démocratie, car la démocratie est très importante pour tous. »
\end{textebox}

\textbf{Questions de compréhension}
\begin{enumerate}
\item Quels sont les deux éléments que possèdent les citoyens dans un pays démocratique ?
\item Que peuvent faire les citoyens chaque année ?
\item Pourquoi les jeunes doivent-ils apprendre la démocratie ?
\end{enumerate}

\textbf{Questions de langue}
\begin{enumerate}
\item Releve dans le texte deux connecteurs et explique leur rôle.
\item Dans la phrase \zh{公民有言论自由，也有责任。}, quel mot exprime l'addition ? Construis une phrase avec ce mot.
\item Traduis en français : \zh{我们应该互相尊重，互相学习。}
\end{enumerate}
\end{exercice}

\begin{exercice}[Thème et version type Bac]
\textbf{A. Thème} — Traduis les phrases suivantes en chinois (caractères obligatoires, pinyin apprécié).

\begin{enumerate}
\item Les citoyens ont le droit de vote et la liberté d'expression.
\item Parce que la démocratie est importante, nous devons participer aux élections.
\item Les jeunes Camerounais doivent aussi apprendre la démocratie.
\end{enumerate}

\textbf{B. Version} — Traduis le texte suivant en français.

\begin{textebox}[Texte à traduire]
民主国家里，公民有言论自由，也有责任。我们应该互相尊重，互相学习。\\
Mínzhǔ guójiā li, gōngmín yǒu yánlùn zìyóu, yě yǒu zérèn. Wǒmen yīnggāi hùxiāng zūnzhòng, hùxiāng xuéxí.
\end{textebox}
\end{exercice}

\begin{exercice}[Expression écrite type Bac — notée sur 20]
\textbf{Sujet :} « La démocratie et les jeunes Camerounais » (\zh{民主和喀麦隆青年}).

Rédige un texte de \textbf{5 à 8 phrases} en chinois (entre 60 et 100 caractères). Tu dois :

\begin{enumerate}
\item dire pourquoi la démocratie est importante ;
\item parler des droits et des responsabilités des jeunes ;
\item utiliser au moins \textbf{trois connecteurs} vus dans la leçon (\zh{首先}、\zh{其次}、\zh{但是}、\zh{所以}、\zh{因为}…) ;
\item utiliser au moins \textbf{cinq mots du lexique} du thème (\zh{民主}、\zh{公民}、\zh{权利}、\zh{投票}、\zh{选举}、\zh{责任}、\zh{自由}、\zh{尊重}…) ;
\item écrire les caractères correctement, en respectant l'ordre des traits.
\end{enumerate}

\textbf{Grille d'évaluation :} respect du sujet et du nombre de phrases (4 pts) ; exactitude des structures grammaticales (6 pts) ; richesse et exactitude du vocabulaire du thème (4 pts) ; utilisation des connecteurs (3 pts) ; qualité de l'écriture des caractères (3 pts).
\end{exercice}

\newpage

\coursec{Corrigés des exercices}

\begin{corrige}[Exercice 1 — QCM : vocabulaire de la démocratie]
\begin{enumerate}
\item \textbf{b.} la démocratie. \zh{民主} (mínzhǔ) = \zh{民} « peuple » + \zh{主} « maître » : « le peuple est maître ».
\item \textbf{c.} \zh{投票权} (tóupiàoquán). \zh{权利} signifie « droit » en général, \zh{自由} signifie « liberté ».
\item \textbf{a.} l'élection. \zh{选举} (xuǎnjǔ) = \zh{选} « choisir » + \zh{举} « lever ».
\item \textbf{a.} \zh{言论自由} (yánlùn zìyóu) = \zh{言论} « parole, discours » + \zh{自由} « liberté ».
\end{enumerate}
\end{corrige}

\begin{corrige}[Exercice 2 — Vrai ou faux]
\begin{enumerate}
\item \textbf{Vrai.} \zh{公民} (gōngmín) signifie bien « citoyen » : \zh{公} « public » + \zh{民} « peuple ».
\item \textbf{Faux.} \zh{互相} (hùxiāng) signifie « mutuellement, l'un l'autre » ; « toujours » se dit \zh{总是} (zǒngshì). La phrase signifie : « Nous devons nous respecter mutuellement. »
\item \textbf{Vrai.} \zh{责任} (zérèn) signifie « responsabilité, devoir ».
\item \textbf{Faux.} La subordonnée de cause avec \zh{因为} se place normalement \emph{avant} la phrase principale : \zh{因为民主很重要，所以我要学习。} « Parce que la démocratie est importante, je veux l'apprendre. » C'est l'inverse de l'ordre le plus fréquent en français.
\end{enumerate}
\end{corrige}

\begin{corrige}[Exercice 3 — Vocabulaire : complète le tableau]
\begin{center}
\begin{tabularx}{\linewidth}{|X|X|X|X|}\hline
\rowcolor{popblueL} Caractères & Pinyin & Nature & Traduction \\ \hline
\zh{民主} & mínzhǔ & n. & la démocratie \\ \hline
\zh{权利} & quánlì & n. & le droit \\ \hline
\zh{投票} & tóupiào & v. & voter \\ \hline
\zh{选举} & xuǎnjǔ & n. & l'élection \\ \hline
\zh{责任} & zérèn & n. & la responsabilité \\ \hline
\zh{尊重} & zūnzhòng & v. & respecter \\ \hline
\end{tabularx}
\end{center}
\textbf{Points de vigilance :} ne pas confondre \zh{主} (zhǔ) et \zh{王} (wáng) — \zh{主} a un point en haut ; \zh{权} (clé 木 « bois ») ne s'écrit pas \zh{劝} ; \zh{票} et non \zh{飘}.
\end{corrige}

\begin{corrige}[Exercice 4 — Associe les caractères au pinyin]
\begin{enumerate}
\item \zh{民主} → \textbf{c.} mínzhǔ
\item \zh{权利} → \textbf{a.} quánlì
\item \zh{投票} → \textbf{e.} tóupiào
\item \zh{选举} → \textbf{b.} xuǎnjǔ
\item \zh{责任} → \textbf{d.} zérèn
\end{enumerate}
\end{corrige}

\begin{corrige}[Exercice 5 — Traduis en français]
\begin{enumerate}
\item \zh{公民有投票的权利。} — « Les citoyens ont le droit de vote. » Structure : Sujet + \zh{有} + complément ; \zh{投票的} est un complément du nom \zh{权利} relié par \zh{的}.
\item \zh{我们应该参加选举。} — « Nous devons participer aux élections. » Le verbe modal \zh{应该} « devoir » se place avant le verbe \zh{参加}.
\item \zh{民主国家里，人民有言论自由。} — « Dans un pays démocratique, le peuple a la liberté d'expression. » \zh{里} marque le lieu : « à l'intérieur de, dans ».
\item \zh{年轻人也有责任。} — « Les jeunes aussi ont des responsabilités. » Attention : \zh{也} se place \emph{avant} le verbe, jamais en fin de phrase comme « aussi » en français.
\end{enumerate}
\end{corrige}

\begin{corrige}[Exercice 6 — Texte à trous : les connecteurs]
\zh{民主很重要。首先，人民可以选择领导人。其次，公民有自由。但是，民主也需要责任。所以，我们应该学习民主。}\\
\emph{Mínzhǔ hěn zhòngyào. Shǒuxiān, rénmín kěyǐ xuǎnzé lǐngdǎorén. Qícì, gōngmín yǒu zìyóu. Dànshì, mínzhǔ yě xūyào zérèn. Suǒyǐ, wǒmen yīnggāi xuéxí mínzhǔ.}\\
« La démocratie est très importante. D'abord, le peuple peut choisir ses dirigeants. Ensuite, les citoyens ont la liberté. Mais la démocratie exige aussi des responsabilités. C'est pourquoi nous devons apprendre la démocratie. »

\textbf{Justification :} \zh{首先} « d'abord » et \zh{其次} « ensuite » ordonnent les arguments ; \zh{但是} « mais » introduit l'opposition (liberté ↔ responsabilité) ; \zh{所以} « donc » conclut. Le mot \zh{也} était déjà présent dans le texte et n'était pas à placer ici (piège de l'énoncé : il y avait cinq connecteurs proposés pour quatre trous).
\end{corrige}

\begin{corrige}[Exercice 7 — Remets les mots dans l'ordre]
\begin{enumerate}
\item \zh{公民有言论自由。} (Gōngmín yǒu yánlùn zìyóu.) — « Les citoyens ont la liberté d'expression. » Ordre : Sujet + verbe \zh{有} + complément.
\item \zh{我们应该互相尊重。} (Wǒmen yīnggāi hùxiāng zūnzhòng.) — « Nous devons nous respecter mutuellement. » Ordre : Sujet + modal \zh{应该} + \zh{互相} + verbe.
\item \zh{年轻人应该参加选举。} (Niánqīngrén yīnggāi cānjiā xuǎnjǔ.) — « Les jeunes doivent participer aux élections. »
\item \zh{因为民主很重要，所以我们要学习。} (Yīnwèi mínzhǔ hěn zhòngyào, suǒyǐ wǒmen yào xuéxí.) — « Parce que la démocratie est très importante, nous devons l'apprendre. » La cause \zh{因为} précède la conséquence \zh{所以}.
\end{enumerate}
\end{corrige}

\begin{corrige}[Exercice 8 — Discrimination des caractères]
\begin{enumerate}
\item 民（\textbf{主} / 王）→ \zh{民主}. \zh{主} a un point au-dessus de \zh{王}.
\item （\textbf{权} / 劝）利 → \zh{权利}. \zh{权} a la clé 木 (bois) à gauche ; \zh{劝} a 又 seul.
\item 投（\textbf{票} / 飘）→ \zh{投票}. \zh{票} (clé 示 en bas) ≠ \zh{飘} « flotter au vent » (avec 风).
\item （\textbf{选} / 远）举 → \zh{选举}. Les deux ont la clé 辶 (mouvement), mais \zh{选} contient \zh{先} et \zh{远} contient \zh{元}.
\item 责（\textbf{任} / 住）→ \zh{责任}. \zh{任} = 亻+ \zh{壬} ; \zh{住} « habiter » = 亻+ \zh{主}.
\end{enumerate}
\textbf{Méthode de dictée :} dicter lentement caractère par caractère, vérifier la clé (radical) de chaque caractère, puis comparer avec le tableau de la leçon.
\end{corrige}

\begin{corrige}[Exercice 9 — Compréhension de texte et questions de langue]
\textbf{Questions de compréhension}
\begin{enumerate}
\item Les citoyens ont la liberté d'expression (\zh{言论自由}) et aussi des responsabilités (\zh{责任}).
\item Chaque année, ils peuvent participer aux élections et choisir leurs dirigeants (\zh{参加选举，选择自己的领导人}).
\item Parce que la démocratie est très importante pour tous (\zh{因为民主对大家很重要}).
\end{enumerate}
\textbf{Questions de langue}
\begin{enumerate}
\item Deux connecteurs : \zh{也} « aussi » (addition : liberté + responsabilité) et \zh{因为} « parce que » (cause : pourquoi apprendre la démocratie). On peut aussi citer \zh{互相} « mutuellement » comme marqueur de réciprocité.
\item Le mot d'addition est \zh{也} « aussi ». Exemple : \zh{我是学生，他也是学生。} (Wǒ shì xuésheng, tā yě shì xuésheng.) « Je suis élève, lui aussi est élève. »
\item \zh{我们应该互相尊重，互相学习。} — « Nous devons nous respecter mutuellement et apprendre les uns des autres. »
\end{enumerate}
\textbf{Barème indicatif :} compréhension 6 pts (2 pts par question) ; langue 4 pts (connecteurs 1 pt, \zh{也} + phrase 1,5 pt, traduction 1,5 pt).
\end{corrige}

\begin{corrige}[Exercice 10 — Thème et version type Bac]
\textbf{A. Thème}
\begin{enumerate}
\item \zh{公民有投票的权利和言论自由。} (Gōngmín yǒu tóupiào de quánlì hé yánlùn zìyóu.) — \zh{和} relie les deux compléments.
\item \zh{因为民主很重要，所以我们应该参加选举。} (Yīnwèi mínzhǔ hěn zhòngyào, suǒyǐ wǒmen yīnggāi cānjiā xuǎnjǔ.)
\item \zh{喀麦隆青年也应该学习民主。} (Kāmàilóng qīngnián yě yīnggāi xuéxí mínzhǔ.) — On accepte aussi \zh{喀麦隆的年轻人}.
\end{enumerate}
\textbf{B. Version} — « Dans un pays démocratique, les citoyens ont la liberté d'expression et aussi des responsabilités. Nous devons nous respecter mutuellement et apprendre les uns des autres. »

\textbf{Barème indicatif :} thème 6 pts (2 pts par phrase : structure 1 pt, vocabulaire et caractères 1 pt) ; version 4 pts (fidélité 2 pts, qualité du français 2 pts).

\textbf{Points de vigilance :} \zh{也} avant le verbe ; \zh{因为…所以…} dans cet ordre ; ne pas oublier \zh{的} dans \zh{投票的权利}.
\end{corrige}

\begin{corrige}[Exercice 11 — Expression écrite type Bac]
\textbf{Proposition de rédaction modèle}

\begin{textebox}[La démocratie et les jeunes Camerounais]
民主很重要。首先，公民有投票的权利，也有言论自由。其次，年轻人应该参加选举，选择自己的领导人。但是，民主也需要责任。所以，喀麦隆青年应该学习民主，互相尊重。\\
Mínzhǔ hěn zhòngyào. Shǒuxiān, gōngmín yǒu tóupiào de quánlì, yě yǒu yánlùn zìyóu. Qícì, niánqīngrén yīnggāi cānjiā xuǎnjǔ, xuǎnzé zìjǐ de lǐngdǎorén. Dànshì, mínzhǔ yě xūyào zérèn. Suǒyǐ, Kāmàilóng qīngnián yīnggāi xuéxí mínzhǔ, hùxiāng zūnzhòng.\\
« La démocratie est très importante. D'abord, les citoyens ont le droit de vote et aussi la liberté d'expression. Ensuite, les jeunes doivent participer aux élections et choisir leurs dirigeants. Mais la démocratie exige aussi des responsabilités. C'est pourquoi les jeunes Camerounais doivent apprendre la démocratie et se respecter mutuellement. »
\end{textebox}

Ce modèle compte 5 phrases, environ 75 caractères, 4 connecteurs (\zh{首先}、\zh{其次}、\zh{但是}、\zh{所以}) et 8 mots du lexique (\zh{民主}、\zh{公民}、\zh{投票}、\zh{权利}、\zh{言论自由}、\zh{选举}、\zh{责任}、\zh{尊重}).

\textbf{Barème indicatif (sur 20) :} respect du sujet et du nombre de phrases : 4 pts ; exactitude grammaticale (ordre des mots, \zh{也}, \zh{因为…所以…}, \zh{的}) : 6 pts ; vocabulaire du thème (au moins 5 mots corrects) : 4 pts ; connecteurs (au moins 3) : 3 pts ; écriture des caractères : 3 pts.

\textbf{Points de vigilance :}
\begin{itemize}
\item Ne pas écrire « aussi » en fin de phrase : \zh{也} se place toujours avant le verbe.
\item Ne pas traduire mot à mot « la démocratie est importante pour les jeunes » : dire \zh{民主对年轻人很重要} avec \zh{对}.
\item Soigner les caractères pièges : \zh{主} (avec le point), \zh{权} (clé 木), \zh{票}, \zh{选}, \zh{任}.
\item Une phrase = une idée ; rester dans les structures vues en classe.
\end{itemize}
\end{corrige}

\newpage

\coursec{Fiche bilan}

\begin{popfigure}
\begin{tikzpicture}[mindmap, grow cyclic, every node/.style={concept, font=\small},
root concept/.append style={concept color=popblue, font=\bfseries},
level 1/.append style={level distance=3.2cm, sibling angle=60},
level 1 concept/.append style={font=\footnotesize}]
\node[root concept]{Expression écrite : la démocratie \zh{民主}}
  child[concept color=poporange]{ node[align=center]{Lexique clé\\ \zh{民主 选举 投票}\\ \zh{权利 自由 法律}} }
  child[concept color=popgreen]{ node[align=center]{Écriture des caractères\\ clés et traits\\ \zh{民 主 选 权}} }
  child[concept color=poppurple]{ node[align=center]{Grammaire\\ 是……的\\ 应该 / 可以\\ 为了……} }
  child[concept color=poppink]{ node[align=center]{Structure du texte\\ introduction\\ arguments, conclusion} }
  child[concept color=popgold]{ node[align=center]{Connecteurs\\ \zh{首先 然后}\\ \zh{但是 所以 总之}} }
  child[concept color=popblueL!80!popblue]{ node[align=center]{Traduction\\ thème et version} };
\end{tikzpicture}
\legende{Carte mentale de la leçon : écrire un texte simple sur la démocratie.}
\end{popfigure}

\begin{bilan}
\textbf{1. Lexique clé à connaître par cœur}
\begin{center}
\begin{tabularx}{\linewidth}{|X|X|X|X|}
\hline
\rowcolor{popblueL} Caractères & Pinyin & Nature & Traduction \\
\hline
民主 & mínzhǔ & nom / adj. & démocratie, démocratique \\
\hline
选举 & xuǎnjǔ & nom / verbe & élection, élire \\
\hline
投票 & tóupiào & verbe & voter \\
\hline
公民 & gōngmín & nom & citoyen \\
\hline
权利 & quánlì & nom & droit (d'une personne) \\
\hline
自由 & zìyóu & nom / adj. & liberté, libre \\
\hline
法律 & fǎlǜ & nom & loi, législation \\
\hline
意见 & yìjiàn & nom & opinion, avis \\
\hline
参加 & cānjiā & verbe & participer à \\
\hline
重要 & zhòngyào & adj. & important \\
\hline
\end{tabularx}
\end{center}

\textbf{2. Règles de grammaire à connaître par cœur}
\begin{center}
\begin{tabularx}{\linewidth}{|X|X|X|}
\hline
\rowcolor{popblueL} Structure & Exemple & Traduction \\
\hline
Sujet + 应该 + verbe (devoir, il faut) & \zh{公民应该参加选举。} Gōngmín yīnggāi cānjiā xuǎnjǔ. & Les citoyens doivent participer aux élections. \\
\hline
Sujet + 可以 + verbe (pouvoir, permission) & \zh{每个人都可以发表意见。} Měi ge rén dōu kěyǐ fābiǎo yìjiàn. & Chacun peut exprimer son opinion. \\
\hline
为了 + but + , + action & \zh{为了国家的发展，我们要遵守法律。} Wèile guójiā de fāzhǎn, wǒmen yào zūnshǒu fǎlǜ. & Pour le développement du pays, nous devons respecter la loi. \\
\hline
每 + classificateur + nom + 都 & \zh{每个公民都有权利。} Měi ge gōngmín dōu yǒu quánlì. & Chaque citoyen a des droits. \\
\hline
是……的 (mise en relief) & \zh{民主是很重要的。} Mínzhǔ shì hěn zhòngyào de. & La démocratie est (vraiment) importante. \\
\hline
\end{tabularx}
\end{center}

\textbf{3. Connecteurs pour structurer le texte}
\begin{itemize}
\item \zh{首先} shǒuxiān : d'abord ; \zh{然后} ránhòu : ensuite ; \zh{最后} zuìhòu : enfin.
\item \zh{但是} dànshì : mais ; \zh{所以} suǒyǐ : donc ; \zh{因为} yīnwèi : parce que.
\item \zh{总之} zǒngzhī : en résumé (pour la conclusion).
\end{itemize}

\textbf{4. Méthodes express}
\begin{itemize}
\item \cle{Plan en 3 parties} : introduction (définir le sujet) $\rightarrow$ 2 ou 3 arguments avec connecteurs $\rightarrow$ conclusion avec \zh{总之}.
\item \cle{Thème (français $\rightarrow$ chinois)} : repérer sujet + verbe + complément ; placer le complément de temps ou de lieu avant le verbe ; vérifier les tons et les classificateurs.
\item \cle{Version (chinois $\rightarrow$ français)} : traduire d'abord le sens global, mot à mot ensuite ; attention à \zh{的} qui marque la possession ou la qualification.
\item \cle{Écriture} : respecter l'ordre des traits (haut $\rightarrow$ bas, gauche $\rightarrow$ droite) ; ne pas confondre \zh{民} mín (peuple) et \zh{主} zhǔ (maître, principal) : \zh{民主} = « le peuple dirige ».
\end{itemize}
\end{bilan}

\begin{aretenir}[Checklist avant le Bac]
\begin{itemize}
\item[$\square$] Je connais le lexique du thème : \zh{民主、选举、投票、公民、权利、自由、法律、意见}.
\item[$\square$] J'écris correctement \zh{民} (5 traits) et \zh{主} (5 traits) sans les confondre.
\item[$\square$] Je sais utiliser \zh{应该} et \zh{可以} avant le verbe.
\item[$\square$] Je maîtrise la structure \zh{为了} + but + action.
\item[$\square$] J'utilise \zh{每} + classificateur + \zh{都} correctement.
\item[$\square$] Je structure mon texte : introduction, arguments, conclusion avec \zh{总之}.
\item[$\square$] Je place au moins 3 connecteurs : \zh{首先、但是、所以}.
\item[$\square$] Je place les compléments de temps et de lieu avant le verbe.
\item[$\square$] Je vérifie les tons du pinyin dans mes traductions.
\item[$\square$] Je relis mon texte : caractères complets, ponctuation chinoise ，。, aucune faute de structure.
\end{itemize}
\end{aretenir}

\findecours

\end{document}
